% Les acides α-aminés et la synthèse des peptides — Chimie Tle TI — Cours signé OnBuch+
% Programme officiel MINESEC, Terminale TI. Compiler avec : tectonic L04-acides-alpha-amines-et-peptides.tex
\newcommand{\DOCMATIERE}{Chimie}
\newcommand{\DOCNIVEAU}{Tle TI}
\newcommand{\DOCMODULE}{Module 1 · Chimie organique}
\newcommand{\DOCLECON}{04}
\newcommand{\DOCTITRE}{Les acides α-aminés et la synthèse des peptides}
% OnBuch+ — préambule « cours signés » ENRICHI (compilé avec tectonic / XeLaTeX)
% Système visuel « Pop » d'OnBuch+ : fond crème (jamais blanc), bordures encre
% épaisses, ombres « dures » décalées, titres Archivo Black en capitales.
% Version enrichie pour les cours de chimie : chimie (mhchem, chemfig),
% graphes (pgfplots), boîtes pédagogiques supplémentaires (définition,
% propriété, méthode, attention, le savais-tu, exercice résolu, exercice,
% TP, bilan), page de garde refaite et signature OnBuch+ ancrée en bas.
%
% Macros attendues dans main.tex AVANT \input{preamble} :
%   \DOCMATIERE  \DOCNIVEAU  \DOCMODULE  \DOCLECON (numéro)  \DOCTITRE
\documentclass[11pt]{article}

\usepackage{fontspec}
\usepackage[french]{babel}
\usepackage[a4paper,margin=2cm,top=3.4cm,bottom=2.8cm,headheight=1.4cm,footskip=1.3cm]{geometry}
\usepackage{amsmath,amssymb,amsfonts}
\usepackage[table]{xcolor} % [table] : fournit aussi \rowcolors (alternance de lignes)
\usepackage{pagecolor}
\usepackage{tikz}
\usetikzlibrary{arrows.meta,positioning,calc,shapes.geometric,shapes.misc,shapes.arrows,decorations.pathmorphing,decorations.markings,patterns,shadows,mindmap,trees,fit,backgrounds,matrix,intersections}
\usepackage{pgfplots}
\pgfplotsset{compat=1.18}
\usepackage[version=4]{mhchem}
\usepackage{chemfig}
\usepackage{enumitem}
\usepackage{fancyhdr}
% [most] charge toutes les bibliothèques tcolorbox (listings, minted, raster,
% documentation...) et gonfle inutilement la mémoire TeX sur un document long
% (~300 boîtes) : on ne charge que ce qui est réellement utilisé.
\usepackage[skins,breakable]{tcolorbox}
\usepackage{graphicx}
\usepackage{colortbl}
\usepackage{booktabs}
\usepackage{tabularx}
\usepackage{array}
\usepackage{multicol}
\usepackage{siunitx}
% parse-numbers=false : accepte aussi \SI{2,5 \times 10^{-3}}{...}
\sisetup{locale=FR,output-decimal-marker={,},group-minimum-digits=4,parse-numbers=false}
\DeclareSIUnit\ohm{\ensuremath{\symup{\Omega}}} % U+2126 absent de la police
\usepackage{qrcode}
% \degree : commande fréquemment utilisée par les modèles pour un angle ou une
% température en dehors de \SI{}{} (non fournie par les paquets chargés).
\providecommand{\degree}{^{\circ}}
\usepackage{lastpage}
\usepackage{hyperref}
\hypersetup{hidelinks}

% ============================================================
% Palette « Pop » OnBuch+ — mêmes hex que la classe P (plus_ui.dart)
% ============================================================
\definecolor{poporange}{HTML}{F59321}
\definecolor{popdark}{HTML}{E07A0C}
\definecolor{popcream}{HTML}{FFF6E8}
\definecolor{popink}{HTML}{1C1714}
\definecolor{poppurple}{HTML}{7A5AE0}
\definecolor{popgreen}{HTML}{2E9E6A}
\definecolor{poppink}{HTML}{E0457B}
\definecolor{popblue}{HTML}{2F7DE1}
\definecolor{popgold}{HTML}{C98A12}
\definecolor{popmuted}{HTML}{8A7F76}
\definecolor{popline}{HTML}{EADFD2}
\definecolor{popgray}{HTML}{8A7F76}
\colorlet{popgrey}{popgray}
% Alias tolérants (le modèle nomme parfois les couleurs différemment)
\colorlet{popwhite}{white}
\colorlet{popblack}{popink}
\colorlet{popred}{poppink}
\colorlet{popyellow}{popgold}
% Teintes claires pour fonds de tableaux / zones
\colorlet{popblueL}{popblue!10!white}
\colorlet{poporangeL}{poporange!14!white}
\colorlet{popgreenL}{popgreen!12!white}
\colorlet{poppurpleL}{poppurple!12!white}
\colorlet{poppinkL}{poppink!12!white}
\colorlet{popgoldL}{popgold!14!white}

\pagecolor{popcream}

% ============================================================
% Typographie
% ============================================================
\defaultfontfeatures{Ligatures=TeX}
\setmainfont{PlusJakartaSans-Medium.ttf}[Path=../../../fonts/,BoldFont=PlusJakartaSans-Bold.ttf]
% Maths en sans-serif (Fira Math) avec chiffres et lettres droites de Plus
% Jakarta, pour que les formules se fondent dans le texte.
\usepackage[math-style=ISO,bold-style=ISO]{unicode-math}
\setmathfont{FiraMath-Regular.otf}[Path=../../../fonts/]
\setmathfont{PlusJakartaSans-Medium.ttf}[Path=../../../fonts/,range={up/num,up/{Latin,latin},\mathup}]
\usepackage{icomma} % virgule décimale sans espace en mode math
% Caractères absents de la police de texte (sinon carré vide) : rendus par
% leur équivalent mathématique, en texte comme en formule.
\usepackage{newunicodechar}
\newunicodechar{α}{\ensuremath{\alpha}}
\newunicodechar{Α}{\ensuremath{\alpha}}
\newunicodechar{β}{\ensuremath{\beta}}
\newunicodechar{δ}{\ensuremath{\delta}}
\newunicodechar{✓}{\popcheck{popgreen}}
\newfontfamily\popdisplay{ArchivoBlack-Regular.ttf}[Path=../../../fonts/]
\newfontfamily\poplogo{SpaceGrotesk-Bold.ttf}[Path=../../../fonts/]
\newfontfamily\popsemi{PlusJakartaSans-SemiBold.ttf}[Path=../../../fonts/]
\newfontfamily\popxbold{PlusJakartaSans-ExtraBold.ttf}[Path=../../../fonts/]
\newfontfamily\popxboldls{PlusJakartaSans-ExtraBold.ttf}[Path=../../../fonts/,LetterSpace=3.0]

\setlist[enumerate]{leftmargin=1.7em,itemsep=5pt,topsep=4pt}
\setlist[itemize]{leftmargin=1.7em,itemsep=4pt,topsep=4pt,label={\color{poporange}\textbullet}}
\setlength{\parindent}{0pt}
\setlength{\parskip}{5pt}
\renewcommand{\arraystretch}{1.25}

% chemfig : liaisons un peu plus épaisses, encre Pop
\setchemfig{atom sep=2em,bond style={line width=0.9pt,popink},cram width=3pt}

% pgfplots : style Pop par défaut
\pgfplotsset{
  every axis/.append style={axis line style={popink,line width=1pt},tick style={popink},
    grid style={popline,line width=0.5pt},label style={font=\small},tick label style={font=\footnotesize}},
  popcurve/.style={poporange,line width=1.8pt,no markers,smooth},
}

% ============================================================
% Pill / squiggle
% ============================================================
\newcommand{\poppill}[3]{%
  \tikz[baseline=-0.55ex]\node[draw=popink,line width=1.2pt,rounded corners=2.6pt,%
    fill=#2,inner xsep=6.5pt,inner ysep=3pt]{%
    \popxboldls\fontsize{7}{7}\selectfont\color{#3}\MakeUppercase{#1}};%
}
\newcommand{\popsquiggle}[1]{%
  \tikz[baseline=-0.4ex]{
    \draw[#1,line width=2.6pt,line cap=round]
      plot [smooth] coordinates {(0,0.09) (0.34,0.01) (0.68,0.21) (1.02,0.01) (1.36,0.10)};
  }%
}
% Mot-clé surligné (marqueur orange sous le texte)
\newcommand{\cle}[1]{{\popxbold\color{popdark}#1}}

% ============================================================
% En-tête / pied
% ============================================================
\pagestyle{fancy}
\fancyhf{}
\renewcommand{\headrulewidth}{0pt}
\renewcommand{\footrulewidth}{0pt}
\lhead{\raisebox{-0.15ex}{\poplogo\fontsize{13}{13}\selectfont\color{popink}OnBuch\color{poporange}+}}
\rhead{\poppill{\DOCMATIERE\ \textbullet\ \DOCNIVEAU\ \textbullet\ Leçon \DOCLECON}{white}{popink}}
\lfoot{\popxbold\fontsize{7.6}{7.6}\selectfont\color{popmuted}\MakeUppercase{Cours signé OnBuch+}\ \textbullet\ {\popsemi\DOCTITRE}}
\rfoot{\tikz[baseline=-0.6ex]\node[rounded corners=8.5pt,fill=popink,minimum height=17pt,minimum width=17pt,inner xsep=5pt,inner sep=0]{\popxbold\fontsize{8}{8}\selectfont\color{popcream}\thepage\,/\,\pageref*{LastPage}};}

% ============================================================
% Titres
% ============================================================
\newcommand{\chaptitre}[2]{%
  \begin{tcolorbox}[enhanced,colback=poporange,colframe=popink,boxrule=2.6pt,arc=11pt,
      drop shadow=popink,left=15pt,right=15pt,top=12pt,bottom=11pt,before skip=3pt,after skip=18pt]
    \poppill{#2}{popink}{popcream}\par\vspace{9pt}
    {\raggedright\hyphenpenalty=10000\exhyphenpenalty=10000\popdisplay\fontsize{22}{24}\selectfont\color{popcream}\MakeUppercase{#1}\par}
  \end{tcolorbox}
}
% Section numérotée : pastille orange avec le numéro + titre Archivo
\newcounter{popsec}
\newcommand{\coursec}[1]{%
  \stepcounter{popsec}\setcounter{popsub}{0}%
  \par\vspace{16pt}\noindent
  \begin{minipage}{\linewidth}
  \tikz[baseline=-0.7ex]\node[draw=popink,line width=1.6pt,fill=poporange,rounded corners=4pt,minimum size=22pt,inner sep=0]{\popdisplay\fontsize{12}{12}\selectfont\color{popcream}\thepopsec};%
  \hspace{8pt}{\popdisplay\fontsize{15}{17}\selectfont\color{popink}\MakeUppercase{#1}}\par
  \vspace{4pt}\noindent{\color{poporange}\rule{36pt}{3.6pt}}
  \end{minipage}\par\nopagebreak\vspace{8pt}
}
\newcounter{popsub}
\newcommand{\courssub}[1]{%
  \stepcounter{popsub}%
  \par\vspace{9pt}\noindent{\popxbold\fontsize{11.5}{13}\selectfont\color{popink}{\color{poporange}\thepopsec.\thepopsub}\hspace{6pt}#1}\par\nopagebreak\vspace{4pt}
}

% ============================================================
% Boîtes pédagogiques — même recette : fond blanc, bordure épaisse colorée,
% ombre dure encre, badge plein. Titre optionnel [#1].
% ============================================================
\tcbset{popbase/.style={breakable,enhanced,colback=white,boxrule=2.3pt,arc=9pt,
  shadow={3pt}{-3pt}{0pt}{popink},left=14pt,right=14pt,top=11pt,bottom=11pt,before skip=11pt,after skip=15pt,
  fontupper=\popsemi\fontsize{10.6}{14.5}\selectfont\color{popink},
  fontlower=\popsemi\fontsize{10.6}{14.5}\selectfont\color{popink}}}
% Coche dessinée (le glyphe ✓ n'existe pas dans les polices du document)
\newcommand{\popcheck}[1]{\tikz[baseline=-0.5ex]\draw[#1,line width=1.6pt,line cap=round,line join=round](-0.13,0.0)--(-0.04,-0.09)--(0.13,0.1);}
\newcommand{\popboxhead}[3]{\poppill{#1}{#2}{white}\ifx\relax#3\relax\else\hspace{6pt}{\popxbold\fontsize{10.6}{12}\selectfont\color{popink}#3}\fi\par\vspace{7pt}}

\newtcolorbox{aretenir}[1][]{popbase,colframe=popgreen,before upper={\popboxhead{À retenir}{popgreen}{#1}}}
\newtcolorbox{exemplebox}[1][]{popbase,colframe=poporange,before upper={\popboxhead{Exemple}{poporange}{#1}}}
\newtcolorbox{definition}[1][]{popbase,colframe=popblue,before upper={\popboxhead{Définition}{popblue}{#1}}}
\newtcolorbox{propriete}[1][]{popbase,colframe=poppurple,before upper={\popboxhead{Propriété}{poppurple}{#1}}}
\newtcolorbox{methode}[1][]{popbase,colframe=popgold,before upper={\popboxhead{Méthode}{popgold}{#1}}}
\newtcolorbox{attention}[1][]{popbase,colframe=poppink,before upper={\popboxhead{Attention}{poppink}{#1}}}
\newtcolorbox{experience}[1][]{popbase,colframe=popblue,colback=popblueL,before upper={\popboxhead{Expérience}{popblue}{#1}}}
% « Le savais-tu ? » : carte violette pleine (contexte, culture, Cameroun)
\newtcolorbox{savaistu}[1][]{popbase,colback=poppurple,colframe=popink,
  fontupper=\popsemi\fontsize{10.4}{14.2}\selectfont\color{white},
  before upper={\poppill{Le savais-tu\,?}{white}{poppurple}\ifx\relax#1\relax\else\hspace{6pt}{\popxbold\color{white}#1}\fi\par\vspace{7pt}}}
% Exercice résolu : énoncé en haut, \tcblower puis la solution sur fond crème
\newcounter{popexo}\newcounter{popexr}
\newtcolorbox{exoresolu}[1][]{popbase,colframe=poporange,colbacklower=popcream,
  segmentation style={popink,line width=1.2pt,dashed},
  before upper={\stepcounter{popexr}\popboxhead{Exercice résolu \thepopexr}{poporange}{#1}},
  before lower={\poppill{Solution}{popink}{popcream}\par\vspace{6pt}}}
% Exercice (non résolu en place) — la correction est dans la partie Corrigés
\newtcolorbox{exercice}[1][]{popbase,colframe=popink,
  before upper={\stepcounter{popexo}\popboxhead{Exercice \thepopexo}{popink}{#1}}}
% Corrigé
\newtcolorbox{corrige}[1][]{popbase,colframe=popgreen,colback=popgreenL,
  before upper={\popboxhead{Corrigé}{popgreen}{#1}}}
% Objectifs / prérequis / bilan
\newtcolorbox{objectifs}{popbase,colframe=popink,colback=poporangeL,
  before upper={\popboxhead{Objectifs de la leçon}{popink}{}}}
\newtcolorbox{prerequis}{popbase,colframe=popink,colback=popblueL,
  before upper={\popboxhead{Prérequis}{popblue}{}}}
\newtcolorbox{bilan}{popbase,colframe=popink,colback=white,boxrule=2.8pt,
  before upper={\popboxhead{Fiche bilan}{poporange}{L'essentiel à connaître pour l'examen}}}
% Figure encadrée avec légende
\newtcolorbox{popfigure}[1][]{enhanced,breakable=false,colback=white,colframe=popline,boxrule=1.4pt,arc=8pt,
  left=10pt,right=10pt,top=10pt,bottom=8pt,before skip=10pt,after skip=12pt,halign=center,
  lower separated=false,halign lower=center,
  fontlower=\popsemi\fontsize{9}{11.5}\selectfont\color{popmuted},
  #1}
\newcommand{\legende}[1]{\par\vspace{6pt}{\popsemi\fontsize{9}{11.5}\selectfont\color{popmuted}#1}}

% ============================================================
% Signature OnBuch+
% ============================================================
\newcommand{\poplogoinline}[1]{{\poplogo\fontsize{#1}{#1}\selectfont\color{popink}OnBuch\color{poporange}+}}
\newcommand{\signatureonbuch}{%
  \begin{tcolorbox}[enhanced,colback=popink,colframe=popink,boxrule=2.2pt,arc=10pt,
      drop shadow=poporange,left=14pt,right=14pt,top=10pt,bottom=10pt,before skip=0pt,after skip=0pt]
    \tikz[baseline=-0.7ex]\node[circle,fill=popgreen,draw=popcream,line width=1.3pt,minimum size=22pt,inner sep=0]{\popcheck{white}};%
    \hspace{9pt}\begin{minipage}[c]{0.62\linewidth}
      {\popxbold\fontsize{12}{13}\selectfont\color{popcream}Signé\ \textbullet\ {\poplogo OnBuch\color{poporange}+}}\par\vspace{2pt}
      {\raggedright\popsemi\fontsize{8}{10.5}\selectfont\color{popline}\MakeUppercase{Équipe pédagogique OnBuch+}\\\MakeUppercase{Conforme au programme officiel MINESEC}\par}
    \end{minipage}\hfill
    {\poplogo\fontsize{22}{22}\selectfont\color{popcream}OnBuch\color{poporange}+}
  \end{tcolorbox}%
}

% ============================================================
% Page de garde
% #1 titre · #2 matière · #3 niveau · #4 module · #5 n° leçon · #6 durée ·
% #7 description / objectifs (texte ou itemize)
% ============================================================
\newcommand{\pagegarde}[7]{%
  \thispagestyle{empty}
  \newgeometry{a4paper,margin=1.7cm,top=1.6cm,bottom=1.4cm}%
  \begin{tikzpicture}[remember picture,overlay]
    \fill[poporange!18] ([xshift=-3.2cm,yshift=-3.6cm]current page.north east) circle (4.6cm);
    \fill[poppurple!14] ([xshift=2.4cm,yshift=7.6cm]current page.south west) circle (3.8cm);
  \end{tikzpicture}%
  {\poplogo\fontsize{30}{30}\selectfont\color{popink}OnBuch\color{poporange}+}\hfill
  \raisebox{0.6ex}{\poppill{Cours signé \textbullet\ Premium}{popink}{popcream}}\par
  \vspace{1pt}\popsquiggle{poppurple}\par
  \vspace{8pt}
  \begin{tcolorbox}[enhanced,colback=poporange,colframe=popink,boxrule=2.8pt,arc=13pt,
      drop shadow=popink,left=18pt,right=18pt,top=12pt,bottom=13pt,after skip=10pt]
    \poppill{#2 \textbullet\ #3}{popink}{popcream}\hspace{5pt}\poppill{#4}{white}{popink}\par\vspace{8pt}
    {\popdisplay\fontsize{14}{14}\selectfont\color{popink}LEÇON #5}\par\vspace{4pt}
    {\raggedright\hyphenpenalty=10000\exhyphenpenalty=10000\popdisplay\fontsize{25}{27}\selectfont\color{popcream}\MakeUppercase{#1}\par}
    \vspace{8pt}
    {\popxbold\fontsize{9.5}{9.5}\selectfont\color{popink}\MakeUppercase{Durée conseillée : #6}}
  \end{tcolorbox}
  \begin{tcolorbox}[enhanced,colback=white,colframe=popink,boxrule=2.2pt,arc=10pt,
      drop shadow=popink,left=16pt,right=16pt,top=10pt,bottom=10pt,after skip=8pt]
    \poppill{Dans cette leçon}{poppurple}{white}\par\vspace{5pt}
    \popsemi\fontsize{9.3}{12.6}\selectfont\color{popink}#7
  \end{tcolorbox}
  \noindent\begin{minipage}[t]{0.487\linewidth}
    \begin{tcolorbox}[enhanced,colback=popgreen,colframe=popink,boxrule=2.2pt,arc=10pt,
        drop shadow=popink,left=11pt,right=11pt,top=10pt,bottom=10pt,height=3.1cm,valign=center]
      \tcbox[colback=white,colframe=popink,boxrule=1.3pt,arc=5pt,boxsep=0pt,nobeforeafter,
        left=3pt,right=3pt,top=3pt,bottom=3pt,valign=center]{\qrcode[height=1.9cm]{https://play.google.com/store/apps/details?id=com.luvvix.onbuch&hl=fr}}%
      \hspace{8pt}\begin{minipage}[c]{0.5\linewidth}
        {\popdisplay\fontsize{11}{12.5}\selectfont\color{white}\MakeUppercase{Télécharge OnBuch}}\par\vspace{4pt}
        {\raggedright\popsemi\fontsize{8.4}{11}\selectfont\color{white}Gratuit \textbullet\ Google Play\\Tuteur IA, annales, concours, résultats\par}
      \end{minipage}
    \end{tcolorbox}
  \end{minipage}\hfill
  \begin{minipage}[t]{0.487\linewidth}
    \begin{tcolorbox}[enhanced,colback=poppurple,colframe=popink,boxrule=2.2pt,arc=10pt,
        drop shadow=popink,left=11pt,right=11pt,top=10pt,bottom=10pt,height=3.1cm,valign=center]
      \tikz[baseline=-0.7ex]\node[circle,fill=white,draw=popink,line width=1.3pt,minimum size=1.6cm,inner sep=0]{\popdisplay\fontsize{24}{24}\selectfont\color{poppurple}+};%
      \hspace{8pt}\begin{minipage}[c]{0.6\linewidth}
        {\popdisplay\fontsize{11}{12.5}\selectfont\color{white}\MakeUppercase{Passe à OnBuch+}}\par\vspace{4pt}
        {\raggedright\popsemi\fontsize{8.4}{11}\selectfont\color{white}Tous les cours signés, Tuteur IA illimité, fiches PDF — onglet OnBuch+ de l'app\par}
      \end{minipage}
    \end{tcolorbox}
  \end{minipage}
  \vspace{8pt}
  \popcontenu
  \vfill
  \signatureonbuch
  \restoregeometry
  \newpage
}

% Rangée « Ce cours contient » (page de garde)
\newcommand{\poptile}[3]{% #1 couleur #2 gros texte #3 légende
  \begin{tcolorbox}[enhanced,colback=#1,colframe=popink,boxrule=2pt,arc=9pt,shadow={2.5pt}{-2.5pt}{0pt}{popink},
      left=6pt,right=6pt,top=9pt,bottom=9pt,halign=center,valign=center,height=1.75cm,width=\linewidth,nobeforeafter]
    {\popdisplay\fontsize{12.5}{13}\selectfont\color{white}\MakeUppercase{#2}}\par\vspace{3pt}
    {\centering\popsemi\fontsize{7.8}{9.5}\selectfont\color{white}#3\par}
  \end{tcolorbox}}
\newcommand{\popcontenu}{%
  \par\noindent{\popxboldls\fontsize{7.5}{7.5}\selectfont\color{popmuted}CE COURS SIGNÉ CONTIENT}\par\vspace{7pt}
  \noindent\begin{minipage}[t]{0.235\linewidth}\poptile{popblue}{Cours}{complet, illustré, pas à pas}\end{minipage}\hfill
  \begin{minipage}[t]{0.235\linewidth}\poptile{poporange}{Méthodes}{et exercices résolus}\end{minipage}\hfill
  \begin{minipage}[t]{0.235\linewidth}\poptile{poppink}{Exercices}{type examen + corrigés}\end{minipage}\hfill
  \begin{minipage}[t]{0.235\linewidth}\poptile{popgreen}{Fiche bilan}{l'essentiel à retenir}\end{minipage}\par}

% Sommaire
\newtcolorbox{sommaire}{enhanced,colback=white,colframe=popink,boxrule=2.2pt,arc=10pt,
  drop shadow=popink,left=18pt,right=18pt,top=13pt,bottom=13pt,before skip=6pt,after skip=20pt,
  before upper={\poppill{Sommaire}{poppurple}{white}\par\vspace{9pt}\popsemi\fontsize{10.6}{14.5}\selectfont\color{popink}}}

% Signature de fin de cours — poussée tout en bas de la dernière page
\newcommand{\findecours}{%
  \par\vfill\nopagebreak
  \begin{center}\popsquiggle{poporange}\end{center}\vspace{4pt}
  \signatureonbuch
}
% Glyphes absents de Fira Math : redéfinis à partir de glyphes disponibles
\AtBeginDocument{%
  \def\setminus{\mathbin{\backslash}}\let\smallsetminus\setminus
  \def\vdots{\mathord{\vbox{\baselineskip3.2pt\lineskiplimit0pt\kern4pt\hbox{.}\hbox{.}\hbox{.}}}}%
  \def\ddots{\mathinner{\mkern1mu\raise6.5pt\hbox{.}\mkern2mu\raise3.6pt\hbox{.}\mkern2mu\raise0.7pt\hbox{.}\mkern1mu}}%
  \def\triangle{\mathord{\scalebox{0.9}{$\symup{\Delta}$}}}%
  \def\nabla{\mathord{\raisebox{\depth}{\scalebox{1}[-1]{$\symup{\Delta}$}}}}%
  \def\checkmark{\popcheck{popdark}}%
  \def\star{\mathbin{\ast}}\def\bigstar{\mathord{\ast}}%
  \def\diamond{\mathbin{\scalebox{0.6}{$\Diamond$}}}%
  \def\bigcup{\mathop{\mathchoice{\raisebox{-0.3em}{\scalebox{1.7}{$\cup$}}}{\scalebox{1.25}{$\cup$}}{\cup}{\cup}}\displaylimits}%
  \def\bigcap{\mathop{\mathchoice{\raisebox{-0.3em}{\scalebox{1.7}{$\cap$}}}{\scalebox{1.25}{$\cap$}}{\cap}{\cap}}\displaylimits}%
  \def\lesseqqgtr{\mathrel{\lessgtr}}\def\bigoplus{\mathop{\oplus}\displaylimits}%
}
\providecommand{\ssi}{si et seulement si} % abréviation fréquente

%% FIGFIX-BEGIN v4 — correctifs globaux des figures (docs/onbuch-generation/tools/figfix)
\usepackage{adjustbox}\usepackage{etoolbox}
% toute figure TikZ/pgfplots plus large que la ligne est réduite à la largeur disponible
\BeforeBeginEnvironment{tikzpicture}{\begin{adjustbox}{max width=\linewidth}}
\AfterEndEnvironment{tikzpicture}{\end{adjustbox}}
% légendes pgfplots lisibles par-dessus les courbes
\pgfplotsset{every axis/.append style={legend style={fill=white,fill opacity=0.92,text opacity=1,draw=black!15,inner sep=3pt}}}
% étiquettes d'arêtes (syntaxe edge["..."]) avec fond blanc
\tikzset{every edge quotes/.append style={fill=white,inner sep=1.2pt}}
%% FIGFIX-END
\begin{document}

\pagegarde{Les acides α-aminés et la synthèse des peptides}{Chimie}{Tle TI}{Module 1 · Chimie organique}{04}{2 h}{Cette leçon présente les acides α-aminés, briques élémentaires des protéines : structure, nomenclature, stéréochimie et comportement amphotère en solution. Elle aboutit à la formation de la liaison peptidique et à la synthèse d'un dipeptide, compétence classique de l'épreuve de chimie du Baccalauréat TI.
\vspace{4pt}
\begin{multicols}{2}\small
\begin{itemize}
\item À la fin de cette leçon, je sais écrire la formule générale d'un acide α-aminé et identifier son carbone asymétrique
\item À la fin de cette leçon, je sais nommer un acide α-aminé à partir de sa formule et inversement
\item À la fin de cette leçon, je sais réaliser la représentation de Fischer d'un acide α-aminé et distinguer les configurations L et D
\item À la fin de cette leçon, je sais représenter l'amphion (zwitterion) d'un acide α-aminé
\item À la fin de cette leçon, je sais écrire les équations-bilans montrant le caractère acide et le caractère basique de l'amphion
\item À la fin de cette leçon, je sais écrire l'équation-bilan de l'équilibre entre la forme moléculaire et la forme ionique d'un acide α-aminé
\end{itemize}\end{multicols}}

\begin{sommaire}
\begin{multicols}{2}
\begin{enumerate}
\item Structure et définition des acides α-aminés
\item Nomenclature et représentation de Fischer
\item Propriétés acido-basiques : l'amphion
\item Comportement amphotère de l'amphion
\item La liaison peptidique et les dipeptides
\item Synthèse dirigée d'un dipeptide (complément Bac)
\item Travaux pratiques
\item Méthodes et exercices résolus
\item Exercices
\item Corrigés des exercices
\item Fiche bilan
\end{enumerate}
\end{multicols}
\end{sommaire}

\begin{objectifs}
\begin{itemize}
\item À la fin de cette leçon, je sais écrire la formule générale d'un acide α-aminé et identifier son carbone asymétrique
\item À la fin de cette leçon, je sais nommer un acide α-aminé à partir de sa formule et inversement
\item À la fin de cette leçon, je sais réaliser la représentation de Fischer d'un acide α-aminé et distinguer les configurations L et D
\item À la fin de cette leçon, je sais représenter l'amphion (zwitterion) d'un acide α-aminé
\item À la fin de cette leçon, je sais écrire les équations-bilans montrant le caractère acide et le caractère basique de l'amphion
\item À la fin de cette leçon, je sais écrire l'équation-bilan de l'équilibre entre la forme moléculaire et la forme ionique d'un acide α-aminé
\item À la fin de cette leçon, je sais identifier la liaison peptidique dans une molécule et nommer un peptide
\item À la fin de cette leçon, je sais écrire l'équation-bilan de formation d'un dipeptide à partir de deux acides α-aminés
\end{itemize}
\end{objectifs}

\begin{prerequis}
\begin{itemize}
\item Fonctions organiques : groupe carboxyle –COOH (acide carboxylique) et groupe amine –NH2 (classe de Première)
\item Réactions acide-base : couples acide/base, transfert de proton H+
\item Stéréochimie : carbone asymétrique, énantiomères, représentation de Cram (début d'année de Terminale)
\item Écriture et équilibrage d'équations-bilans en chimie organique
\item Nomenclature des composés organiques oxygénés
\end{itemize}
\end{prerequis}

\newpage

\coursec{Structure et définition des acides \texorpdfstring{$\alpha$}{alpha}-aminés}

Au marché Mokolo à Yaoundé, les vendeuses vantent le ndolé et le eru, riches en protéines, pour « donner la force » aux écoliers. Et elles ont raison : les protéines sont indispensables à la construction des muscles, des cheveux, de la peau, des enzymes. Mais que deviennent ces protéines dans l'organisme ? Lors de la digestion, elles sont découpées en petites molécules appelées \cle{acides $\alpha$-aminés}, véritables briques de la vie. Comment ces briques sont-elles construites ? Pourquoi possèdent-elles des propriétés si particulières ? Et comment s'assemblent-elles pour former les peptides puis les protéines ? C'est ce que nous allons découvrir dans cette leçon, qui nous mènera jusqu'à la synthèse d'un dipeptide. Commençons par le commencement : la structure de ces fameuses briques.

\courssub{Les protéines, des polymères naturels d'acides \texorpdfstring{$\alpha$}{alpha}-aminés}

Imagine un collier de perles : chaque perle est une petite molécule, et le collier entier est une très longue chaîne. Une \cle{protéine} fonctionne exactement ainsi : c'est un \cle{polymère} naturel, c'est-à-dire une macromolécule formée par l'enchaînement d'un très grand nombre de petites molécules identiques ou semblables, appelées \emph{monomères}. Dans le cas des protéines, les monomères sont les acides $\alpha$-aminés.

Comment le chimiste peut-il le vérifier ? En cassant le collier, perle par perle. C'est exactement ce que réalise la réaction d'\cle{hydrolyse} : en chauffant une protéine en milieu acide concentré (par exemple dans une solution d'acide chlorhydrique \ce{HCl} à \SI{6}{mol.L^{-1}}, à environ \SI{110}{\celsius} pendant plusieurs heures), on rompt toutes les liaisons qui unissent les monomères, et l'on obtient un mélange d'acides $\alpha$-aminés libres.

\begin{experience}[Hydrolyse d'une protéine]
\textbf{Dispositif et protocole :} on introduit dans un ballon un fragment de protéine (blanc d'œuf, gélatine) avec une solution concentrée d'acide chlorhydrique. On chauffe à reflux pendant plusieurs heures, puis on neutralise et on analyse le mélange par chromatographie sur couche mince (CCM) en présence de témoins.\par
\textbf{Observations :} la chromatographie révèle de nombreuses taches, chacune correspondant à un acide $\alpha$-aminé différent. On en dénombre jusqu'à une vingtaine de types différents.\par
\textbf{Interprétation :} la protéine est bien un assemblage d'acides $\alpha$-aminés. L'hydrolyse a rompu les liaisons qui les unissaient (nous verrons plus loin qu'il s'agit des \emph{liaisons peptidiques}).
\end{experience}

\begin{aretenir}[L'essentiel sur les protéines]
Les protéines sont des \cle{polymères naturels} constitués d'enchaînements d'acides $\alpha$-aminés. Leur hydrolyse totale fournit un mélange d'environ \textbf{20 acides $\alpha$-aminés différents}, que l'on appelle les acides $\alpha$-aminés \emph{naturels} (ou \emph{protéinogènes}). Toutes les protéines du vivant, du ndolé à l'insuline, sont construites avec ce même jeu de 20 briques, assemblées dans des ordres et des longueurs différents.
\end{aretenir}

\courssub{Définition et formule générale d'un acide \texorpdfstring{$\alpha$}{alpha}-aminé}

Regardons maintenant une de ces briques de très près. Le nom « acide $\alpha$-aminé » dit déjà presque tout :
\begin{itemize}
\item \textbf{« acide »} : la molécule possède un groupe \cle{carboxyle} \ce{-COOH}, qui confère un caractère acide ;
\item \textbf{« aminé »} : la molécule possède un groupe \cle{amine} \ce{-NH2}, qui confère un caractère basique ;
\item \textbf{« $\alpha$ »} : les deux groupes sont portés par le \emph{même} atome de carbone, celui qui est voisin du groupe carboxyle, appelé \cle{carbone $\alpha$}.
\end{itemize}

\begin{definition}[Acide $\alpha$-aminé]
Un \cle{acide $\alpha$-aminé} est un composé organique possédant un groupe amine \ce{-NH2} et un groupe carboxyle \ce{-COOH} fixés sur le \textbf{même atome de carbone}, appelé \cle{carbone $\alpha$} (noté $C_\alpha$). Sa formule générale est :
\[ \ce{H2N-CHR-COOH} \quad \text{ou} \quad \chemfig{H_2N-\chembelow{C}{H}(-[6]R)-C(=[2]O)-OH} \]
où \ce{R} est la \cle{chaîne latérale}, un groupe alkyle ou un groupe fonctionnel qui distingue les acides $\alpha$-aminés les uns des autres.
\end{definition}

\begin{popfigure}
\begin{center}
\begin{tikzpicture}
\node at (0,0) {\chemfig{H_2N-[2]C(-[2]H)(-[6]R)-[2]C(=[2]O)-[0]OH}};
\node[popink,font=\bfseries] at (-2.9,0.15) {\ce{NH2}};
\node[popgreen,font=\bfseries] at (-0.9,1.15) {H};
\node[poppurple,font=\bfseries] at (-0.9,-1.15) {R};
\node[poporange,font=\bfseries] at (2.3,0.15) {COOH};
\node[popink,font=\Large\bfseries] at (-0.55,0.55) {*};
\node[font=\small\itshape,popdark] at (-0.9,-1.9) {carbone $\alpha$ (asymétrique)};
\draw[-{Stealth},popdark] (-0.9,-1.65) -- (-0.75,-0.35);
\node[font=\small,popink] at (-4.1,-0.7) {groupe amine};
\draw[-{Stealth},popink] (-3.6,-0.45) -- (-3.1,-0.05);
\node[font=\small,poporange] at (4.1,-0.7) {groupe carboxyle};
\draw[-{Stealth},poporange] (3.6,-0.45) -- (3.1,-0.05);
\node[font=\small,poppurple] at (1.6,-1.6) {chaîne latérale};
\draw[-{Stealth},poppurple] (0.6,-1.35) -- (-0.6,-1.05);
\end{tikzpicture}
\end{center}
\legende{Formule générale d'un acide $\alpha$-aminé : le carbone $\alpha$ (marqué d'un astérisque) porte quatre groupes : \ce{-NH2}, \ce{-COOH}, \ce{-H} et \ce{-R}.}
\end{popfigure}

\begin{attention}[Ne pas confondre $\alpha$, $\beta$, $\gamma$...]
La lettre grecque indique la \textbf{position} du groupe amine par rapport au groupe carboxyle. Dans un acide $\alpha$-aminé, \ce{NH2} est sur le carbone \emph{directement voisin} de \ce{COOH}. Si l'amine est sur le carbone suivant, on parle d'acide $\beta$-aminé (exemple : \ce{H2N-CH2-CH2-COOH}, l'acide $\beta$-alaniné), qui \textbf{n'est pas} un acide $\alpha$-aminé et n'entre pas dans la constitution des protéines. Vérifie toujours que \ce{NH2} et \ce{COOH} sont sur le \emph{même} carbone.
\end{attention}

\courssub{Le carbone \texorpdfstring{$\alpha$}{alpha}, un carbone asymétrique}

Voici maintenant la propriété structurale la plus importante des acides $\alpha$-aminés. Compte les groupes fixés sur le carbone $\alpha$ : un hydrogène \ce{H}, un groupe amine \ce{NH2}, un groupe carboxyle \ce{COOH} et une chaîne latérale \ce{R}. Quatre groupes \emph{différents} !

\begin{definition}[Carbone asymétrique]
Un \cle{carbone asymétrique} est un atome de carbone tétraédrique lié à \textbf{quatre atomes ou groupes d'atomes différents}. On le repère par un astérisque : $C^{*}$.
\end{definition}

\begin{methode}[Repérer un carbone asymétrique dans un acide $\alpha$-aminé]
\begin{enumerate}
\item Écris la formule semi-développée de la molécule et repère le carbone $\alpha$ : c'est celui qui porte à la fois \ce{-NH2} et \ce{-COOH}.
\item Liste les quatre groupes fixés sur ce carbone : \ce{H}, \ce{NH2}, \ce{COOH}, \ce{R}.
\item Compare-les : s'ils sont tous différents, le carbone $\alpha$ est asymétrique. Dans un acide $\alpha$-aminé, c'est \textbf{presque toujours} le cas.
\item Vérifie le cas particulier : si $\ce{R} = \ce{H}$ (glycine), le carbone $\alpha$ porte \emph{deux} atomes d'hydrogène identiques : il n'est \textbf{pas} asymétrique.
\end{enumerate}
\end{methode}

La conséquence est capitale : une molécule possédant un carbone asymétrique est \cle{chirale}, c'est-à-dire qu'elle n'est pas superposable à son image dans un miroir, comme ta main gauche n'est pas superposable à ta main droite. Elle existe alors sous deux formes appelées \cle{énantiomères}, images l'une de l'autre dans un miroir.

\begin{aretenir}[Chiralité des acides $\alpha$-aminés]
Tout acide $\alpha$-aminé dont la chaîne latérale \ce{R} est différente de \ce{H} possède un carbone $\alpha$ asymétrique : il est \cle{chiral} et existe sous deux énantiomères. Nous verrons dans la section suivante comment les distinguer (configurations L et D) grâce à la représentation de Fischer.
\end{aretenir}

\begin{attention}[Erreur fréquente : « tous les acides $\alpha$-aminés sont chiraux »]
\textbf{Faux !} La \cle{glycine}, de formule \ce{H2N-CH2-COOH}, possède une chaîne latérale réduite à un simple atome d'hydrogène ($\ce{R} = \ce{H}$). Son carbone $\alpha$ porte donc \textbf{deux} atomes d'hydrogène identiques : il n'est pas asymétrique, et la glycine est \textbf{achirale} (superposable à son image miroir). C'est le seul des 20 acides $\alpha$-aminés naturels dans ce cas. Au Baccalauréat, cette exception est un piège classique : cite-la systématiquement.
\end{attention}

\courssub{Exemples usuels d'acides \texorpdfstring{$\alpha$}{alpha}-aminés}

Les quelque 20 acides $\alpha$-aminés naturels ne diffèrent les uns des autres que par leur chaîne latérale \ce{R}. C'est elle qui dicte leurs propriétés : taille, polarité, caractère hydrophile ou hydrophobe. En voici cinq à connaître parfaitement.

\begin{popfigure}
\begin{center}
\renewcommand{\arraystretch}{2.2}
\begin{tabularx}{\linewidth}{|l|c|>{\centering\arraybackslash}X|}
\hline
\rowcolor{popblueL} \textbf{Nom} & \textbf{Chaîne latérale R} & \multicolumn{1}{c|}{\textbf{Formule semi-développée}} \\
\hline
Glycine (Gly) & \ce{H} & \chemfig{H_2N-CH_2-C(=[2]O)-OH} \\
\hline
Alanine (Ala) & \ce{CH3} & \chemfig{H_2N-CH(-[6]CH_3)-C(=[2]O)-OH} \\
\hline
Valine (Val) & \ce{CH(CH3)2} & \chemfig{H_2N-CH(-[6]CH(-[5]CH_3)(-[7]CH_3))-C(=[2]O)-OH} \\
\hline
Sérine (Ser) & \ce{CH2OH} & \chemfig{H_2N-CH(-[6]CH_2-OH)-C(=[2]O)-OH} \\
\hline
Phénylalanine (Phe) & \ce{CH2-C6H5} & \chemfig{H_2N-CH(-[6]CH_2-**6(-----))-C(=[2]O)-OH} \\
\hline
\end{tabularx}
\end{center}
\legende{Cinq acides $\alpha$-aminés usuels. Seule la chaîne latérale \ce{R} change ; le motif \ce{H2N-CH-COOH} est commun à tous.}
\end{popfigure}

Remarque la diversité des chaînes latérales : la sérine porte un groupe hydroxyle \ce{-OH} (chaîne polaire), la valine et la phénylalanine portent des chaînes hydrocarbonées (apolaires, hydrophobes). Cette variété explique l'immense diversité des protéines.

\begin{exemplebox}[Identifier un acide $\alpha$-aminé]
La molécule de formule semi-développée \ce{H2N-CH(CH2OH)-COOH} est-elle un acide $\alpha$-aminé ? Si oui, lequel ?\par
\textbf{Réponse :} le carbone porteur de \ce{NH2} porte aussi directement le groupe \ce{COOH} : c'est bien un carbone $\alpha$. La chaîne latérale est $\ce{R} = \ce{CH2OH}$ : il s'agit de la \textbf{sérine}. Son carbone $\alpha$ porte quatre groupes différents (\ce{H}, \ce{NH2}, \ce{COOH}, \ce{CH2OH}) : il est asymétrique, donc la sérine est chirale.
\end{exemplebox}

\begin{exoresolu}[Carbone asymétrique et formule brute]
\textbf{Énoncé.} L'alanine a pour formule semi-développée \ce{H2N-CH(CH3)-COOH}.
\begin{enumerate}
\item Montrer que son carbone $\alpha$ est asymétrique.
\item Déterminer sa formule brute et calculer sa masse molaire moléculaire.
\item On dissout $m = \SI{4,45}{g}$ d'alanine dans de l'eau distillée pour obtenir $V = \SI{500}{mL}$ de solution. Calculer la concentration molaire $C$ de la solution.
\end{enumerate}
Données : $M(\ce{C}) = \SI{12}{g.mol^{-1}}$, $M(\ce{H}) = \SI{1}{g.mol^{-1}}$, $M(\ce{O}) = \SI{16}{g.mol^{-1}}$, $M(\ce{N}) = \SI{14}{g.mol^{-1}}$.
\tcblower
\textbf{Solution détaillée.}
\begin{enumerate}
\item Le carbone $\alpha$ est le carbone central. Il est lié à quatre groupes : \ce{-H}, \ce{-NH2}, \ce{-COOH} et \ce{-CH3}. Ces quatre groupes sont \textbf{tous différents} : le carbone $\alpha$ est donc \textbf{asymétrique} et l'alanine est chirale.
\item Comptons les atomes : 3 carbones, 7 hydrogènes ($2 + 1 + 3 + 1$), 1 azote, 2 oxygènes. Formule brute : \ce{C3H7NO2}.
\[ M = 3 \times 12 + 7 \times 1 + 14 + 2 \times 16 = 36 + 7 + 14 + 32 = \SI{89}{g.mol^{-1}}. \]
\item Quantité de matière dissoute :
\[ n = \frac{m}{M} = \frac{4,45}{89} = 5,0 \times 10^{-2}~\text{mol}. \]
Concentration molaire :
\[ C = \frac{n}{V} = \frac{5,0 \times 10^{-2}}{0,500} = \SI{0,10}{mol.L^{-1}}. \]
\end{enumerate}
\textbf{Conclusion :} la solution d'alanine a une concentration de $\SI{0,10}{mol.L^{-1}}$.
\end{exoresolu}

\begin{savaistu}[Les mêmes 20 briques pour tout le vivant]
Les protéines du ndolé, du poisson fumé ou du eru consommés au Cameroun sont constituées des \textbf{mêmes 20 acides $\alpha$-aminés} que toutes les protéines du vivant : celles de l'être humain, du baobab, des bactéries du sol ou du plancton de l'océan. Seuls changent l'ordre d'enchaînement des briques et la longueur de la chaîne, un peu comme les mêmes 26 lettres de l'alphabet permettent d'écrire tous les livres du monde. Certains de ces acides $\alpha$-aminés, dits \emph{essentiels} (comme la valine et la phénylalanine), ne peuvent pas être fabriqués par notre organisme : ils doivent obligatoirement être apportés par l'alimentation. D'où l'importance d'une alimentation variée et riche en protéines !
\end{savaistu}

\begin{exercice}[À toi de jouer]
On considère les trois molécules suivantes :
\[ \text{(A) } \ce{H2N-CH2-CH2-COOH} \qquad \text{(B) } \ce{H2N-CH(CH3)-COOH} \qquad \text{(C) } \ce{H2N-CH2-COOH} \]
\begin{enumerate}
\item Laquelle de ces molécules n'est \textbf{pas} un acide $\alpha$-aminé ? Justifier.
\item Parmi les acides $\alpha$-aminés, lequel n'est pas chiral ? Justifier en examinant les groupes portés par le carbone $\alpha$.
\item Donner le nom usuel des acides $\alpha$-aminés identifiés.
\end{enumerate}
\end{exercice}

\begin{corrige}[Exercice : À toi de jouer]
\begin{enumerate}
\item La molécule (A) n'est pas un acide $\alpha$-aminé : le groupe \ce{NH2} et le groupe \ce{COOH} sont portés par des carbones \emph{différents} (l'amine est en position $\beta$). C'est un acide $\beta$-aminé.
\item La molécule (C) possède $\ce{R} = \ce{H}$ : son carbone $\alpha$ porte deux atomes d'hydrogène identiques, il n'est pas asymétrique. La molécule (C) n'est donc pas chirale. La molécule (B), elle, porte quatre groupes différents (\ce{H}, \ce{NH2}, \ce{COOH}, \ce{CH3}) : elle est chirale.
\item (B) est l'\textbf{alanine} et (C) est la \textbf{glycine}.
\end{enumerate}
\end{corrige}

\begin{aretenir}[Bilan de la section]
\begin{itemize}
\item Les protéines sont des polymères naturels dont l'hydrolyse fournit environ 20 acides $\alpha$-aminés.
\item Un acide $\alpha$-aminé possède un groupe \ce{-NH2} et un groupe \ce{-COOH} portés par le même carbone, le carbone $\alpha$ ; sa formule générale est \ce{H2N-CHR-COOH}.
\item Le carbone $\alpha$ porte quatre groupes différents (\ce{H}, \ce{NH2}, \ce{COOH}, \ce{R}) : c'est un carbone asymétrique, et la molécule est chirale, \textbf{sauf la glycine} ($\ce{R} = \ce{H}$).
\item Les acides $\alpha$-aminés ne diffèrent que par leur chaîne latérale \ce{R}, qui détermine leurs propriétés.
\end{itemize}
Dans la section suivante, nous apprendrons à nommer ces molécules et à représenter leurs deux énantiomères grâce à la représentation de Fischer.
\end{aretenir}

\coursec{Nomenclature et représentation de Fischer}

Dans la section précédente, nous avons établi la structure générale des acides $\alpha$-aminés : un carbone $\alpha$ portant un groupe carboxyle ($-\ce{COOH}$), un groupe amine ($-\ce{NH2}$), un atome d'hydrogène et une chaîne latérale $\ce{R}$. Nous allons maintenant apprendre à les nommer rigoureusement et, surtout, à représenter leur stéréochimie, car la très grande majorité de ces molécules sont chirales. C'est une compétence indispensable pour comprendre la structure des protéines que nous construirons plus tard.

\courssub{Nomenclature systématique et noms usuels}

\paragraph{Nomenclature systématique (IUPAC).} Selon la nomenclature officielle, un acide $\alpha$-aminé est un \textbf{acide 2-aminoalcanoïque}. La numérotation de la chaîne carbonée part obligatoirement du carbone du groupe carboxyle (C1). Le groupe amine se trouve donc en position 2. La chaîne latérale $\ce{R}$ détermine la longueur et la ramification de la chaîne principale.

\begin{exemplebox}[Nommage systématique de l'alanine]
L'alanine a pour formule semi-développée $\ce{CH3-CH(NH2)-COOH}$.
\begin{enumerate}
    \item La chaîne carbonée la plus longue contenant le $\ce{COOH}$ compte 3 atomes de carbone $\rightarrow$ \emph{propanoïque}.
    \item Le groupe amine est sur le C2 $\rightarrow$ \emph{2-amino}.
    \item Nom systématique : \textbf{acide 2-aminopropanoïque}.
\end{enumerate}
\end{exemplebox}

\begin{definition}[Noms usuels et abréviations]
Bien que la nomenclature systématique soit rigoureuse, les chimistes et biochimistes utilisent quasi exclusivement des \textbf{noms usuels} (triviaux) et des \textbf{abréviations à trois lettres} (code à une lettre aussi utilisé en biochimie). Il est impératif de connaître les plus courants pour le Baccalauréat.
\end{definition}

\begin{center}
\begin{tabularx}{\linewidth}{|l|l|l|l|X|}
\hline
\rowcolor{popblueL}
\textbf{Nom usuel} & \textbf{Nom systématique (IUPAC)} & \textbf{Formule semi-développée} & \textbf{Abréviation (3 lettres)} & \textbf{Chaîne latérale $\ce{R}$} \\
\hline
Glycine & Acide aminoéthanoïque & $\ce{H2N-CH2-COOH}$ & Gly & $\ce{-H}$ \\
Alanine & Acide 2-aminopropanoïque & $\ce{CH3-CH(NH2)-COOH}$ & Ala & $\ce{-CH3}$ \\
Valine & Acide 2-amino-3-méthylbutanoïque & $\ce{(CH3)2CH-CH(NH2)-COOH}$ & Val & $\ce{-CH(CH3)2}$ \\
Leucine & Acide 2-amino-4-méthylpentanoïque & $\ce{(CH3)2CH-CH2-CH(NH2)-COOH}$ & Leu & $\ce{-CH2-CH(CH3)2}$ \\
Sérine & Acide 2-amino-3-hydroxypropanoïque & $\ce{HO-CH2-CH(NH2)-COOH}$ & Ser & $\ce{-CH2OH}$ \\
Cystéine & Acide 2-amino-3-sulfanylpropanoïque & $\ce{HS-CH2-CH(NH2)-COOH}$ & Cys & $\ce{-CH2SH}$ \\
\hline
\end{tabularx}
\end{center}

\paragraph{Remarques importantes.}
\begin{itemize}
    \item La \textbf{glycine} ($\ce{R=H}$) est le seul acide $\alpha$-aminé \textbf{non chiral} (son carbone $\alpha$ porte deux hydrogènes). Elle n'a pas d'énantiomères.
    \item Les cinq autres exemples ci-dessus possèdent un \textbf{carbone asymétrique} (le C$\alpha$ lié à 4 groupes différents : $\ce{H}$, $\ce{NH2}$, $\ce{COOH}$, $\ce{R}$). Ils existent donc sous forme de deux énantiomères (configurations \textbf{L} et \textbf{D}).
    \item Les abréviations à trois lettres (Gly, Ala, Val, Leu, Ser, Cys) sont les « briques » avec lesquelles nous écrirons les séquences peptidiques (ex. : Ala-Gly pour un dipeptide).
\end{itemize}

\courssub{La représentation de Fischer : convention et lecture}

Pour représenter la configuration absolue des acides $\alpha$-aminés chiraux, on utilise la \textbf{projection de Fischer}, une convention graphique bidimensionnelle normalisée.

\begin{definition}[Projection de Fischer]
C'est une représentation plane d'une molécule chirale où :
\begin{itemize}
    \item La chaîne carbonée principale est disposée \textbf{verticalement}.
    \item Le groupe le plus oxydé (le groupe carboxyle $\ce{-COOH}$ pour les acides aminés) est placé \textbf{en haut}.
    \item Le carbone asymétrique est au centre de la croix.
    \item \textbf{Les liaisons horizontales sortent du plan de la feuille (vers l'observateur).}
    \item \textbf{Les liaisons verticales rentrent dans le plan de la feuille (loin de l'observateur).}
\end{itemize}
\end{definition}

\begin{attention}[Piège majeur : lecture des traits]
Ne confondez jamais : \textbf{Horizontal = Sort (vers vous)} / \textbf{Vertical = Rentrent (loin de vous)}. C'est l'inverse de la représentation de Cram (traits pleins / traits pointillés). Une erreur de lecture inverse la configuration absolue !
\end{attention}

\begin{popfigure}
\begin{tikzpicture}[scale=1.2, every node/.style={font=\small}]
    % L-Alanine (NH2 à GAUCHE)
    \begin{scope}[xshift=-4cm]
        % Vertical bonds (going back) - dashed
        \draw[popdark, densely dashed, thick] (0,1.5) -- (0,0.4);
        \draw[popdark, densely dashed, thick] (0,-0.4) -- (0,-1.5);
        % Horizontal bonds (coming out) - wedges (thick with triangle head)
        \draw[popblue, line width=3pt, -{Triangle[length=3pt,width=4pt]}] (-0.8,0) -- (0,0); % NH2 left, out
        \draw[poporange, line width=3pt, -{Triangle[length=3pt,width=4pt]}] (0,0) -- (0.8,0); % H right, out
        % Labels
        \node[above] at (0,1.5) {\ce{COOH}};
        \node[below] at (0,-1.5) {\ce{CH3}};
        \node[left] at (-1,0) {\ce{NH2}};
        \node[right] at (1,0) {\ce{H}};
        % Chiral center dot
        \fill[popdark] (0,0) circle (2pt);
        \node[below=0.6cm] at (0,-1.5) {\textbf{L-Alanine}};
        \node[above=0.3cm] at (0,1.8) {$\ce{NH2}$ à \textbf{GAUCHE}};
    \end{scope}

    % D-Alanine (NH2 à DROITE)
    \begin{scope}[xshift=4cm]
        % Vertical bonds (going back) - dashed
        \draw[popdark, densely dashed, thick] (0,1.5) -- (0,0.4);
        \draw[popdark, densely dashed, thick] (0,-0.4) -- (0,-1.5);
        % Horizontal bonds (coming out) - wedges
        \draw[poporange, line width=3pt, -{Triangle[length=3pt,width=4pt]}] (-0.8,0) -- (0,0); % H left, out
        \draw[popblue, line width=3pt, -{Triangle[length=3pt,width=4pt]}] (0,0) -- (0.8,0); % NH2 right, out
        % Labels
        \node[above] at (0,1.5) {\ce{COOH}};
        \node[below] at (0,-1.5) {\ce{CH3}};
        \node[left] at (-1,0) {\ce{H}};
        \node[right] at (1,0) {\ce{NH2}};
        % Chiral center dot
        \fill[popdark] (0,0) circle (2pt);
        \node[below=0.6cm] at (0,-1.5) {\textbf{D-Alanine}};
        \node[above=0.3cm] at (0,1.8) {$\ce{NH2}$ à \textbf{DROITE}};
    \end{scope}
\end{tikzpicture}
\legende{Projections de Fischer de la L-alanine et de la D-alanine. Les traits horizontaux (épaissis, pointe vers le centre) sortent du plan ; les traits verticaux (pointillés) rentrent dans le plan.}
\end{popfigure}

\courssub{Configurations L et D : la règle de l'alanine}

La configuration absolue (L ou D) est définie par référence à la \textbf{glycéraldéhyde} (référence historique). Pour les acides $\alpha$-aminés, la règle pratique est la suivante :

\begin{methode}[Identifier la configuration L ou D d'un acide $\alpha$-aminé sur une projection de Fischer]
\begin{enumerate}
    \item Orienter la chaîne carbonée \textbf{verticalement} avec le groupe \textbf{$\ce{-COOH}$ en haut} (C1).
    \item Placer la chaîne latérale $\ce{R}$ \textbf{en bas}.
    \item Observer la position du groupe \textbf{$\ce{-NH2}$} sur le carbone asymétrique (au centre) :
    \begin{itemize}
        \item $\ce{-NH2}$ à \textbf{GAUCHE} $\rightarrow$ Configuration \textbf{L} (série L).
        \item $\ce{-NH2}$ à \textbf{DROITE} $\rightarrow$ Configuration \textbf{D} (série D).
    \end{itemize}
    \item (L'hydrogène $\ce{H}$ occupe automatiquement la place restante).
\end{enumerate}
\end{methode}

\begin{aretenir}[Règle mnémotechnique : « COOH en haut, R en bas, NH2 à gauche = L »]
Pour les acides $\alpha$-aminés \textbf{naturels} (ceux qui constituent les protéines des êtres vivants, de la bactérie à l'homme), la configuration est quasi exclusivement \textbf{L}. La D-alanine existe dans certaines parois bactériennes (peptidoglycanes) ou certains peptides antibiotiques (gramicidine), mais jamais dans les protéines ribosomales.
\end{aretenir}

\begin{popfigure}
\begin{tikzpicture}[scale=1.1, every node/.style={font=\small, transform shape}]
    % Fischer representation with 3D hints only (clearer than Cram conversion for this level)
    \begin{scope}[xshift=-3cm]
        \node[align=center] at (0,2.2) {\textbf{Projection de Fischer (Convention 3D)}};
        % Vertical bonds (going back) - dashed
        \draw[popdark, densely dashed, thick] (0,1.5) -- (0,0.4);
        \draw[popdark, densely dashed, thick] (0,-0.4) -- (0,-1.5);
        % Horizontal bonds (coming out) - solid wedges
        \draw[popblue, line width=3pt, -{Triangle[length=3pt,width=4pt]}] (-0.8,0) -- (0,0);
        \draw[poporange, line width=3pt, -{Triangle[length=3pt,width=4pt]}] (0,0) -- (0.8,0);
        \fill[popdark] (0,0) circle (2pt);
        \node[above] at (0,1.5) {\ce{COOH}};
        \node[below] at (0,-1.5) {\ce{R}};
        \node[left] at (-1,0) {\ce{NH2}};
        \node[right] at (1,0) {\ce{H}};
        \node[align=center] at (0,-2.0) {$\ce{NH2}$ à \textbf{GAUCHE} $\rightarrow$ \textbf{L-configuration}\\
        (Acide aminé naturel)};
    \end{scope}

    \begin{scope}[xshift=3cm]
        \node[align=center] at (0,2.2) {\textbf{Représentation de Cram (équivalente)}};
        % Cram for L-amino acid: COOH up (plane), R down (plane), H wedge (out), NH2 dash (back) is ONE convention.
        % Better: Show the tetrahedron with wedges/dashes matching Fischer: Horizontal=Wedge, Vertical=Dash.
        % But standard Cram has 2 in plane. Let's draw a standard "Fischer-to-Cram" view:
        % Carbon center. COOH up (dash/back), R down (dash/back), H left (wedge/out), NH2 right (wedge/out) -> Impossible in standard Cram (2 wedges, 2 dashes).
        % Standard pedagogical representation in FR Terminale: Show the cross with wedges/dashes.
        % I will draw a perspective tetrahedron (Cram) where the horizontal plane is the paper plane.
        \node at (0,0) {\chemfig{
            % Central C
            C
            (-[2,0.8,,,draw=popdark, densely dashed]COOH) % Up, back (dashed)
            (-[-90,0.8,,,draw=popdark, densely dashed]R)   % Down, back (dashed)
            (-[180,0.8,,,draw=popblue, line width=2pt, -{Triangle[length=2pt,width=3pt]}]NH_2) % Left, out (wedge)
            (-[0,0.8,,,draw=poporange, line width=2pt, -{Triangle[length=2pt,width=3pt]}]H)    % Right, out (wedge)
        }};
        \node[align=center] at (0,-2.0) {Groupes verticaux (\ce{COOH}, \ce{R}) : \textbf{rentrent} (pointillés)\\
        Groupes horizontaux (\ce{NH2}, \ce{H}) : \textbf{sortent} (pleins)};
    \end{scope}
\end{tikzpicture}
\legende{Correspondance 3D de la projection de Fischer (L-configuration). En Fischer, les groupes horizontaux sortent (représentés en traits pleins épais), les groupes verticaux rentrent (pointillés). En Cram, on retrouve cette géométrie : liaisons pointillées vers l'arrière, liaisons en triangle vers l'avant.}
\end{popfigure}

\courssub{Exemple résolu : de la formule à la projection de Fischer}

\begin{exoresolu}[Nommage et projection de Fischer de la Valine]
\textbf{Énoncé :} La valine a pour formule semi-développée $\ce{(CH3)2CH-CH(NH2)-COOH}$.
\begin{enumerate}
    \item Donner son nom systématique.
    \item Donner son abréviation à trois lettres.
    \item Dessiner la projection de Fischer de la \textbf{L-valine}.
    \item La L-valine est-elle un énantiomère de la D-valine ? Justifier.
\end{enumerate}

\textbf{Solution détaillée :}
\begin{enumerate}
    \item \textbf{Nom systématique :} La chaîne principale contenant le $\ce{COOH}$ a 4 atomes de carbone (butanoïque). Le groupe amine est en C2. Il y a un groupe méthyle en C3. $\rightarrow$ \textbf{Acide 2-amino-3-méthylbutanoïque}.
    \item \textbf{Abréviation :} \textbf{Val}.
    \item \textbf{Projection de Fischer de la L-valine :}
    \begin{itemize}
        \item Étape 1 : Chaîne verticale, $\ce{COOH}$ en haut.
        \item Étape 2 : Chaîne latérale $\ce{R = -CH(CH3)2}$ (isopropyle) en bas.
        \item Étape 3 : Configuration L $\rightarrow$ $\ce{NH2}$ à \textbf{GAUCHE}, $\ce{H}$ à \textbf{DROITE}.
    \end{itemize}
    \item \textbf{Oui.} La L-valine et la D-valine sont deux stéréoisomères non superposables, images l'un de l'autre dans un miroir (carbone asymétrique unique, configuration inversée). Ce sont donc des \textbf{énantiomères}.
\end{enumerate}
\end{exoresolu}

\begin{popfigure}
\begin{tikzpicture}[scale=1.2]
    % Vertical bonds (going back) - dashed
    \draw[popdark, densely dashed, thick] (0,1.5) -- (0,0.3);
    \draw[popdark, densely dashed, thick] (0,-0.3) -- (0,-1.5);
    % Horizontal bonds (coming out) - wedges
    \draw[popblue, line width=3pt, -{Triangle[length=3pt,width=4pt]}] (-0.8,0) -- (0,0); % NH2 left, out
    \draw[poporange, line width=3pt, -{Triangle[length=3pt,width=4pt]}] (0,0) -- (0.8,0); % H right, out
    \fill[popdark] (0,0) circle (2pt);
    \node[above] at (0,1.5) {\ce{COOH}};
    \node[below] at (0,-1.5) {\ce{CH(CH3)2}};
    \node[left] at (-1,0) {\ce{NH2}};
    \node[right] at (1,0) {\ce{H}};
    \node[below=0.6cm] at (0,-1.5) {\textbf{L-Valine}};
\end{tikzpicture}
\legende{Projection de Fischer de la L-valine (configuration L : $\ce{NH2}$ à gauche).}
\end{popfigure}

\begin{attention}[Erreur classique à l'écrit : inversion COOH / R]
Si vous placez la chaîne latérale $\ce{R}$ en haut et le $\ce{COOH}$ en bas, la règle « $\ce{NH2}$ à gauche = L » \textbf{ne marche plus} (elle donne le résultat inverse). \textbf{Toujours} mettre le groupe le plus oxydé ($\ce{COOH}$, $\ce{CHO}$, $\ce{CO-R}$) en haut. C'est la règle absolue de la projection de Fischer.
\end{attention}

\begin{savaistu}[Le goût des énantiomères : une affaire de récepteurs chiraux]
Les récepteurs gustatifs de la langue sont des protéines chirales (faites d'acides aminés L). Ils reconnaissent sélectivement les énantiomères.
\begin{itemize}
    \item La \textbf{L-alanine} a un \textbf{goût sucré} (utilisée comme édulcorant dans certains aliments diététiques).
    \item La \textbf{D-alanine} a aussi un goût sucré, mais \textbf{moins intense}.
    \item La \textbf{D-phénylalanine} est \textbf{amère}, alors que la L-phénylalanine est sans goût particulier.
    \item Certains acides aminés D (ex. D-sérine) peuvent être \textbf{neurotoxiques} à haute dose car ils perturbent les récepteurs NMDA du cerveau.
\end{itemize}
C'est pourquoi l'industrie agroalimentaire et pharmaceutique (usines de Douala, Yaoundé) doit contrôler rigoureusement la pureté énantiomérique (excès énantiomérique $ee$) des acides aminés utilisés comme additifs (aspartame, glutamate monosodique) ou principes actifs.
\end{savaistu}

\courssub{Exercices d'application}

\begin{exercice}[Entraînement : Sérine et Leucine]
\begin{enumerate}
    \item Écrire la formule semi-développée de la \textbf{sérine} ($\ce{R = -CH2OH}$) et de la \textbf{leucine} ($\ce{R = -CH2-CH(CH3)2}$).
    \item Donner leur nom systématique IUPAC.
    \item Dessiner la projection de Fischer de la \textbf{L-sérine} et de la \textbf{L-leucine}.
    \item Pourquoi la glycine n'a-t-elle pas de configuration L ou D ?
\end{enumerate}
\end{exercice}

\begin{corrige}[Exercice n°1]
\begin{enumerate}
    \item \textbf{Sérine :} $\ce{HO-CH2-CH(NH2)-COOH}$. \textbf{Leucine :} $\ce{(CH3)2CH-CH2-CH(NH2)-COOH}$.
    \item \textbf{Sérine :} Acide 2-amino-3-hydroxypropanoïque. \textbf{Leucine :} Acide 2-amino-4-méthylpentanoïque.
    \item \textbf{L-Sérine :} $\ce{COOH}$ haut, $\ce{CH2OH}$ bas, $\ce{NH2}$ gauche, $\ce{H}$ droite. \\
    \textbf{L-Leucine :} $\ce{COOH}$ haut, $\ce{CH2-CH(CH3)2}$ bas, $\ce{NH2}$ gauche, $\ce{H}$ droite.
    \item La glycine a $\ce{R = H}$. Son carbone $\alpha$ porte deux atomes d'hydrogène ($\ce{H}$ et $\ce{R}$). Il n'est donc \textbf{pas asymétrique} (pas de chiralité). Un seul stéréoisomère existe.
\end{enumerate}
\end{corrige}

\begin{aretenir}[Bilan de la section]
\begin{itemize}
    \item Nom systématique : \textbf{Acide 2-aminoalcanoïque}.
    \item Noms usuels + abréviations 3 lettres (Gly, Ala, Val, Leu, Ser, Cys) à connaître par cœur.
    \item Projection de Fischer : Chaîne verticale, $\ce{COOH}$ en haut. \textbf{Horizontaux = Sortent}, \textbf{Verticaux = Rentrent}.
    \item Configuration \textbf{L} : $\ce{NH2}$ à \textbf{GAUCHE}. Configuration \textbf{D} : $\ce{NH2}$ à \textbf{DROITE}.
    \item Acides aminés protéiques = Série \textbf{L} (sauf glycine achirale).
\end{itemize}
\end{aretenir}

\coursec{Propriétés acido-basiques : l'amphion}

Dans la section précédente, nous avons appris à nommer les acides $\alpha$-aminés et à les représenter en projection de Fischer, en distinguant les configurations \cle{L} et \cle{D}. Mais une question fondamentale reste en suspens : sous quelle forme un acide $\alpha$-aminé existe-t-il \emph{réellement} ? La formule $\ce{H2N-CHR-COOH}$ que nous avons écrite jusqu'ici cache une réalité surprenante : la molécule « neutre » est presque un fantôme. C'est ce que nous allons découvrir maintenant, et c'est l'une des notions les plus importantes de cette leçon pour le Baccalauréat TI.

\courssub{Une molécule à double visage : acide et base à la fois}

\textbf{Intuition.} Imagine une pièce où se trouvent face à face un donneur compulsif et un receveur insatiable. Le donneur ne peut pas s'empêcher de donner, le receveur ne peut pas s'empêcher de prendre. Que va-t-il se passer ? Un transfert, inévitablement. C'est exactement la situation à l'intérieur d'une molécule d'acide $\alpha$-aminé.

Regardons attentivement les deux groupes fonctionnels portés par le même carbone $\alpha$ :

\begin{itemize}
\item Le groupe \cle{carboxyle} $\ce{-COOH}$ est un \textbf{acide} au sens de Brønsted : c'est un \emph{donneur de proton} $\ce{H+}$. On connaît bien ce comportement depuis l'étude des acides carboxyliques :
\[ \ce{R-COOH + H2O <=> R-COO- + H3O+} \]
\item Le groupe \cle{amine} $\ce{-NH2}$ est une \textbf{base} au sens de Brønsted : c'est un \emph{accepteur de proton} $\ce{H+}$, grâce au doublet non liant de l'atome d'azote :
\[ \ce{R-NH2 + H2O <=> R-NH3+ + HO-} \]
\end{itemize}

\begin{aretenir}[L'essentiel sur la double nature]
Un acide $\alpha$-aminé possède \textbf{à la fois} un groupe acide ($\ce{-COOH}$, donneur de $\ce{H+}$) et un groupe basique ($\ce{-NH2}$, accepteur de $\ce{H+}$). Ces deux groupes étant portés par la \emph{même} molécule, ils peuvent réagir \emph{l'un sur l'autre}.
\end{aretenir}

\courssub{Le transfert interne de proton}

\textbf{Intuition.} Puisque le donneur ($\ce{-COOH}$) et le receveur ($\ce{-NH2}$) sont dans la même molécule, séparés par un seul atome de carbone, le transfert n'a pas besoin d'un intermédiaire : le proton passe \emph{directement} du groupe carboxyle au groupe amine. On parle de \cle{transfert interne de proton} (ou transfert intramoléculaire).

Le groupe carboxyle cède son proton : $\ce{-COOH -> -COO-}$. Le groupe amine le capte : $\ce{-NH2 -> -NH3+}$. Le résultat est un ion remarquable.

\begin{definition}[L'amphion (ou zwitterion)]
L'\cle{amphion}, appelé aussi \cle{zwitterion} ou \emph{ion dipolaire}, est l'ion de formule générale $\ce{^{+}H3N-CHR-COO^{-}}$ résultant du \textbf{transfert interne d'un proton} du groupe carboxyle $\ce{-COOH}$ vers le groupe amine $\ce{-NH2}$ d'un acide $\alpha$-aminé. C'est un ion \textbf{globalement neutre} (charge électrique totale nulle) qui porte \textbf{deux charges opposées} : une charge positive sur l'azote et une charge négative sur l'oxygène.
\end{definition}

\begin{popfigure}
\begin{center}
\chemfig{H_2N-[:30](-[2]R)(-[6]H)-[:330]C(=[2]O)(-[6]OH)}
\quad
{\color{popink}\Large $\rightleftharpoons$}
\quad
\chemfig{\chemabove{N}{\scriptstyle\oplus}(-[6]H)(-[2]H)-[:30](-[2]R)(-[6]H)-[:330]C(=[2]O)(-[6]\chembelow{O}{\scriptstyle\ominus})}
\end{center}
\legende{Transfert interne de proton : la forme moléculaire $\ce{H2N-CHR-COOH}$ (à gauche) se transforme en amphion $\ce{^{+}H3N-CHR-COO^{-}}$ (à droite). Le proton $\ce{H+}$ migre de l'oxygène du groupe carboxyle vers l'azote du groupe amine.}
\end{popfigure}

\begin{methode}[Écrire l'équilibre forme moléculaire / forme ionique]
Pour écrire correctement l'équilibre entre la forme moléculaire et la forme ionique d'un acide $\alpha$-aminé :
\begin{enumerate}
\item Écris la \textbf{forme moléculaire} à gauche : $\ce{H2N-CHR-COOH}$ (aucune charge).
\item Place une \textbf{double flèche} $\rightleftharpoons$ : il s'agit d'un \emph{équilibre chimique}, pas d'une réaction totale.
\item Écris la \textbf{forme ionique} (amphion) à droite : $\ce{^{+}H3N-CHR-COO^{-}}$.
\item Vérifie le \textbf{placement des charges} : le $\oplus$ est sur l'azote (qui a gagné un proton), le $\ominus$ est sur l'oxygène (qui a perdu un proton).
\item Vérifie la \textbf{neutralité globale} : une charge $+$ et une charge $-$ se compensent, la charge totale est nulle des deux côtés de la flèche.
\end{enumerate}
\end{methode}

L'équation-bilan de l'équilibre s'écrit donc :

\[ \ce{H2N-CHR-COOH <=> ^{+}H3N-CHR-COO^{-}} \]

\begin{attention}[Le piège classique du Bac : les charges inversées]
L'erreur la plus fréquente est d'écrire l'amphion avec les charges inversées, par exemple $\ce{^{-}H2N-CHR-COOH^{+}}$. C'est \textbf{doublement faux} :
\begin{itemize}
\item C'est bien l'\textbf{azote} qui porte la charge \textbf{positive} : il a \emph{gagné} un proton ($\ce{-NH2 + H+ -> -NH3+}$).
\item C'est bien l'\textbf{oxygène} qui porte la charge \textbf{négative} : le groupe carboxyle a \emph{perdu} son proton ($\ce{-COOH -> -COO- + H+}$).
\end{itemize}
Retiens l'image : le proton \emph{voyage} de l'oxygène vers l'azote. Celui qui \emph{reçoit} devient positif, celui qui \emph{donne} devient négatif. De plus, un groupe $\ce{-COOH^{+}}$ ou $\ce{-NH2^{-}}$ n'a aucune réalité chimique dans ce contexte.
\end{attention}

\courssub{Exemple détaillé : l'amphion de l'alanine}

Prenons l'alanine ($R = \ce{CH3}$), l'acide $\alpha$-aminé le plus simple après la glycine. Sa forme moléculaire est $\ce{H2N-CH(CH3)-COOH}$. Après transfert interne du proton, on obtient :

\begin{popfigure}
\begin{center}
\chemfig{\chemabove{N}{\scriptstyle\oplus}(-[6]H)(-[2]H)-[:30](-[2]CH_3)(-[6]H)-[:330]C(=[2]O)(-[6]\chembelow{O}{\scriptstyle\ominus})}
\end{center}
\legende{Amphion de l'alanine : $\ce{^{+}H3N-CH(CH3)-COO^{-}}$. La charge $\oplus$ est portée par l'azote de l'ancien groupe amine devenu ammonium $\ce{-NH3+}$ ; la charge $\ominus$ est portée par l'oxygène du groupe carboxylate $\ce{-COO-}$. La charge électrique totale de l'ion est nulle.}
\end{popfigure}

L'équation-bilan de l'équilibre pour l'alanine s'écrit :

\[ \ce{H2N-CH(CH3)-COOH <=> ^{+}H3N-CH(CH3)-COO^{-}} \]

Commentaire : cet équilibre est un équilibre \emph{acido-basique intramoléculaire}, quasi instantané, qui ne nécessite ni catalyseur ni chauffage. Il est \textbf{très fortement déplacé vers la droite} : en solution aqueuse comme à l'état solide, l'alanine existe presque exclusivement sous forme d'amphion.

\courssub{Où se trouve l'équilibre ? La forme amphion domine}

\textbf{Intuition.} Un transfert de proton entre un acide carboxylique ($pK_a \approx 2{,}3$ pour le couple $\ce{-COOH}/\ce{-COO-}$) et une amine (dont l'acide conjugué $\ce{-NH3+}$ a un $pK_a \approx 9{,}7$) est thermodynamiquement très favorable : l'acide le plus fort cède son proton à la base la plus forte. L'équilibre est donc massivement déplacé vers l'amphion.

\begin{propriete}[Prédominance de l'amphion]
En \textbf{solution aqueuse} comme à l'\textbf{état solide}, un acide $\alpha$-aminé existe \textbf{essentiellement sous forme d'amphion} $\ce{^{+}H3N-CHR-COO^{-}}$. La forme moléculaire $\ce{H2N-CHR-COOH}$ est pratiquement absente. C'est pourquoi on dit que les acides $\alpha$-aminés sont des \cle{sels internes} : ils se comportent comme des sels, alors que le cation et l'anion appartiennent à la même entité.
\end{propriete}

\begin{exoresolu}[Équilibre de la glycine et bilan de charge]
\textbf{Énoncé.} La glycine (ou acide aminoéthanoïque) a pour formule semi-développée $\ce{H2N-CH2-COOH}$.
\begin{enumerate}
\item Écris l'équation-bilan de l'équilibre entre sa forme moléculaire et sa forme ionique.
\item On dissout $m = \SI{7,5}{g}$ de glycine ($M = \SI{75}{g.mol^{-1}}$) dans de l'eau distillée de façon à obtenir $V = \SI{200}{mL}$ de solution. En supposant la glycine entièrement sous forme d'amphion, calcule la concentration molaire en amphion dans la solution.
\item La solution obtenue conduit-elle le courant électrique ? Justifie.
\end{enumerate}
\tcblower
\textbf{Solution détaillée.}
\begin{enumerate}
\item Le transfert interne de proton s'écrit avec une double flèche (équilibre) :
\[ \ce{H2N-CH2-COOH <=> ^{+}H3N-CH2-COO^{-}} \]
La charge $\oplus$ est sur l'azote, la charge $\ominus$ sur l'oxygène ; la charge totale est nulle de part et d'autre.
\item Quantité de matière de glycine :
\[ n = \frac{m}{M} = \frac{7{,}5}{75} = \SI{0,10}{mol} \]
Concentration de la solution :
\[ C = \frac{n}{V} = \frac{0{,}10}{0{,}200} = \SI{0,50}{mol.L^{-1}} \]
La glycine étant entièrement sous forme d'amphion : $[\text{amphion}] = \SI{0,50}{mol.L^{-1}}$.
\item \textbf{Oui}, la solution conduit le courant électrique. Bien que chaque amphion soit globalement neutre, il s'agit d'un \emph{ion dipolaire} porteur de deux charges localisées : les amphions sont mobiles et hydratés en solution, ils assurent le transport des charges sous l'effet d'un champ électrique alternatif, exactement comme les ions d'un sel dissous. C'est une preuve expérimentale directe de l'existence de la forme ionique.
\end{enumerate}
\end{exoresolu}

\courssub{Conséquences physiques : des molécules qui se comportent comme des sels}

La structure ionique de l'amphion explique des propriétés physiques tout à fait inhabituelles pour des composés organiques de petite taille.

\textbf{Des points de fusion élevés.} Un alcane, un alcool ou un acide carboxylique de masse molaire comparable fond facilement (souvent en dessous de $\SI{100}{\celsius}$). Or les acides $\alpha$-aminés fondent (en se décomposant) à des températures très élevées : l'alanine se décompose vers $\SI{297}{\celsius}$, la glycine vers $\SI{233}{\celsius}$. Pourquoi ? À l'état solide, les amphions s'organisent en \textbf{réseau cristallin ionique} : les charges $\oplus$ d'un amphion attirent les charges $\ominus$ des amphions voisins, comme dans un cristal de sel de cuisine $\ce{NaCl}$. Briser ces attractions électrostatiques demande beaucoup d'énergie.

\textbf{Une bonne solubilité dans l'eau.} L'eau, solvant polaire, solvate très efficacement les espèces chargées : les molécules d'eau entourent les pôles $\ce{-NH3+}$ et $\ce{-COO-}$ de l'amphion (hydratation). Les acides $\alpha$-aminés sont donc nettement plus solubles dans l'eau que les molécules organiques neutres de taille équivalente, et très peu solubles dans les solvants organiques non polaires (éther, benzène).

\textbf{Le caractère de sels internes.} Ces deux constats (fusion élevée, solubilité aqueuse) sont la signature des composés \emph{ioniques}. On qualifie les acides $\alpha$-aminés de \cle{sels internes} : tout se passe comme si la molécule avait formé un sel \emph{avec elle-même}.

\begin{center}
\begin{tabularx}{\linewidth}{|l|X|X|}
\hline
\rowcolor{popblueL} \textbf{Propriété} & \textbf{Acide $\alpha$-aminé (amphion)} & \textbf{Composé organique neutre comparable} \\
\hline
État à \SI{25}{\celsius} & Solide cristallisé & Souvent liquide ou solide fusible \\
\hline
Point de fusion & Élevé (souvent $>\SI{200}{\celsius}$, avec décomposition) & Faible ou modéré \\
\hline
Solubilité dans l'eau & Bonne (espèce chargée, hydratée) & Faible à moyenne \\
\hline
Solubilité dans l'éther & Très faible & Bonne \\
\hline
Conduction électrique en solution & Oui (ions dipolaires mobiles) & Non \\
\hline
\end{tabularx}
\end{center}

\begin{savaistu}[D'où vient le mot « zwitterion » ?]
Le mot \emph{zwitterion} vient de l'allemand \emph{Zwitter}, qui signifie « hybride », « bâtard », « être double ». L'amphion est littéralement un \textbf{hybride entre un cation et un anion} : mi-cation ($\ce{-NH3+}$), mi-anion ($\ce{-COO-}$), réunis dans une seule et même espèce. Ce terme a été introduit par les chimistes allemands à la fin du \textsc{xix}\textsuperscript{e} siècle, lorsqu'ils ont découvert avec étonnement que les acides aminés se comportaient comme des sels. Cette double nature ionique est d'ailleurs exploitée au quotidien : c'est elle qui permet aux protéines de nos aliments (viande, haricots niébé, arachides) de se dissoudre partiellement lors de la cuisson et d'être digérées par nos enzymes.
\end{savaistu}

\begin{aretenir}[Ce qu'il faut retenir de cette section]
\begin{itemize}
\item Un acide $\alpha$-aminé possède un groupe \textbf{acide} ($\ce{-COOH}$) et un groupe \textbf{basique} ($\ce{-NH2}$) sur la même molécule.
\item Il se produit un \textbf{transfert interne de proton} du carboxyle vers l'amine, donnant l'\textbf{amphion} $\ce{^{+}H3N-CHR-COO^{-}}$, ion globalement neutre portant deux charges opposées.
\item L'équilibre $\ce{H2N-CHR-COOH <=> ^{+}H3N-CHR-COO^{-}}$ est très déplacé vers l'amphion : c'est la forme \textbf{majoritaire} en solution et à l'état solide.
\item Conséquences : points de fusion élevés, bonne solubilité dans l'eau, comportement de \textbf{sels internes}.
\item Piège à éviter : le $\oplus$ est sur l'\textbf{azote}, le $\ominus$ sur l'\textbf{oxygène} — jamais l'inverse.
\end{itemize}
\end{aretenir}

Nous savons désormais que l'acide $\alpha$-aminé vit sous la forme d'un ion dipolaire. Mais cet amphion garde un secret : il peut encore \emph{donner} un proton ou \emph{en capter} un selon le milieu dans lequel on le plonge. C'est ce comportement \emph{amphotère} que nous étudierons dans la section suivante.

\coursec{Comportement amphotère de l'amphion}

Dans la section précédente, nous avons découvert que l'acide $\alpha$-aminé n'existe pratiquement jamais sous sa forme moléculaire \chemfig{H_2N-CHR-COOH} : par transfert interne de proton, il adopte la forme ionique appelée \cle{amphion} (ou zwitterion), \chemfig{^{+}H_3N-CHR-COO^{-}}, qui porte à la fois une charge positive et une charge négative. Une question surgit alors naturellement : que se passe-t-il si l'on plonge cet amphion dans un milieu acide, comme le suc gastrique, ou dans un milieu basique, comme une lessive de savonnerie artisanale ? La réponse est remarquable : l'amphion sait \emph{s'adapter} au milieu. C'est tout l'objet de cette section.

\courssub{Un ampholyte : l'intuition d'abord}

Imagine un élève bilingue qui parle français avec les francophones et anglais avec les anglophones : il change de « comportement » selon son interlocuteur. L'amphion fait de même avec les protons. Il possède deux groupes de nature opposée :

\begin{itemize}
\item le groupe \chemfig{COO^{-}} (carboxylate), qui est une \textbf{base} : il peut \emph{capter} un proton \ce{H+} pour redonner \chemfig{COOH} ;
\item le groupe \chemfig{NH_3^{+}} (ammonium), qui est un \textbf{acide} : il peut \emph{céder} un proton \ce{H+} pour redonner \chemfig{NH_2}.
\end{itemize}

L'amphion peut donc, selon le milieu, jouer le rôle de base ou le rôle d'acide. On dit que c'est un \cle{ampholyte} (ou espèce \cle{amphotère}).

\begin{definition}[Ampholyte]
Un \cle{ampholyte} (ou espèce amphotère) est une espèce chimique qui peut se comporter \textbf{soit comme un acide} (en cédant un proton \ce{H+}), \textbf{soit comme une base} (en captant un proton \ce{H+}), selon la nature du milieu dans lequel elle se trouve. L'amphion \chemfig{^{+}H_3N-CHR-COO^{-}} d'un acide $\alpha$-aminé est un ampholyte, tout comme l'ion hydrogénocarbonate \ce{HCO3-} ou l'eau \ce{H2O}.
\end{definition}

\begin{aretenir}[L'idée directrice]
Selon le pH du milieu, un acide $\alpha$-aminé existe sous \textbf{trois formes} différentes :
\begin{itemize}
\item en milieu \textbf{acide} : forme \cle{cationique} \chemfig{^{+}H_3N-CHR-COOH} ;
\item en milieu \textbf{neutre} (ou proche de la neutralité) : forme \cle{amphion} \chemfig{^{+}H_3N-CHR-COO^{-}} ;
\item en milieu \textbf{basique} : forme \cle{anionique} \chemfig{H_2N-CHR-COO^{-}}.
\end{itemize}
\end{aretenir}

\courssub{Caractère basique de l'amphion : le comportement en milieu acide}

Plaçons l'amphion en milieu acide, c'est-à-dire en présence d'ions oxonium \ce{H3O+} (ou, pour simplifier l'écriture, de protons \ce{H+}). Quel groupe va réagir ? C'est le groupe qui a « envie » de proton, c'est-à-dire la base du couple : le groupe carboxylate \chemfig{COO^{-}}. Il capte un proton et redevient un groupe carboxyle \chemfig{COOH}.

L'équation-bilan s'écrit :

\[ \ce{^{+}H3N-CHR-COO- + H+ -> ^{+}H3N-CHR-COOH} \]

ou, de façon plus rigoureuse, avec l'ion oxonium :

\[ \ce{^{+}H3N-CHR-COO- + H3O+ -> ^{+}H3N-CHR-COOH + H2O} \]

\begin{exemplebox}[Lecture de l'équation]
\begin{itemize}
\item \textbf{Réactifs} : l'amphion et un proton (apporté par l'acide du milieu).
\item \textbf{Produit} : un \cle{cation}, car la molécule ne porte plus que la charge positive du groupe \chemfig{NH_3^{+}} ; la charge négative a disparu lors de la capture du proton.
\item \textbf{Caractéristiques} : réaction acido-basique, donc \emph{rapide}, quasi \emph{totale} si l'acide ajouté est fort, et se déroulant à température ambiante, sans catalyseur.
\item \textbf{Rôle joué} : l'amphion capte un proton $\Rightarrow$ il se comporte comme une \textbf{base}.
\end{itemize}
\end{exemplebox}

Le couple acide/base mis en jeu ici est :

\[ \underbrace{\chemfig{^{+}H_3N-CHR-COOH}}_{\text{acide}} \;/\; \underbrace{\chemfig{^{+}H_3N-CHR-COO^{-}}}_{\text{base conjuguée}} \qquad \text{(couple } \ce{COOH}/\ce{COO-} \text{ de la molécule)} \]

\courssub{Caractère acide de l'amphion : le comportement en milieu basique}

Plaçons maintenant l'amphion en milieu basique, c'est-à-dire en présence d'ions hydroxyde \ce{OH-}. Cette fois, c'est le groupe qui possède un proton « disponible » qui réagit : le groupe ammonium \chemfig{NH_3^{+}}. Il cède son proton à l'ion hydroxyde et redevient un groupe amine \chemfig{NH_2}.

L'équation-bilan s'écrit :

\[ \ce{^{+}H3N-CHR-COO- + OH- -> H2N-CHR-COO- + H2O} \]

\begin{exemplebox}[Lecture de l'équation]
\begin{itemize}
\item \textbf{Réactifs} : l'amphion et un ion hydroxyde (apporté par la base du milieu, par exemple la soude).
\item \textbf{Produits} : un \cle{anion}, car la molécule ne porte plus que la charge négative du groupe \chemfig{COO^{-}} ; et une molécule d'eau.
\item \textbf{Caractéristiques} : réaction acido-basique, \emph{rapide} et quasi \emph{totale}, à température ambiante, sans catalyseur.
\item \textbf{Rôle joué} : l'amphion cède un proton $\Rightarrow$ il se comporte comme un \textbf{acide}.
\end{itemize}
\end{exemplebox}

Le couple acide/base mis en jeu ici est :

\[ \underbrace{\chemfig{^{+}H_3N-CHR-COO^{-}}}_{\text{acide}} \;/\; \underbrace{\chemfig{H_2N-CHR-COO^{-}}}_{\text{base conjuguée}} \qquad \text{(couple } \ce{NH3+}/\ce{NH2} \text{ de la molécule)} \]

Remarque bien que l'amphion apparaît dans \textbf{deux couples} : comme base dans le premier, comme acide dans le second. C'est précisément la signature d'un ampholyte.

\begin{methode}[Écrire correctement les équations-bilans]
Pour écrire l'équation-bilan de la réaction de l'amphion avec un acide ou une base, procède en quatre étapes :
\begin{enumerate}
\item \textbf{Identifie le milieu} : acide (présence de \ce{H+} ou \ce{H3O+}) ou basique (présence de \ce{OH-}).
\item \textbf{Choisis le bon groupe réactif} :
\begin{itemize}
\item milieu acide $\Rightarrow$ l'amphion joue le rôle de \textbf{base} : c'est le groupe \chemfig{COO^{-}} qui \emph{captе} \ce{H+} ;
\item milieu basique $\Rightarrow$ l'amphion joue le rôle d'\textbf{acide} : c'est le groupe \chemfig{NH_3^{+}} qui \emph{cède} \ce{H+}.
\end{itemize}
\item \textbf{Écris l'équation} en ne modifiant \emph{que} le groupe qui réagit ; le reste de la molécule est spectateur.
\item \textbf{Vérifie la conservation} : mêmes atomes et même charge totale des deux côtés de la flèche. Par exemple, dans $\ce{^{+}H3N-CHR-COO- + H+ -> ^{+}H3N-CHR-COOH}$, la charge totale est $+1$ à gauche ($0 + 1$) et $+1$ à droite : l'équation est équilibrée.
\end{enumerate}
\end{methode}

\begin{attention}[L'erreur classique : faire réagir le mauvais groupe]
Beaucoup d'élèves, en milieu acide, font céder un proton par \chemfig{NH_3^{+}} alors que le milieu \emph{apporte} des protons ! Retiens cette logique simple :
\begin{itemize}
\item en milieu \textbf{acide} (excès de \ce{H+}), on ne peut que \emph{capter} des protons : c'est \chemfig{COO^{-}} qui réagit $\Rightarrow$ on obtient le \textbf{cation} ;
\item en milieu \textbf{basique} (excès de \ce{OH-}), on ne peut que \emph{céder} des protons : c'est \chemfig{NH_3^{+}} qui réagit $\Rightarrow$ on obtient l'\textbf{anion}.
\end{itemize}
Autre piège : ne jamais écrire la forme moléculaire \chemfig{H_2N-CHR-COOH} comme produit de ces réactions. Elle n'existe pratiquement pas en solution ; les seules formes réelles sont le cation, l'amphion et l'anion.
\end{attention}

\courssub{Schéma général : les trois formes de l'acide $\alpha$-aminé}

Les deux équilibres précédents se résument dans le schéma suivant, à connaître parfaitement pour le Baccalauréat :

\begin{popfigure}
\begin{center}
\begin{tikzpicture}[node distance=1.1cm and 2.6cm]
\node[draw=popblue, fill=popblueL, rounded corners, thick, align=center, inner sep=6pt] (cat) {\chemfig{^{+}H_3N-CHR-COOH}\\[2pt]\textbf{cation}\\ (milieu acide)};
\node[draw=poppurple, fill=poppurpleL, rounded corners, thick, align=center, inner sep=6pt, right=of cat] (amph) {\chemfig{^{+}H_3N-CHR-COO^{-}}\\[2pt]\textbf{amphion}\\ (pH intermédiaire)};
\node[draw=poporange, fill=poporangeL, rounded corners, thick, align=center, inner sep=6pt, right=of amph] (an) {\chemfig{H_2N-CHR-COO^{-}}\\[2pt]\textbf{anion}\\ (milieu basique)};
\draw[-{Stealth}, very thick, popblue] (cat.15) -- node[above, font=\small]{\chemfig{OH^{-}}} (amph.165);
\draw[-{Stealth}, very thick, popblue] (amph.195) -- node[below, font=\small]{\chemfig{H^{+}}} (cat.-15);
\draw[-{Stealth}, very thick, poporange] (amph.15) -- node[above, font=\small]{\chemfig{OH^{-}}} (an.165);
\draw[-{Stealth}, very thick, poporange] (an.195) -- node[below, font=\small]{\chemfig{H^{+}}} (amph.-15);
\end{tikzpicture}
\end{center}
\legende{Les trois formes d'un acide $\alpha$-aminé en solution aqueuse. Chaque ajout de \ce{OH-} arrache un proton ; chaque ajout de \ce{H+} en redonne un.}
\end{popfigure}

\begin{exoresolu}[Réactions de l'amphion de l'alanine]
\textbf{Énoncé.} On dissout de l'alanine ($\mathrm{R = CH_3}$) dans l'eau : elle s'y trouve sous forme d'amphion \chemfig{^{+}H_3N-CH(CH_3)-COO^{-}}.
\begin{enumerate}
\item On ajoute quelques gouttes d'acide chlorhydrique concentré. Écris l'équation-bilan de la réaction observée et nomme l'espèce obtenue.
\item On repart de la solution initiale et on ajoute cette fois de la soude. Écris l'équation-bilan et nomme l'espèce obtenue.
\item Conclus sur la nature chimique de l'amphion.
\end{enumerate}
\tcblower
\textbf{Solution détaillée.}
\begin{enumerate}
\item L'acide chlorhydrique apporte des ions \ce{H3O+} : le milieu devient acide. L'amphion joue le rôle de \textbf{base} ; c'est son groupe \chemfig{COO^{-}} qui capte un proton :
\[ \ce{^{+}H3N-CH(CH3)-COO- + H3O+ -> ^{+}H3N-CH(CH3)-COOH + H2O} \]
L'espèce obtenue est le \textbf{cation} de l'alanine (charge globale $+1$).
\item La soude apporte des ions \ce{OH-} : le milieu devient basique. L'amphion joue le rôle d'\textbf{acide} ; c'est son groupe \chemfig{NH_3^{+}} qui cède un proton :
\[ \ce{^{+}H3N-CH(CH3)-COO- + OH- -> H2N-CH(CH3)-COO- + H2O} \]
L'espèce obtenue est l'\textbf{anion} de l'alanine (charge globale $-1$).
\item L'amphion de l'alanine réagit aussi bien avec un acide qu'avec une base : c'est un \textbf{ampholyte} (espèce amphotère). Cette propriété est commune à tous les acides $\alpha$-aminés.
\end{enumerate}
\end{exoresolu}

\courssub{Diagramme de prédominance : quelle forme selon le pH ?}

Chacun des deux couples possède une constante d'acidité, caractérisée par son $pK_a$ :
\begin{itemize}
\item couple cation/amphion : $\ce{^{+}H3N-CHR-COOH}/\ce{^{+}H3N-CHR-COO-}$, de $pK_{a1}$ voisin de $2$ à $2{,}5$ ;
\item couple amphion/anion : $\ce{^{+}H3N-CHR-COO-}/\ce{H2N-CHR-COO-}$, de $pK_{a2}$ voisin de $9$ à $10$.
\end{itemize}

Pour l'\textbf{alanine}, les valeurs exactes sont $pK_{a1} = 2{,}3$ et $pK_{a2} = 9{,}7$. On en déduit le diagramme de prédominance :

\begin{popfigure}
\begin{center}
\begin{tikzpicture}
\begin{axis}[
width=0.95\linewidth, height=5.2cm,
xmin=0, xmax=14, ymin=0, ymax=1,
xlabel={pH}, xlabel style={font=\small},
ytick=\empty, xtick={0,2.3,6,9.7,14},
axis lines=left, clip=false,
]
\addplot[draw=none, fill=popblueL, domain=0:2.3] {1} \closedcycle;
\addplot[draw=none, fill=poppurpleL, domain=2.3:9.7] {1} \closedcycle;
\addplot[draw=none, fill=poporangeL, domain=9.7:14] {1} \closedcycle;
\addplot[popcurve, domain=0:14, samples=2] {0};
\draw[thick, popblue, dashed] (axis cs:2.3,0) -- (axis cs:2.3,1);
\draw[thick, poporange, dashed] (axis cs:9.7,0) -- (axis cs:9.7,1);
\node[align=center, font=\small, popdark] at (axis cs:1.15,0.5) {\textbf{cation}\\ \chemfig{^{+}H_3N-CHR-COOH}};
\node[align=center, font=\small, popdark] at (axis cs:6,0.5) {\textbf{amphion}\\ \chemfig{^{+}H_3N-CHR-COO^{-}}};
\node[align=center, font=\small, popdark] at (axis cs:11.85,0.5) {\textbf{anion}\\ \chemfig{H_2N-CHR-COO^{-}}};
\node[font=\footnotesize, popblue, below] at (axis cs:2.3,0) {$pK_{a1}=2{,}3$};
\node[font=\footnotesize, poporange, above] at (axis cs:9.7,1) {$pK_{a2}=9{,}7$};
\end{axis}
\end{tikzpicture}
\end{center}
\legende{Diagramme de prédominance simplifié des trois formes de l'alanine en fonction du pH.}
\end{popfigure}

\begin{exemplebox}[Exemple chiffré : l'alanine selon le pH (complément Bac)]
Pour l'alanine ($pK_{a1} = 2{,}3$ ; $pK_{a2} = 9{,}7$) :
\begin{itemize}
\item à $\mathrm{pH} < 2{,}3$ (par exemple $\mathrm{pH} = 1$, comme dans l'estomac) : la \textbf{forme cationique} \chemfig{^{+}H_3N-CH(CH_3)-COOH} prédomine ;
\item autour de $\mathrm{pH} = 6$ (entre les deux $pK_a$) : l'\textbf{amphion} \chemfig{^{+}H_3N-CH(CH_3)-COO^{-}} est très largement majoritaire ; c'est la situation dans la plupart des cellules ;
\item à $\mathrm{pH} > 9{,}7$ (par exemple $\mathrm{pH} = 12$) : la \textbf{forme anionique} \chemfig{H_2N-CH(CH_3)-COO^{-}} prédomine.
\end{itemize}
Au voisinage du pH intermédiaire $pH_i = \dfrac{pK_{a1}+pK_{a2}}{2} = \dfrac{2{,}3+9{,}7}{2} = 6{,}0$ (appelé \emph{point isoélectrique}), l'espèce majoritaire est l'amphion, de charge globale nulle.
\end{exemplebox}

\begin{aretenir}[Bilan de la section]
\begin{itemize}
\item L'amphion est un \cle{ampholyte} : il capte un proton en milieu acide (rôle de base) et en cède un en milieu basique (rôle d'acide).
\item Milieu acide : $\ce{^{+}H3N-CHR-COO- + H+ -> ^{+}H3N-CHR-COOH}$ (cation).
\item Milieu basique : $\ce{^{+}H3N-CHR-COO- + OH- -> H2N-CHR-COO- + H2O}$ (anion).
\item Deux couples sont mis en jeu : $\ce{^{+}H3N-CHR-COOH}/\ce{^{+}H3N-CHR-COO-}$ et $\ce{^{+}H3N-CHR-COO-}/\ce{H2N-CHR-COO-}$.
\item Selon le pH, l'acide $\alpha$-aminé existe sous forme \textbf{cationique}, d'\textbf{amphion} ou \textbf{anionique}.
\end{itemize}
\end{aretenir}

\begin{savaistu}[Des tampons biologiques qui protègent la vie]
Le comportement amphotère des acides aminés explique l'une de leurs fonctions biologiques les plus précieuses : ils agissent comme des \textbf{solutions tampons}. Puisqu'ils peuvent neutraliser aussi bien un acide qu'une base, ils s'opposent aux variations brutales de pH. Le pH du sang humain doit rester compris entre $7{,}35$ et $7{,}45$ : un écart de quelques dixièmes seulement peut être mortel. Les protéines (assemblages d'acides $\alpha$-aminés, comme nous le verrons avec la liaison peptidique) et l'hémoglobine participent activement à cette stabilisation, aux côtés du couple \ce{H2CO3}/\ce{HCO3-}. Chez les sportifs de l'équipe nationale d'athlétisme, l'effort intense produit de l'acide lactique : c'est en partie grâce à ces systèmes tampons que le pH sanguin ne s'effondre pas pendant la course !
\end{savaistu}

\begin{exercice}[Pour t'entraîner]
On considère la glycine ($\mathrm{R = H}$), le plus simple des acides $\alpha$-aminés, dont les constantes sont $pK_{a1} = 2{,}4$ et $pK_{a2} = 9{,}8$.
\begin{enumerate}
\item Écris l'équation-bilan de la réaction de l'amphion de la glycine avec l'ion oxonium \ce{H3O+}.
\item Écris l'équation-bilan de sa réaction avec l'ion hydroxyde \ce{OH-}.
\item Quelle forme prédomine à $\mathrm{pH} = 2$ ? À $\mathrm{pH} = 7$ ? À $\mathrm{pH} = 11$ ? Justifie à l'aide d'un diagramme de prédominance.
\end{enumerate}
\end{exercice}

\begin{corrige}[Exercice : glycine en solution]
\begin{enumerate}
\item En milieu acide, le groupe \chemfig{COO^{-}} capte un proton :
\[ \ce{^{+}H3N-CH2-COO- + H3O+ -> ^{+}H3N-CH2-COOH + H2O} \]
\item En milieu basique, le groupe \chemfig{NH_3^{+}} cède un proton :
\[ \ce{^{+}H3N-CH2-COO- + OH- -> H2N-CH2-COO- + H2O} \]
\item À $\mathrm{pH} = 2 < pK_{a1} = 2{,}4$ : le \textbf{cation} \chemfig{^{+}H_3N-CH_2-COOH} prédomine. À $\mathrm{pH} = 7$, compris entre $2{,}4$ et $9{,}8$ : l'\textbf{amphion} prédomine. À $\mathrm{pH} = 11 > pK_{a2} = 9{,}8$ : l'\textbf{anion} \chemfig{H_2N-CH_2-COO^{-}} prédomine. Le diagramme de prédominance est identique à celui de l'alanine, avec les frontières placées en $2{,}4$ et $9{,}8$.
\end{enumerate}
\end{corrige}

\coursec{La liaison peptidique et les dipeptides}

Après avoir étudié le comportement amphotère de l'amphion, qui explique la réactivité chimique des acides $\alpha$-aminés en solution aqueuse (équilibres acide/base, point isoélectrique), nous allons maintenant voir comment ces molécules s'assemblent entre elles. C'est la clé de la chimie des protéines : la fonction carboxyle d'un acide $\alpha$-aminé réagit sur la fonction amine d'un autre pour former une liaison covalente robuste, la \cle{liaison peptidique}. Ce processus, une condensation, libère une molécule d'eau et construit la chaîne carbonée des peptides et des protéines.

\courssub{Formation de la liaison peptidique : une réaction de condensation}

\paragraph{Principe chimique.}
Lorsqu'on met en présence deux acides $\alpha$-aminés, le groupe carboxyle ($-\ce{COOH}$) du premier — qui possède un caractère acide — peut réagir sur le groupe amine ($-\ce{NH2}$) du second — qui possède un caractère basique. Cette réaction est une \textbf{condensation} : les deux molécules s'unissent avec élimination d'une molécule d'eau ($\ce{H2O}$). Le groupement $-\ce{OH}$ du carboxyle et un atome d'hydrogène de l'amine partent former l'eau ; il reste une liaison covalente entre le carbone carbonylé du premier acide aminé et l'azote de l'amine du second.

\begin{definition}[Liaison peptidique]
La \textbf{liaison peptidique} est la liaison covalente \ce{-CO-NH-} (fonction amide) formée entre le groupe carboxyle d'un acide $\alpha$-aminé et le groupe amine d'un autre acide $\alpha$-aminé, avec élimination d'une molécule d'eau. Elle unit deux \textbf{résidus} d'acides $\alpha$-aminés.
\end{definition}

\paragraph{Équation-bilan générale.}
Soient deux acides $\alpha$-aminés de chaînes latérales R et R' (qui peuvent être identiques ou différentes). L'équation-bilan de la condensation s'écrit :
\[
\mathrm{H_2N\!-\!CHR\!-\!COOH} + \mathrm{H_2N\!-\!CHR'\!-\!COOH} \longrightarrow \mathrm{H_2N\!-\!CHR\!-\!CO\!-\!NH\!-\!CHR'\!-\!COOH} + \mathrm{H_2O}
\]
Le produit obtenu est un \textbf{dipeptide}. Il conserve une extrémité amine libre (extrémité \textbf{N-terminale}) et une extrémité carboxyle libre (extrémité \textbf{C-terminale}), ce qui lui confère encore un caractère amphotère et permet l'allongement de la chaîne.

\begin{methode}[Écrire l'équation-bilan de formation d'un dipeptide]
\begin{enumerate}
    \item Écrire les formules semi-développées des deux acides $\alpha$-aminés réactifs, en plaçant celui qui fournira l'extrémité N-terminale à gauche (son $\ce{NH2}$ ne réagit pas) et celui qui fournira l'extrémité C-terminale à droite (son $\ce{COOH}$ ne réagit pas).
    \item Identifier le groupe $\ce{-COOH}$ de l'acide aminé de gauche et le groupe $\ce{-NH2}$ de l'acide aminé de droite.
    \item Faire se rejoindre ces deux groupes : l'atome de carbone du $\ce{COOH}$ se lie à l'atome d'azote du $\ce{NH2}$.
    \item Éliminer une molécule d'eau $\ce{H2O}$ (un $\ce{H}$ du $\ce{NH2}$ et le $\ce{OH}$ du $\ce{COOH}$).
    \item Écrire la chaîne peptidique résultante en faisant apparaître la liaison \ce{-CO-NH-} (souvent notée \ce{-CONH-} ou \ce{-CO-NH-}).
    \item Vérifier la conservation des atomes (C, H, N, O) et de la charge.
\end{enumerate}
\end{methode}

\begin{popfigure}
\centering
\begin{tikzpicture}[node distance=1.5cm, >=Stealth, baseline=(ala.base)]
    % Reactant 1: Alanine
    \node (ala) {\chemfig{H_2N-[:30]CH(-[:90]CH_3)-[:30]C(=[:90]O)-[:-30]OH}};
    % Plus sign
    \node[right=1.2cm of ala] (plus) {\Large$+$};
    % Reactant 2: Glycine
    \node[right=1.2cm of plus] (gly) {\chemfig{H_2N-[:30]CH_2-[:30]C(=[:90]O)-[:-30]OH}};
    % Arrow
    \node[right=1.2cm of gly] (arrow) {\Large$\rightarrow$};
    % Product: Ala-Gly
    \node[right=1.2cm of arrow] (dipep) {\chemfig{H_2N-[:30]CH(-[:90]CH_3)-[:30]C(=[:90]O)-[:-30]N(-[:-90]H)-[:30]CH_2-[:30]C(=[:90]O)-[:-30]OH}};
    % Water
    \node[right=1.2cm of dipep] (water) {\Large$+$};
    \node[right=0.5cm of water] (h2o) {\chemfig{H-O-H}};
    
    % Labels under reactants
    \node[below=0.5cm of ala, font=\small\bfseries] {Alanine (Ala)};
    \node[below=0.5cm of gly, font=\small\bfseries] {Glycine (Gly)};
    \node[below=0.5cm of dipep, font=\small\bfseries] {Ala--Gly};
    \node[below=0.5cm of h2o, font=\small\bfseries] {Eau};
    
    % Highlight peptide bond in product using a background rectangle (no decorations library needed)
    \begin{scope}[on background layer]
        % We approximate the position of the peptide bond in the chemfig output (fragile but okay for visual)
        % Better: draw the peptide bond separately? No, keep chemfig simple.
        % Instead, add a label pointing to it.
    \end{scope}
    \draw[popblue, thick, <->] ($(dipep.south) + (-1.2cm, -0.8cm)$) -- ($(dipep.south) + (0.5cm, -0.8cm)$) node[midway, below, popblue, font=\small] {Liaison peptidique \ce{-CO-NH-}};
\end{tikzpicture}
\legende{Formation du dipeptide Ala–Gly par condensation entre l'alanine (à gauche) et la glycine (à droite). La liaison peptidique \ce{-CO-NH-} formée est indiquée par la flèche bleue ; une molécule d'eau est éliminée.}
\end{popfigure}

\begin{exoresolu}[Formation et masse molaire du dipeptide Ala–Gly]
\textbf{Énoncé.} L'alanine (Ala, formule brute $\ce{C3H7NO2}$, $M = \SI{89,1}{g.mol^{-1}}$) et la glycine (Gly, $\ce{C2H5NO2}$, $M = \SI{75,1}{g.mol^{-1}}$) réagissent pour former le dipeptide Ala–Gly. Écrire l'équation-bilan de la réaction et vérifier la conservation de la masse molaire.

\textbf{Solution détaillée.}
\begin{enumerate}
    \item \textbf{Équation-bilan :} L'alanine fournit l'extrémité N-terminale, la glycine l'extrémité C-terminale.
    \[
    \ce{H2N-CH(CH3)-COOH + H2N-CH2-COOH -> H2N-CH(CH3)-CO-NH-CH2-COOH + H2O}
    \]
    \item \textbf{Formule brute du dipeptide Ala–Gly :} On additionne les formules brutes des deux acides aminés et on soustrait $\ce{H2O}$.
    $\ce{C3H7NO2 + C2H5NO2 - H2O = C5H10N2O3}$.
    \item \textbf{Masse molaire théorique :} $M(\text{Ala–Gly}) = M(\text{Ala}) + M(\text{Gly}) - M(\ce{H2O})$.
    $M(\text{Ala–Gly}) = \SI{89,1}{} + \SI{75,1}{} - \SI{18,0}{} = \SI{146,2}{g.mol^{-1}}$.
    \item \textbf{Vérification :} Calcul direct depuis la formule brute $\ce{C5H10N2O3}$ :
    $5 \times \SI{12,0}{} + 10 \times \SI{1,0}{} + 2 \times \SI{14,0}{} + 3 \times \SI{16,0}{} = \SI{60,0}{} + \SI{10,0}{} + \SI{28,0}{} + \SI{48,0}{} = \SI{146,0}{g.mol^{-1}}$ (arrondi cohérent).
    La masse est conservée : la masse des réactifs ($\SI{164,2}{g.mol^{-1}}$) égale la masse des produits ($\SI{146,2}{} + \SI{18,0}{} = \SI{164,2}{g.mol^{-1}}$).
\end{enumerate}
\end{exoresolu}

\courssub{Structure et nomenclature des dipeptides}

\paragraph{Anatomie d'un dipeptide.}
Un dipeptide possède une structure linéaire caractéristique. Il est essentiel de savoir l'orienter et d'identifier ses parties fonctionnelles :
\begin{itemize}
    \item L'\textbf{extrémité N-terminale} (ou amino-terminale) : porte le groupe amine libre $\ce{-NH2}$ (ou $\ce{-NH3+}$ en solution). C'est le « début » de la chaîne.
    \item L'\textbf{extrémité C-terminale} (ou carboxy-terminale) : porte le groupe carboxyle libre $\ce{-COOH}$ (ou $\ce{-COO-}$ en solution). C'est la « fin » de la chaîne.
    \item La \textbf{liaison peptidique} \ce{-CO-NH-} : elle est rigide, plane et à caractère partiel de double liaison (résonance), ce qui limite la rotation autour de l'axe C–N et structure l'espace tridimensionnel de la protéine.
\end{itemize}

\begin{popfigure}
\centering
\begin{tikzpicture}[node distance=0.8cm, >=Stealth, font=\small]
    % Backbone nodes
    \node (Nterm) {$\ce{H2N}$};
    \node[right=0.3cm of Nterm] (CA1) {$\ce{CH(R1)}$};
    \node[right=1.0cm of CA1] (C1) {$\ce{C(=O)}$};
    \node[right=0.7cm of C1] (Npep) {$\ce{NH}$};
    \node[right=1.0cm of Npep] (CA2) {$\ce{CH(R2)}$};
    \node[right=1.0cm of CA2] (Cterm) {$\ce{COOH}$};

    % Bonds
    \draw[thick, popdark] (Nterm) -- (CA1) -- (C1) -- (Npep) -- (CA2) -- (Cterm);
    \draw[thick, popdark] (C1) -- ++(0,0.5) node[above] {$\ce{O}$};
    \draw[thick, popdark] (CA1) -- ++(0,-0.5) node[below] {$\ce{R1}$};
    \draw[thick, popdark] (CA2) -- ++(0,-0.5) node[below] {$\ce{R2}$};
    \draw[thick, popdark] (Npep) -- ++(0,0.5) node[above] {$\ce{H}$};

    % Colored background boxes for regions
    \fill[popblueL, opacity=0.5, rounded corners=3pt] ($(Nterm.south west)+(-0.1,-0.1)$) rectangle ($(C1.north east)+(0.1,0.1)$);
    \fill[poporangeL, opacity=0.5, rounded corners=3pt] ($(C1.south west)+(-0.1,-0.1)$) rectangle ($(Npep.north east)+(0.1,0.1)$);
    \fill[popgreenL, opacity=0.5, rounded corners=3pt] ($(Npep.south west)+(-0.1,-0.1)$) rectangle ($(Cterm.north east)+(0.1,0.1)$);
    
    % Redraw bonds on top of fill
    \draw[thick, popdark] (Nterm) -- (CA1) -- (C1) -- (Npep) -- (CA2) -- (Cterm);
    \draw[thick, popdark] (C1) -- ++(0,0.5) node[above] {$\ce{O}$};
    \draw[thick, popdark] (CA1) -- ++(0,-0.5) node[below] {$\ce{R1}$};
    \draw[thick, popdark] (CA2) -- ++(0,-0.5) node[below] {$\ce{R2}$};
    \draw[thick, popdark] (Npep) -- ++(0,0.5) node[above] {$\ce{H}$};

    % Labels
    \node[above=0.7cm of CA1, align=center, text width=2.5cm, popblue] (labN) {\textbf{Extrémité N-terminale}\\\ce{NH2} libre};
    \node[above=0.7cm of Npep, align=center, text width=2.5cm, poporange] (labPep) {\textbf{Liaison peptidique}\\\ce{-CO-NH-}};
    \node[above=0.7cm of CA2, align=center, text width=2.5cm, popgreen] (labC) {\textbf{Extrémité C-terminale}\\\ce{COOH} libre};

    \draw[popblue, thick, ->] (labN.south) -- ++(0,-0.3);
    \draw[poporange, thick, ->] (labPep.south) -- ++(0,-0.3);
    \draw[popgreen, thick, ->] (labC.south) -- ++(0,-0.3);
\end{tikzpicture}
\legende{Structure schématique d'un dipeptide. Les zones colorées délimitent l'extrémité N-terminale (bleu), la liaison peptidique (orange) et l'extrémité C-terminale (vert). R1 et R2 sont les chaînes latérales des deux résidus.}
\end{popfigure}

\paragraph{Nomenclature des dipeptides.}
La nomenclature suit le sens chimique de la chaîne : on cite les résidus dans l'ordre \textbf{N-terminal $\rightarrow$ C-terminal}. On utilise les abréviations à trois lettres (ou à une lettre) des acides $\alpha$-aminés constitutifs, séparées par un trait d'union.

\begin{methode}[Nommer un peptide]
\begin{enumerate}
    \item Identifier l'extrémité N-terminale (groupement $\ce{-NH2}$ libre ou $\ce{-NH3+}$).
    \item Partir de cette extrémité et avancer le long de la chaîne en suivant les liaisons peptidiques \ce{-CO-NH-}.
    \item Noter l'abréviation de chaque résidu d'acide $\alpha$-aminé rencontré dans l'ordre.
    \item Relier les abréviations par des traits d'union. Exemple : \textbf{Gly–Ala} (glycyl-alanine).
\end{enumerate}
\end{methode}

\paragraph{Isomérie de séquence.}
Avec deux acides $\alpha$-aminés \textbf{différents} A et B, on peut former \textbf{deux dipeptides distincts} : A–B et B–A. Ce sont des isomères de constitution (isomères de séquence) : ils ont la même formule brute mais des propriétés physico-chimiques et biologiques différentes.
\begin{itemize}
    \item \textbf{Ala–Gly} : extrémité N-terminale = Alanine, C-terminale = Glycine.
    \item \textbf{Gly–Ala} : extrémité N-terminale = Glycine, C-terminale = Alanine.
\end{itemize}
Si le mélange réactionnel contient A et B sans protection sélective, on obtient un mélange de \textbf{quatre} dipeptides possibles : A–A, A–B, B–A, B–B. C'est pourquoi la synthèse dirigée en laboratoire (et la biosynthèse au ribosome) nécessite des stratégies de protection/déprotection des fonctions pour ne former qu'un seul produit.

\begin{attention}[Pièges classiques à éviter]
\begin{itemize}
    \item \textbf{Oublier l'eau :} L'équation-bilan \textbf{doit} faire apparaître $\ce{+ H2O}$ dans les produits. C'est une condensation.
    \item \textbf{Sens de la liaison :} La liaison peptidique s'écrit \ce{-CO-NH-} (carbonyl $\rightarrow$ azote). L'écrire \ce{-NH-CO-} inverse le sens de la chaîne et change l'identité du dipeptide (confusion N/C terminale).
    \item \textbf{Nommage inversé :} Écrire « Gly–Ala » pour une structure où l'alanine est en N-terminal. Le premier nom cité est \emph{toujours} le résidu N-terminal.
    \item \textbf{Chiralté :} La condensation ne modifie pas la configuration du carbone $\alpha$ (pas de racémisation spontanée dans les conditions douces), mais en synthèse chimique, l'activation du carboxyle peut provoquer une racémisation partielle $\rightarrow$ utiliser des groupes protecteurs adaptés.
\end{itemize}
\end{attention}

\courssub{Des dipeptides aux protéines : la chaîne polypeptidique}

Le principe de la condensation se répète indéfiniment. L'extrémité C-terminale d'un dipeptide (groupe $\ce{-COOH}$ libre) peut réagir sur l'extrémité N-terminale d'un troisième acide $\alpha$-aminé (groupe $\ce{-NH2}$ libre) pour former un \textbf{tripeptide}, puis un \textbf{tétrapeptide}, etc. On parle de \textbf{polypeptide} pour une chaîne de plusieurs dizaines à centaines de résidus, et de \textbf{protéine} pour un polypeptide plié, fonctionnel, souvent de masse molaire supérieure à \SI{10000}{g.mol^{-1}} (environ 100 résidus).

\begin{popfigure}
\centering
\chemfig{
    H_2N-[:30]CH(-[:90]CH_3)-[:30]C(=[:90]O)-[:-30]N(-[:-90]H)-[:30]CH_2-[:30]C(=[:90]O)-[:-30]N(-[:-90]H)-[:30]CH(-[:90]CH_2SH)-[:30]C(=[:90]O)-[:-30]OH
}
\legende{Structure du tripeptide Ala–Gly–Cys. Les deux liaisons peptidiques \ce{-CO-NH-} relient les trois résidus. On note la conservation des extrémités N-terminale (Ala) et C-terminale (Cys).}
\end{popfigure}

\begin{savaistu}[L'aspartame, un dipeptide édulcorant dans nos boissons]
L'\textbf{aspartame} (E951), édulcorant intense présent dans de nombreuses boissons « light » vendues dans les supermarchés de Douala, Yaoundé ou Bafoussam, est un dérivé du dipeptide \textbf{Asp–Phe} (aspartame = L-aspartyl-L-phénylalanine méthyl ester). Sa puissance sucrante est environ 200 fois celle du saccharose. Cependant, il est déconseillé aux personnes atteintes de \textbf{phénylcétonurie} (maladie génétique de dégradation de la phénylalanine) car l'hydrolyse digestive libère de la phénylalanine. C'est un exemple concret où la stéréochimie (configuration L) et la nature de la liaison peptidique conditionnent une propriété biologique majeure.
\end{savaistu}

\begin{exoresolu}[Synthèse dirigée et identification]
\textbf{Énoncé.} On souhaite synthétiser le dipeptide Val–Ser (Valine–Sérine).
Données : Valine (Val, $\ce{C5H11NO2}$, $M=\SI{117,1}{g.mol^{-1}}$) ; Sérine (Ser, $\ce{C3H7NO3}$, $M=\SI{105,1}{g.mol^{-1}}$).
\begin{enumerate}
    \item Écrire l'équation-bilan de la réaction de condensation.
    \item Donner la formule brute et la masse molaire du dipeptide Val–Ser.
    \item Représenter la structure semi-développée du dipeptide en entourant la liaison peptidique.
    \item Quelle serait la structure du dipeptide isomère Ser–Val ?
\end{enumerate}

\textbf{Solution détaillée.}
\begin{enumerate}
    \item \textbf{Équation-bilan :}
    \[
    \ce{H2N-CH[CH(CH3)2]-COOH + H2N-CH(CH2OH)-COOH -> H2N-CH[CH(CH3)2]-CO-NH-CH(CH2OH)-COOH + H2O}
    \]
    (Valine à gauche $\rightarrow$ N-terminale ; Sérine à droite $\rightarrow$ C-terminale).
    \item \textbf{Formule brute :} $\ce{C5H11NO2 + C3H7NO3 - H2O = C8H16N2O4}$.
    \textbf{Masse molaire :} $M = \SI{117,1}{} + \SI{105,1}{} - \SI{18,0}{} = \SI{204,2}{g.mol^{-1}}$.
    \item \textbf{Représentation :}
    \begin{center}
    \chemfig{
        H_2N-[:30]CH(-[:90]CH(-[:30]CH_3)(-[:-30]CH_3))-[:30]C(=[:90]O)-[:-30]N(-[:-90]H)-[:30]CH(-[:90]CH_2OH)-[:30]C(=[:90]O)-[:-30]OH
    }
    \end{center}
    La liaison peptidique est le groupement \ce{-CO-NH-} central (encadré mentalement entre le carbone carbonylé de la Val et l'azote de la Ser).
    \item \textbf{Ser–Val :} L'extrémité N-terminale est la Sérine, l'extrémité C-terminale est la Valine.
    Structure : $\ce{H2N-CH(CH2OH)-CO-NH-CH[CH(CH3)2]-COOH}$.
    C'est un composé différent de Val–Ser (isomère de séquence).
\end{enumerate}
\end{exoresolu}

\begin{exercice}[Entraînement : Dipeptides d'acides aminés aromatiques]
On considère la phénylalanine (Phe, $\ce{C9H11NO2}$) et la tyrosine (Tyr, $\ce{C9H11NO3}$).
\begin{enumerate}
    \item Combien de dipeptides différents peut-on obtenir théoriquement à partir d'un mélange équimolaire de Phe et Tyr (sans protection) ? Les nommer.
    \item Écrire l'équation-bilan de formation du dipeptide Phe–Tyr.
    \item Calculer la masse molaire de Phe–Tyr ($M_{\text{Phe}}=\SI{165,2}{g.mol^{-1}}$, $M_{\text{Tyr}}=\SI{181,2}{g.mol^{-1}}$).
    \item Dessiner la structure topologique (schéma type) de Tyr–Phe en légendant les extrémités et la liaison peptidique.
\end{enumerate}
\end{exercice}

\begin{corrige}[Exercice : Dipeptides d'acides aminés aromatiques]
\begin{enumerate}
    \item Quatre dipeptides possibles : \textbf{Phe–Phe}, \textbf{Phe–Tyr}, \textbf{Tyr–Phe}, \textbf{Tyr–Tyr}.
    \item $\ce{H2N-CH(CH2Ph)-COOH + H2N-CH(CH2-Ph-OH)-COOH -> H2N-CH(CH2Ph)-CO-NH-CH(CH2-Ph-OH)-COOH + H2O}$.
    \item $M(\text{Phe–Tyr}) = \SI{165,2}{} + \SI{181,2}{} - \SI{18,0}{} = \SI{328,4}{g.mol^{-1}}$.
    \item Schéma : Extrémité N-terminale = Tyr (chaîne latérale $\ce{-CH2-Ph-OH}$) ; Liaison peptidique \ce{-CO-NH-} ; Extrémité C-terminale = Phe (chaîne latérale $\ce{-CH2-Ph}$).
    \begin{center}
    \begin{tikzpicture}[font=\small]
        \node (N) {$\ce{H2N}$};
        \node[right=0.3cm of N] (CA1) {$\ce{CH(CH2-Ph-OH)}$};
        \node[right=1.0cm of CA1] (C1) {$\ce{C(=O)}$};
        \node[right=0.7cm of C1] (Np) {$\ce{NH}$};
        \node[right=1.0cm of Np] (CA2) {$\ce{CH(CH2-Ph)}$};
        \node[right=1.0cm of CA2] (C) {$\ce{COOH}$};
        \draw[thick] (N)--(CA1)--(C1)--(Np)--(CA2)--(C);
        \draw[thick] (C1)--++(0,0.4) node[above]{O};
        \draw[thick] (Np)--++(0,0.4) node[above]{H};
        \node[above=0.6cm of CA1, popblue] {N-term (Tyr)};
        \node[above=0.6cm of Np, poporange] {Liaison peptidique};
        \node[above=0.6cm of CA2, popgreen] {C-term (Phe)};
    \end{tikzpicture}
    \end{center}
\end{enumerate}
\end{corrige}

\coursec{Synthèse dirigée d'un dipeptide (complément Bac)}

Dans la section précédente, tu as découvert la \cle{liaison peptidique} : en chauffant un mélange de deux acides $\alpha$-aminés, on peut former un dipeptide par élimination d'une molécule d'eau. Mais une question se pose immédiatement, et elle est capitale : si l'on mélange simplement deux acides $\alpha$-aminés différents, obtient-on \emph{le} dipeptide souhaité, ou un mélange inextricable de produits ? C'est tout l'enjeu de cette dernière section : comprendre pourquoi la synthèse d'un peptide exige une stratégie, et découvrir le principe de la \cle{synthèse dirigée}, celle-là même qui permet à l'industrie pharmaceutique de fabriquer des médicaments comme l'insuline.

\begin{attention}[Complément utile au Bac]
Ce paragraphe est un \textbf{complément} : il prolonge le programme officiel de Terminale TI sans le dépasser. Il t'aide à comprendre en profondeur la notion de liaison peptidique et il est fréquemment mobilisé dans les sujets de synthèse au Baccalauréat. Retiens l'essentiel : le \emph{problème} (mélange de produits) et la \emph{solution} (protéger, activer, condenser, déprotéger).
\end{attention}

\courssub{Le problème : une condensation non dirigée donne un mélange}

\courssub{Intuition : deux ouvriers, quatre façons de s'associer}

Imagine deux maçons, Ali et Bello, qui doivent poser deux briques l'une sur l'autre. Chacun a une « tête » (le groupe \ce{-NH2}) et un « pied » (le groupe \ce{-COOH}). Si tu les laisses travailler sans consigne, la tête de l'un peut se poser sur le pied de l'autre, mais aussi la tête d'Ali sur le pied d'un autre Ali, ou la tête de Bello sur le pied d'un autre Bello. Résultat : quatre constructions différentes !

Transposons à la chimie. Prenons deux acides $\alpha$-aminés : l'\textbf{alanine} (Ala, $\mathrm{R} = \ce{CH3}$) et la \textbf{glycine} (Gly, $\mathrm{R} = \ce{H}$). Chaque molécule possède :
\begin{itemize}
\item un groupe \textbf{amine} \ce{-NH2} qui joue le rôle de \emph{donneur} (il attaque) ;
\item un groupe \textbf{carboxyle} \ce{-COOH} qui joue le rôle d'\emph{accepteur} (il est attaqué).
\end{itemize}

En chauffant le mélange Ala + Gly, \emph{toutes} les rencontres sont possibles :
\begin{itemize}
\item le \ce{-NH2} de Ala attaque le \ce{-COOH} de Ala $\rightarrow$ dipeptide \textbf{Ala--Ala} ;
\item le \ce{-NH2} de Gly attaque le \ce{-COOH} de Gly $\rightarrow$ dipeptide \textbf{Gly--Gly} ;
\item le \ce{-NH2} de Gly attaque le \ce{-COOH} de Ala $\rightarrow$ dipeptide \textbf{Ala--Gly} ;
\item le \ce{-NH2} de Ala attaque le \ce{-COOH} de Gly $\rightarrow$ dipeptide \textbf{Gly--Ala}.
\end{itemize}

\begin{popfigure}
\begin{center}
\begin{tikzpicture}[
  box/.style={draw=popblue, thick, rounded corners=3pt, fill=popblueL, minimum width=2.6cm, minimum height=0.9cm, font=\small\bfseries, text=popdark},
  lab/.style={font=\footnotesize\itshape, text=popmuted}
]
\node[box] (aa) at (0,0) {Ala--Ala};
\node[box] (ab) at (3.6,0) {Ala--Gly};
\node[box] (ba) at (7.2,0) {Gly--Ala};
\node[box] (bb) at (10.8,0) {Gly--Gly};
\node[lab] at (0,-0.85) {homodipeptide};
\node[lab] at (3.6,-0.85) {hétérodipeptide};
\node[lab] at (7.2,-0.85) {hétérodipeptide};
\node[lab] at (10.8,-0.85) {homodipeptide};
\draw[poporange, very thick, rounded corners=4pt] (1.9,-0.55) rectangle (8.9,0.65);
\node[font=\footnotesize\bfseries, text=poporange] at (5.4,1.05) {deux dipeptides distincts !};
\end{tikzpicture}
\legende{Les quatre dipeptides obtenus par simple chauffage d'un mélange alanine + glycine. Attention : Ala--Gly et Gly--Ala sont deux molécules \emph{différentes} car l'ordre des résidus est inversé.}
\end{center}
\end{popfigure}

\begin{attention}[Piège classique : Ala--Gly n'est pas Gly--Ala]
Beaucoup d'élèves comptent seulement trois dipeptides en oubliant que l'\textbf{ordre} des résidus compte. Dans Ala--Gly, l'extrémité \ce{-NH2} libre (N-terminal) est portée par l'alanine ; dans Gly--Ala, elle est portée par la glycine. Ce sont deux composés différents, de propriétés différentes. Un mélange Ala + Gly donne donc bien \textbf{quatre} dipeptides.
\end{attention}

\begin{methode}[Prévoir les produits d'une condensation non dirigée]
Pour dénombrer les dipeptides formés en mélangeant $n$ acides $\alpha$-aminés \emph{différents} :
\begin{enumerate}
\item Chaque dipeptide est un couple ordonné (résidu N-terminal, résidu C-terminal).
\item Il y a $n$ choix pour la première position et $n$ choix pour la seconde.
\item Nombre total de dipeptides possibles : $n \times n = n^2$.
\end{enumerate}
Exemples : $n = 2$ acides $\rightarrow$ $2^2 = 4$ dipeptides ; $n = 3$ acides $\rightarrow$ $3^2 = 9$ dipeptides. Et pour un tripeptide à partir de $n$ acides différents : $n^3$ possibilités !
\end{methode}

\begin{exoresolu}[Dénombrement des dipeptides]
On chauffe un mélange équimolaire de trois acides $\alpha$-aminés : alanine (Ala), glycine (Gly) et valine (Val). Combien de dipeptides différents peut-on obtenir ? Les citer.
\tcblower
\textbf{Solution.} Avec $n = 3$ acides $\alpha$-aminés différents, le nombre de dipeptides possibles est $n^2 = 9$. En effet, chaque position (N-terminale puis C-terminale) offre 3 choix :
\begin{center}
\begin{tabularx}{\linewidth}{|l|X|X|X|}
\hline
\rowcolor{popblueL}
\textbf{1\ier{} résidu} & \multicolumn{3}{c|}{\textbf{2\ieme{} résidu}} \\
\hline
 & \textbf{Ala} & \textbf{Gly} & \textbf{Val} \\
\hline
\textbf{Ala} & Ala--Ala & Ala--Gly & Ala--Val \\
\hline
\textbf{Gly} & Gly--Ala & Gly--Gly & Gly--Val \\
\hline
\textbf{Val} & Val--Ala & Val--Gly & Val--Val \\
\hline
\end{tabularx}
\end{center}
On obtient 3 homodipeptides (Ala--Ala, Gly--Gly, Val--Val) et 6 hétérodipeptides. Séparer ces 9 produits serait un cauchemar de laboratoire : d'où la nécessité d'une synthèse \emph{dirigée}.
\end{exoresolu}

\courssub{La solution : la synthèse dirigée}

\begin{definition}[Synthèse dirigée]
Une \cle{synthèse dirigée} (ou synthèse sélective) est une synthèse dans laquelle on \textbf{contrôle l'ordre d'enchaînement} des acides $\alpha$-aminés grâce à la \textbf{protection sélective} de certaines fonctions (\ce{-NH2} ou \ce{-COOH}) et à l'\textbf{activation} de la fonction qui doit réagir. On obtient ainsi \emph{un seul} dipeptide, dans l'ordre voulu.
\end{definition}

L'idée est simple et élégante : si l'on veut fabriquer \textbf{uniquement} Ala--Gly, il faut :
\begin{itemize}
\item \emph{interdire} au groupe \ce{-NH2} de l'alanine de réagir (sinon on obtient Gly--Ala et Ala--Ala) : on le \textbf{protège} ;
\item \emph{interdire} au groupe \ce{-COOH} de la glycine de réagir (sinon on obtient Gly--Gly et Gly--Ala) : on le \textbf{protège} ;
\item \emph{encourager} le groupe \ce{-COOH} de l'alanine à réagir : on l'\textbf{active}.
\end{itemize}
Il ne reste alors qu'une seule rencontre possible : le \ce{-COOH} activé de Ala avec le \ce{-NH2} libre de Gly.

\courssub{Les quatre étapes de la synthèse dirigée}

\textbf{Étape 1 : protection du groupe \ce{-NH2} de l'acide A.}\par
On transforme l'amine de l'alanine en une fonction non nucléophile (par exemple un carbamate, noté classiquement Z-- ou Boc--). Le groupe protégé ne peut plus attaquer un carbonyle :
\[ \ce{H2N-CHR-COOH ->[\text{agent protecteur}] P-NH-CHR-COOH} \]
où P désigne le groupe protecteur. Cette réaction est rapide, totale et \emph{réversible} : on pourra retirer P plus tard.

\textbf{Étape 2 : activation du groupe \ce{-COOH} de l'acide A.}\par
Le groupe carboxyle est naturellement peu réactif vis-à-vis des amines. On le transforme en dérivé plus réactif (chlorure d'acyle, ester activé, anhydride mixte), noté \ce{-COX} :
\[ \ce{P-NH-CHR-COOH ->[\text{agent activant}] P-NH-CHR-COX} \]
Le carbone du carbonyle devient beaucoup plus électrophile, donc plus facilement attaqué.

\textbf{Étape 3 : condensation avec l'acide B protégé sur son $-\mathrm{COOH}$.}\par
La glycine, dont le carboxyle est bloqué (par exemple sous forme d'ester $-\mathrm{COOR}'$), apporte son unique fonction libre : le $-\mathrm{NH}_2$. Une seule réaction est possible :
\[ \mathrm{P{-}NH{-}CHR{-}COX} + \mathrm{H_2N{-}CHR'{-}COOR'} \rightarrow \mathrm{P{-}NH{-}CHR{-}CO{-}NH{-}CHR'{-}COOR'} \]
C'est la formation de la \cle{liaison peptidique} $\mathrm{-CO{-}NH-}$, exactement à l'endroit voulu.

\textbf{Étape 4 : déprotection.}\par
On retire enfin les deux groupes protecteurs (P sur l'amine, R' sur l'ester) par hydrolyse ou hydrogénolyse, dans des conditions douces qui \emph{ne cassent pas} la liaison peptidique :
\[ \mathrm{P{-}NH{-}CHR{-}CO{-}NH{-}CHR'{-}COOR'} \xrightarrow{\text{déprotection}} \mathrm{H_2N{-}CHR{-}CO{-}NH{-}CHR'{-}COOH} \]
On isole le dipeptide \textbf{Ala--Gly pur}, avec un excellent rendement.

\begin{popfigure}
\begin{center}
\begin{tikzpicture}[
  etape/.style={draw=popblue, very thick, rounded corners=4pt, fill=popblueL, minimum width=3.1cm, minimum height=2.1cm, align=center, font=\footnotesize, text=popdark},
  titre/.style={font=\footnotesize\bfseries, text=popblue},
  fleche/.style={-{Stealth[length=3mm]}, very thick, poporange}
]
\node[etape] (e1) at (0,0) {\textbf{1. Protection}\\[2pt] \text{-}\ce{NH2} de A bloqué\\ par le groupe P};
\node[etape] (e2) at (4.4,0) {\textbf{2. Activation}\\[2pt] \text{-}\ce{COOH} de A rendu\\ réactif (\text{-COX})};
\node[etape] (e3) at (8.8,0) {\textbf{3. Condensation}\\[2pt] + B (dont \text{-}\ce{COOH}\\ est protégé en \text{-COOR'})};
\node[etape, fill=popgreenL, draw=popgreen] (e4) at (13.2,0) {\textbf{4. Déprotection}\\[2pt] retrait de P et R'\\ $\Rightarrow$ dipeptide A--B};
\draw[fleche] (e1) -- (e2);
\draw[fleche] (e2) -- (e3);
\draw[fleche] (e3) -- (e4);
\node[titre] at (0,1.45) {A : \chemfig[atom sep=1.4em]{H_2N-CH(R)-COOH}};
\node[titre] at (13.2,1.45) {produit unique};
\end{tikzpicture}
\legende{Schéma-bloc des quatre étapes de la synthèse dirigée du dipeptide A--B : protection de l'amine de A, activation de son acide, condensation avec B protégé, puis déprotection.}
\end{center}
\end{popfigure}

\begin{aretenir}[L'idée clé]
La \textbf{protection} et la \textbf{déprotection} des fonctions \ce{-NH2} et \ce{-COOH} garantissent l'\textbf{ordre d'enchaînement des résidus}. Sans protection : $n^2$ dipeptides. Avec protection : \emph{un seul} dipeptide, celui que l'on a choisi. C'est le principe de toute la synthèse peptidique moderne, y compris la synthèse sur support solide (prix Nobel de chimie 1984, Bruce Merrifield).
\end{aretenir}

\begin{exemplebox}[Synthèse dirigée de Ala--Gly : bilan]
Objectif : préparer uniquement le dipeptide Ala--Gly.
\begin{enumerate}
\item On protège \ce{NH2} de l'alanine $\rightarrow$ P--Ala--OH.
\item On active son \ce{COOH} $\rightarrow$ P--Ala--OX.
\item On ajoute la glycine dont le \ce{COOH} est estérifié : \ce{H2N-CH2-COOR}'. Seule la liaison Ala--Gly se forme.
\item On déprotège $\rightarrow$ \ce{H2N-CH(CH3)-CO-NH-CH2-COOH}, c'est-à-dire Ala--Gly pur.
\end{enumerate}
Sans stratégie, le même mélange de départ aurait donné 4 dipeptides à séparer, avec un rendement en Ala--Gly d'environ \SI{25}{\percent} seulement.
\end{exemplebox}

\begin{attention}[Deux erreurs à ne jamais commettre]
\begin{itemize}
\item \textbf{Protéger la mauvaise fonction} : pour obtenir A--B, c'est le \ce{-NH2} de A et le \ce{-COOH} de B qu'il faut bloquer. Si tu bloques le \ce{-COOH} de A, tu obtiendras B--A !
\item \textbf{Oublier l'activation} : un simple chauffage du mélange, même avec des fonctions protégées, est lent et peu sélectif. L'activation du \ce{-COOH} rend la condensation rapide et quasi totale.
\end{itemize}
\end{attention}

\courssub{Application industrielle : des peptides médicaments}

La synthèse dirigée, répétée résidu après résidu, permet de construire des chaînes peptidiques longues et parfaitement ordonnées. C'est ainsi que l'industrie pharmaceutique produit des \textbf{peptides médicamenteux} : hormones, antibiotiques, analogues d'hormones. L'\textbf{ocytocine} (hormone de l'accouchement, 9 acides aminés) fut le premier peptide synthétisé en laboratoire (Vincent du Vigneaud, prix Nobel 1955). Aujourd'hui, ces synthèses sont automatisées sur support solide.

\begin{savaistu}[L'insuline, une protéine vitale]
L'\textbf{insuline}, hormone indispensable aux personnes diabétiques, est une protéine de \textbf{51 acides aminés} répartis sur deux chaînes reliées par des ponts disulfure. Pendant des décennies, on l'extrayait du pancréas de porcs ou de bœufs — y compris dans les abattoirs d'Afrique centrale. Depuis 1982, elle est produite par \textbf{génie génétique} : on insère le gène humain de l'insuline dans une bactérie (\emph{Escherichia coli}) ou une levure, qui fabrique alors l'hormone humaine en grandes quantités, pure et sans risque de contamination. C'est l'un des premiers médicaments issus des biotechnologies — une révolution qui a sauvé des millions de vies, y compris au Cameroun où le diabète touche une part croissante de la population.
\end{savaistu}

\begin{exercice}[À toi de jouer]
On souhaite synthétiser \emph{uniquement} le dipeptide Gly--Ala à partir de glycine et d'alanine.
\begin{enumerate}
\item Quelles fonctions faut-il protéger, et sur quel acide $\alpha$-aminé ?
\item Quelle fonction faut-il activer ?
\item Combien de dipeptides obtiendrait-on sans aucune protection ?
\end{enumerate}
\end{exercice}

\begin{corrige}[Exercice 1]
\begin{enumerate}
\item Le résidu N-terminal doit être la glycine : on protège donc le \ce{-NH2} de la \textbf{glycine}. Le résidu C-terminal est l'alanine : on protège le \ce{-COOH} de l'\textbf{alanine} (par exemple sous forme d'ester).
\item On active le \ce{-COOH} de la glycine (transformation en dérivé réactif \ce{-COX}) pour qu'il soit attaqué par le \ce{-NH2} libre de l'alanine.
\item Sans protection, le mélange Gly + Ala conduit à $n^2 = 2^2 = 4$ dipeptides : Gly--Gly, Ala--Ala, Gly--Ala et Ala--Gly.
\end{enumerate}
\end{corrige}

\begin{aretenir}[Bilan de la section]
\begin{itemize}
\item Chauffer un mélange de $n$ acides $\alpha$-aminés différents donne $n^2$ dipeptides : la condensation simple n'est \textbf{pas sélective}.
\item La \textbf{synthèse dirigée} repose sur quatre étapes : \textbf{protéger} le \ce{-NH2} de A, \textbf{activer} son \ce{-COOH}, \textbf{condenser} avec B (dont le \ce{-COOH} est protégé), puis \textbf{déprotéger}.
\item La protection/déprotection des fonctions garantit l'\textbf{ordre d'enchaînement des résidus} : on obtient un seul dipeptide, celui voulu.
\item Ce principe est à la base de la production industrielle de peptides médicamenteux (ocytocine, insuline).
\end{itemize}
\end{aretenir}

\newpage

\coursec{Travaux pratiques}

\courssub{Objectifs du TP}
\begin{enumerate}
    \item Mettre en évidence expérimentalement le \cle{caractère amphotère} d'un acide α-aminé (la glycine) en montrant qu'il réagit à la fois avec une solution acide et avec une solution basique.
    \item Réaliser l'\cle{hydrolyse acide} d'une protéine alimentaire (blanc d'œuf ou lait en poudre) pour la dépolymériser en acides α-aminés constitutifs.
    \item Identifier la présence des acides α-aminés libérés par l'hydrolyse grâce au \cle{test à la ninhydrine} (coloration violette caractéristique, dite de Ruhemann) et le comparer au \cle{test au réactif de Biuret} (spécifique de la liaison peptidique).
    \item Manipuler en respectant strictement les règles de sécurité au laboratoire (port des EPI, manipulation des acides/bases concentrés par l'enseignant).
\end{enumerate}

\courssub{Matériel et produits}
\begin{center}
\begin{tabularx}{\linewidth}{|l|X|l|}\hline
\textbf{Désignation} & \textbf{Description / Spécifications} & \textbf{Alternative locale (Cameroun)} \\ \hline
Verrerie & 3 tubes à essai (16×160 mm) + support, 1 ballon à fond rond (100 mL), 1 réfrigérant à boules (ou droit), 1 bec Bunsen ou réchaud à gaz, 1 trépied + treillis + triangle en terre, 1 erlenmeyer (100 mL), 1 burette ou pipette graduée (10 mL), 1 pH-mètre ou papier pH universel (1–14), baguettes de verre. & Tubes en verre récupérés (flacons de sirop stérilisés), réfrigérant artisanal (tube verre dans tube PVC + eau courante), réchaud à gaz « Gazelle » ou charbon de bois. \\ \hline
Réactifs & Glycine pure (\ce{NH2-CH2-COOH}) : \SI{2}{g} ; HCl \SI{0,1}{mol.L^{-1}} (\SI{20}{mL}) ; NaOH \SI{0,1}{mol.L^{-1}} (\SI{20}{mL}) ; Eau distillée ; Solution de ninhydrine \SI{0,2}{\%} (m/v) dans l'éthanol ou l'eau (\SI{10}{mL}) ; Réactif de Biuret (\ce{CuSO4.5H2O} + NaOH + tartrate de potassium) (\SI{10}{mL}) ; Blanc d'œuf frais (\SI{5}{mL}) ou lait en poudre (\SI{2}{g}) ; HCl \SI{6}{mol.L^{-1}} (\SI{20}{mL}) pour hydrolyse ; NaOH \SI{6}{mol.L^{-1}} pour neutralisation. & Glycine : commande pharmacie/vétérinaire (Douala/Yaoundé) ; HCl/NaOH concentrés : fournisseurs chimie (SOCACHIM, CIPHAR) ; Ninhydrine : poudre à dissoudre ; Biuret : \ce{CuSO4.5H2O} + NaOH + sel de Rochelle (pharmacie) ; Eau distillée : eau de batterie ou eau de pluie filtrée bouillie. \\ \hline
Équipements de protection & Blouse en coton, lunettes de protection, gants nitrile, manipulation sous hotte ou à l'air libre (zone ventilée). & Lunettes de soudure (teinte claire), gants ménagers épais, blouse d'infirmier. \\ \hline
\end{tabularx}
\end{center}

\courssub{Sécurité}
\begin{attention}[Règles de sécurité impératives — Pictogrammes et gestes]
\textbf{Pictogrammes concernés :}
\begin{itemize}
    \item \textbf{Corrosif (GHS05) :} HCl \SI{6}{mol.L^{-1}} et NaOH \SI{6}{mol.L^{-1}} provoquent de graves brûlures cutanées et oculaires. \emph{Geste :} Manipulation \textbf{uniquement par l'enseignant} pour la préparation des solutions mères. Port obligatoire de gants, lunettes, blouse. En cas de projection : rinçage immédiat \SI{15}{min} à l'eau courante (œil : station de lavage oculaire ou bouteille d'eau), alerte infirmerie.
    \item \textbf{Irritant (GHS07) :} Ninhydrine (coloration cutanée tenace, irritation voies respiratoires), Biuret (\ce{Cu^{2+}} toxique par ingestion). \emph{Geste :} Ne pas respirer les vapeurs (ninhydrine chauffée), gants obligatoires.
    \item \textbf{Inflammable (GHS02) :} Éthanol (solvant ninhydrine) près flamme nue (réchaud/Bunsen). \emph{Geste :} Éloigner le flacon d'éthanol du poste de chauffe ; utiliser un bain-marie électrique si possible pour le test ninhydrine.
\end{itemize}
\textbf{Règles générales :} Pas de manger/boire au labo. Étiqueter tous les béchers/tubes. Déchets acides/bases : neutralisation approximative avant évacuation (fosse septique). Déchets ninhydrine/Biuret : bac récupération « métaux lourds/organiques » (enseignant).
\end{attention}

\courssub{Schéma du dispositif}
\begin{popfigure}
\begin{tikzpicture}[scale=0.9, every node/.style={transform shape},
    vessel/.style={draw=popdark, thick, fill=popblueL!30, rounded corners=2pt},
    heater/.style={draw=popdark, thick, fill=poporangeL!50, rounded corners=3pt},
    clamp/.style={draw=popdark, very thick},
    water/.style={draw=popblue, thick, fill=popblue!20},
    vapor/.style={draw=popmuted, dashed, ->, thick},
    label/.style={font=\small, text=popdark, align=center}]
% Réfrigérant (tube extérieur)
\draw[vessel] (0,5.5) rectangle (1.5,8.5);
\draw[water] (0,8.5) rectangle (1.5,9.0); node[label, above] at (0.75,9.2) {Entrée eau (haut)};
\draw[water] (0,5.0) rectangle (1.5,5.5); node[label, below] at (0.75,4.7) {Sortie eau (bas)};
% Tube intérieur (vapeurs)
\draw[clamp] (0.6,5.5) -- (0.6,1.5);
\draw[clamp] (0.9,5.5) -- (0.9,1.5);
% Ballon
\draw[vessel] (-1.5,1.5) .. controls (-1.5,0.5) and (2.5,0.5) .. (2.5,1.5) -- (2.5,2.5) .. controls (2.5,3.5) and (-1.5,3.5) .. (-1.5,2.5) -- cycle;
\node[label] at (0.5,2.0) {Blanc d'œuf + HCl \SI{6}{mol.L^{-1}}\\+ billes de porcelaine};
% Chauffage
\draw[heater] (-2,0) rectangle (3,1.0);
\node[label] at (0.5,0.5) {Réchaud à gaz / Bunsen};
% Flamme
\draw[poporange, thick, decorate, decoration={snake, amplitude=1pt, segment length=4pt}] (0,0) -- (0,-0.5) node[below, label] {Flamme douce};
% Pinces
\draw[clamp] (-2.5,2.5) -- (-2.5,4) node[left, label] {Pince à boule};
\draw[clamp] (3.5,8.5) -- (3.5,10) node[right, label] {Pince à boule};
% Flèches vapeur
\draw[vapor] (0.75,1.5) -- (0.75,5.5);
\node[label, right] at (1.0,3.5) {Vapeurs d'HCl / H₂O};
% Légendes
\node[label, anchor=west] at (4,7) {\textbf{Dispositif de chauffage à reflux}};
\node[label, anchor=west] at (4,6.3) {1. Ballon à fond rond (100 mL)};
\node[label, anchor=west] at (4,5.6) {2. Réfrigérant à boules (condenseur)};
\node[label, anchor=west] at (4,4.9) {3. Circulation eau froide (bas $\to$ haut)};
\node[label, anchor=west] at (4,4.2) {4. Source de chaleur (réchaud gaz)};
\node[label, anchor=west] at (4,3.5) {5. Billes de porcelaine (ébullition calme)};
\node[label, anchor=west] at (4,2.8) {6. Support + pinces};
\end{tikzpicture}
\legende{Dispositif de chauffage à reflux pour l'hydrolyse acide du blanc d'œuf (Séance 2). Le réfrigérant empêche la perte d'HCl volatil et assure un chauffage prolongé à \SI{100}{\celsius}.}
\end{popfigure}

\courssub{Protocole}
\textbf{Séance 1 : Caractère amphotère de la glycine (30 min)}
\begin{enumerate}
    \item \textbf{Préparation de la solution de glycine :} Peser \SI{0,75}{g} de glycine solide (soit \SI{0,01}{mol}, $M_{\text{Gly}} = \SI{75,07}{g.mol^{-1}}$). La dissoudre dans \SI{25,0}{mL} d'eau distillée dans un erlenmeyer (agiter jusqu'à dissolution complète). On obtient une solution \SI{0,40}{mol.L^{-1}}.
    \item \textbf{Répartition :} Verser \SI{5,0}{mL} de cette solution dans chacun des 3 tubes à essai étiquetés T1, T2, T3.
    \item \textbf{Test acide (T1) :} Ajouter \SI{2,0}{mL} de solution d'HCl \SI{0,1}{mol.L^{-1}}. Agiter. Mesurer le pH (papier pH ou pH-mètre). Noter la valeur.
    \item \textbf{Test basique (T2) :} Ajouter \SI{2,0}{mL} de solution de NaOH \SI{0,1}{mol.L^{-1}}. Agiter. Mesurer le pH. Noter la valeur.
    \item \textbf{Témoin (T3) :} Ajouter \SI{2,0}{mL} d'eau distillée. Agiter. Mesurer le pH. Noter la valeur.
    \item \textbf{Test colorimétrique (optionnel) :} Ajouter 2 gouttes de \cle{Bromothymol Bleu (BBT)} dans chaque tube. Noter les couleurs : jaune (acide), bleu (basique), vert (neutre/tampon).
\end{enumerate}

\textbf{Séance 2 : Hydrolyse acide d'une protéine et détection (1 h 30 min)}
\begin{enumerate}
    \item \textbf{Mise en place du reflux :} Monter le dispositif (schéma ci-dessus). Vérifier l'étanchéité des joints (graisse silicone ou papier parafilm). Ouvrir l'eau du réfrigérant (filet continu, entrée par le bas).
    \item \textbf{Préparation du mélange :} Dans le ballon, introduire \SI{5,0}{mL} de blanc d'œuf frais (ou \SI{2,0}{g} de lait en poudre redissous dans \SI{10}{mL} d'eau) + \SI{20,0}{mL} de HCl \SI{6}{mol.L^{-1}} (manipulation enseignant) + 3-4 billes de porcelaine.
    \item \textbf{Chauffage :} Chauffer modérément (flamme douce ou thermostat 3-4) pour maintenir un bouillonnement calme (\SI{100}{\celsius}) pendant \textbf{20 minutes}. Surveiller le niveau de liquide (ne pas laisser à sec).
    \item \textbf{Refroidissement et neutralisation :} Arrêter le chauffage. Laisser refroidir \SI{10}{min}. Transvaser dans un erlenmeyer. Neutraliser \emph{lentement} et sous agitation par ajout de NaOH \SI{6}{mol.L^{-1}} (manipulation enseignant) jusqu'à pH $\approx$ 7 (papier pH). \textbf{Attention :} réaction exothermique violente, dégagement de chaleur, risque de projection. Ajouter le NaOH goutte à goutte.
    \item \textbf{Test de Biuret (sur hydrolysat neutre) :} Prélever \SI{2}{mL} d'hydrolysat neutre dans un tube (Tube B1). Ajouter \SI{2}{mL} de réactif de Biuret. Agiter. Observer la couleur (violet = liaison peptidique présente). Pour comparaison, faire le test sur \SI{2}{mL} de blanc d'œuf non hydrolysé (Tube B0).
    \item \textbf{Test à la ninhydrine (sur hydrolysat neutre) :} Prélever \SI{1}{mL} d'hydrolysat neutre dans un tube (Tube N1). Ajouter \SI{1}{mL} de solution de ninhydrine \SI{0,2}{\%}. Boucher le tube. Plonger dans un bain-marie bouillant (\SI{100}{\celsius}) pendant \textbf{5 à 10 min}.
    \item \textbf{Témoins ninhydrine :} Tube N0 : \SI{1}{mL} eau + \SI{1}{mL} ninhydrine. Tube N2 : \SI{1}{mL} solution de glycine \SI{0,1}{M} (résidu Séance 1) + \SI{1}{mL} ninhydrine. Traiter identiquement (bain-marie 10 min).
    \item \textbf{Observation finale :} Sortir les tubes, laisser tiédir. Comparer les colorations.
\end{enumerate}

\courssub{Observations et résultats attendus}
\begin{center}
\begin{tabularx}{\linewidth}{|l|X|X|X|}\hline
\textbf{Épreuve} & \textbf{Tube / Condition} & \textbf{Observation (Couleur / pH)} & \textbf{Interprétation immédiate} \\ \hline
\textbf{Séance 1} & T1 : Glycine + HCl \SI{0,1}{M} & pH $\approx$ 2,5 -- 3,5 (BBT : \textbf{Jaune}) & La glycine accepte H⁺ : \textbf{caractère basique} (fonction carboxylate). \\ \hline
& T2 : Glycine + NaOH \SI{0,1}{M} & pH $\approx$ 10,5 -- 11,5 (BBT : \textbf{Bleu}) & La glycine donne H⁺ : \textbf{caractère acide} (fonction ammonium). \\ \hline
& T3 : Glycine + Eau (Témoin) & pH $\approx$ 5,5 -- 6,5 (BBT : \textbf{Vert}) & Forme \textbf{amphion (zwitterion)} prédominante, pH = pI (point isoélectrique). \\ \hline
\textbf{Séance 2} & B0 : Blanc d'œuf natif + Biuret & \textbf{Violet intense} & Protéines natives : nombreuses liaisons peptidiques accessibles. \\ \hline
& B1 : Hydrolysat 20 min + Biuret & \textbf{Violet pâle à incolore} & Hydrolyse partielle : liaisons peptidiques coupées $\to$ peptides courts / AA libres. \\ \hline
& N0 : Eau + Ninhydrine (Témoin) & \textbf{Incolore / Jaune pâle} & Réactif seul ne colore pas (ou jaune par dégradation thermique). \\ \hline
& N2 : Glycine + Ninhydrine & \textbf{Violet intense (pourpre de Ruhemann)} & Test positif caractéristique des acides α-aminés primaires. \\ \hline
& N1 : Hydrolysat + Ninhydrine & \textbf{Violet clair à moyen} & \textbf{Preuve de la libération d'acides α-aminés} par hydrolyse acide. \\ \hline
\end{tabularx}
\end{center}
\begin{aretenir}[Valeurs repères pour le compte-rendu]
\begin{itemize}
    \item pH solution glycine \SI{0,4}{M} (eau) $\approx$ \SI{6,0}{pH} (pI glycine = 5,97).
    \item pH après ajout HCl \SI{0,1}{M} (v/v 5/2) $\approx$ \SI{2,8}{pH}.
    \item pH après ajout NaOH \SI{0,1}{M} (v/v 5/2) $\approx$ \SI{11,2}{pH}.
    \item Temps d'hydrolyse standard Bac : \SI{20}{min} à reflux (HCl \SI{6}{mol.L^{-1}}). Hydrolyse totale $\approx$ \SI{24}{h} (non réalisable en TP).
    \item Ninhydrine : coloration maximale vers \SI{570}{nm} (violet). Proline/Hydroxyproline donnent une couleur \textbf{jaune-orangé} (amine secondaire).
\end{itemize}
\end{aretenir}

\courssub{Exploitation}
\begin{enumerate}
    \item \textbf{Équilibres acido-basiques de la glycine.} Écrire les deux demi-équations-bilans montrant le caractère acide puis basique de l'amphion de la glycine ($\ce{^{+}H3N-CH2-COO-}$). Préciser les couples acide/base mis en jeu.
    \item \textbf{Équilibre forme moléculaire / amphion.} Écrire l'équation-bilan de l'équilibre entre la forme moléculaire neutre ($\ce{H2N-CH2-COOH}$) et la forme ionique (amphion). Justifier pourquoi la forme moléculaire est quasi inexistante en solution aqueuse.
    \item \textbf{Calcul de pH (approché).} On rappelle que pour un acide aminé monoprotique, $\text{pH} = \frac{\text{p}K_{a1} + \text{p}K_{a2}}{2}$ au point isoélectrique. Données : $\text{p}K_{a1}(\ce{COOH}) = 2,34$, $\text{p}K_{a2}(\ce{NH3+}) = 9,60$. Calculer le pI théorique et comparer à la valeur mesurée au T3.
    \item \textbf{Hydrolyse et liaison peptidique.} Écrire l'équation-bilan schématique de l'hydrolyse acide d'un dipeptide $\ce{Ala-Gly}$ (alanine-glycine). Identifier la liaison coupée.
    \item \textbf{Spécificité des tests.} Expliquer pourquoi le test de Biuret devient négatif (ou faiblement positif) après hydrolyse prolongée, alors que le test à la ninhydrine devient fortement positif.
    \item \textbf{Contexte camerounais.} Le \emph{bili-bili} (bière de mil) et le \emph{vin de palme} contiennent des acides aminés issus de l'hydrolyse enzymatique des protéines du mil ou de la sève par les levures/bactéries. Pourquoi la teneur en acides aminés libres augmente-t-elle pendant la fermentation ?
\end{enumerate}

\begin{exoresolu}[Éléments de correction guidée]
\textbf{1. Caractère amphotère (amphion $\ce{^{+}H3N-CH2-COO-}$) :}
\begin{itemize}
    \item \textbf{Caractère acide (don de H⁺ sur fonction ammonium) :}
    \[ \ce{^{+}H3N-CH2-COO- + H2O <=> H2N-CH2-COO- + H3O+} \quad \text{Couple : } \ce{^{+}H3N-CH2-COO-}/\ce{H2N-CH2-COO-} \quad (\text{p}K_{a2} = 9,60) \]
    \item \textbf{Caractère basique (capture de H⁺ sur fonction carboxylate) :}
    \[ \ce{^{+}H3N-CH2-COO- + H3O+ <=> ^{+}H3N-CH2-COOH + H2O} \quad \text{Couple : } \ce{^{+}H3N-CH2-COOH}/\ce{^{+}H3N-CH2-COO-} \quad (\text{p}K_{a1} = 2,34) \]
    \emph{Remarque :} Avec NaOH, on écrit $\ce{^{+}H3N-CH2-COO- + HO- -> H2N-CH2-COO- + H2O}$.
\end{itemize}

\textbf{2. Équilibre forme moléculaire / amphion :}
\[ \ce{H2N-CH2-COOH <=> ^{+}H3N-CH2-COO-} \]
C'est un transfert de proton \emph{intramoléculaire}. La constante d'équilibre $K = \frac{[\text{amphion}]}{[\text{moléculaire}]} = 10^{\text{p}K_{a2} - \text{p}K_{a1}}$ (très grand, $\approx 10^{7}$). La forme zwitterionique est \textbf{prédominante à plus de 99,999\%} en solution aqueuse neutre. La forme neutre n'existe quasi pas en solution (elle existe dans le gaz ou solvants apolaires).

\textbf{3. Calcul du pI :}
\[ \text{pI} = \frac{\text{p}K_{a1} + \text{p}K_{a2}}{2} = \frac{2,34 + 9,60}{2} = \frac{11,94}{2} = \mathbf{5,97} \]
Valeur mesurée au T3 (\SI{5,5}{--}\SI{6,5}) : \textbf{accord excellent}. La légère différence vient de la force ionique, température, et imprécision papier pH.

\textbf{4. Hydrolyse acide de l'Ala-Gly :}
\[ \ce{CH3-CH(NH2)-CO-NH-CH2-COOH + H2O ->[H+/\Delta] CH3-CH(NH2)-COOH + H2N-CH2-COOH} \]
\emph{Alanine} \quad \emph{Glycine}
La liaison coupée est la \cle{liaison peptidique (amide)} : $\ce{-CO-NH-}$. L'hydrolyse acide protonne le carbone carbonyl, facilitant l'attaque nucléophile par l'eau.

\textbf{5. Biuret vs Ninhydrine :}
\begin{itemize}
    \item \textbf{Biuret :} Complexe $\ce{Cu^{2+}}$ avec les \textbf{liaisons peptidiques} ($\ge 2$ liaisons $\ce{-CO-NH-}$ nécessaires pour une coloration franche). Protéine native : $\gg 2$ liaisons $\to$ violet intense. Hydrolysat total : 0 liaison peptidique (acides aminés libres) $\to$ \textbf{négatif}. Hydrolysat partiel (20 min) : peptides courts (quelques liaisons) $\to$ violet pâle.
    \item \textbf{Ninhydrine :} Réagit avec le \textbf{groupe $\alpha$-amine libre} ($\ce{-NH2}$) des acides α-aminés (et N-terminal des peptides). Protéine native : très peu de groupes $\alpha$-amine libres (seuls les N-terminaux) $\to$ faible coloration. Hydrolysat total : \textbf{tous} les acides aminés libèrent leur $\alpha$-amine $\to$ \textbf{violet intense (Ruhemann)}.
\end{itemize}

\textbf{6. Fermentation (Bili-bili / Vin de palme) :}
Les microorganismes (levures \emph{Saccharomyces}, bactéries lactiques) sécrètent des \textbf{protéases} (enzymes) qui hydrolysent les protéines du substrat (mil, maïs, sève de palmier) en peptides puis acides aminés libres. Ces AA servent de \textbf{source d'azote} pour la synthèse protéique des microbes (croissance) et précurseurs d'arômes (acides aminés $\to$ aldéhydes, alcools supérieurs par déamination/décarboxylation : voie d'Ehrlich). L'augmentation des AA libres est donc marqueur de l'activité enzymatique microbienne.
\end{exoresolu}

\courssub{Conclusion}
Ce TP a permis de valider expérimentalement deux notions fondamentales du programme de Terminale TI :
\begin{enumerate}
    \item Le \textbf{caractère amphotère} des acides α-aminés : la glycine, sous sa forme amphion (zwitterion) à pH neutre (pH $\approx$ pI = 6,0), réagit avec l'acide (comportement basique de la fonction carboxylate) et avec la base (comportement acide de la fonction ammonium). Cette double réactivité est la signature chimique des acides α-aminés.
    \item La \textbf{structure polymérique des protéines} : l'hydrolyse acide (HCl \SI{6}{mol.L^{-1}}, reflux 20 min) casse les liaisons peptidiques ($\ce{-CO-NH-}$), libérant les monomères (acides α-aminés). Le virage du test de Biuret (positif $\to$ négatif) couplé au virage de la ninhydrine (négatif $\to$ violet intense) constitue une preuve chimique robuste de cette dépolymérisation.
\end{enumerate}
\begin{aretenir}[Compétences évaluées au Bac TI]
\begin{itemize}
    \item Écrire les équations-bilans des réactions acide/base de l'amphion (identifier les couples).
    \item Calculer le pI à partir des pKa et l'interpréter.
    \item Écrire l'équation-bilan d'hydrolyse d'un dipeptide (ou peptide) en milieu acide.
    \item Interpréter les tests de Biuret (liaison peptidique) et Ninhydrine (amine α-libre) avant/après hydrolyse.
    \item Respecter le protocole sécurité (manipulation acide/base concentré, chauffage à reflux).
\end{itemize}
\end{aretenir}

\begin{savaistu}[Du laboratoire à l'industrie : SONARA, Savonneries, Pharmacopée]
\begin{itemize}
    \item \textbf{SONARA (Limbe) :} Les catalyseurs de craquage (zéolites) fonctionnent par sites acides de Brønsted/Lewis. La compréhension de l'amphotérie (sites acides/bases sur surface solide) est transposée de la chimie moléculaire (acides aminés) à la catalyse hétérogène industrielle.
    \item \textbf{Savonneries artisanales (Douala, Bafoussam) :} L'hydrolyse basique (saponification) des triglycérides (huile de palme) libère du glycérol et des carboxylates (savon). C'est l'analogue exact de l'hydrolyse acide des protéines : rupture de liaison ester vs amide, catalyse par $\ce{HO-}$ vs $\ce{H3O+}$.
    \item \textbf{Pharmacopée camerounaise :} L'écorce de \emph{Prunus africana} (Pygeum) ou les feuilles de \emph{Moringa oleifera} sont riches en acides aminés libres et peptides bioactifs. Les tradipraticiens utilisent des décoctions (hydrolyse thermique aqueuse) pour extraire ces principes actifs. La ninhydrine est utilisée en contrôle qualité (labo LANACOME, IRAD) pour doser l'azote aminé.
    \item \textbf{Agro-alimentaire (Brasseries, Chococam) :} Le contrôle de l'hydrolyse des protéines (malt, cacao) pendant le brassage ou la fermentation du cacao détermine le profil aromatique (précurseurs de Maillard : acides aminés + sucres réducteurs $\to$ arômes torréfiés).
\end{itemize}
\end{savaistu}

\newpage

\coursec{Méthodes et exercices résolus}

\coursec{Introduction : Les acides aminés, briques des protéines}

Les protéines, macromolécules essentielles à la vie (enzymes, anticorps, hémoglobine, collagène), sont des polymères naturels constitués d'unités répétitives : les \cle{acides $\alpha$-aminés}. Au Cameroun, l'industrie agroalimentaire (brasseries, laiteries, transformation du poisson à Kribi ou Douala) et l'industrie pharmaceutique (laboratoires de Douala, Yaoundé) manipulent ces molécules pour la nutrition, la santé ou la synthèse de médicaments comme l'aspartame (édulcorant dipeptidique). Comprendre leur structure, leur comportement amphotère et leur mode de liaison (liaison peptidique) est indispensable pour tout technicien supérieur en chimie industrielle.

\coursec{Structure et définition des acides $\alpha$-aminés}

\courssub{Méthode 1 — Reconnaître un acide $\alpha$-aminé et repérer le carbone asymétrique}

\begin{methode}[Identifier un acide $\alpha$-aminé et son carbone asymétrique]
\begin{enumerate}
\item Cherche dans la molécule la présence simultanée d'un groupe \cle{carboxyle} \ce{-COOH} et d'un groupe \cle{amine} \ce{-NH2}.
\item Vérifie que les deux groupes sont portés par le \textbf{même atome de carbone} : c'est le carbone $\alpha$ (le carbone adjacent au \ce{-COOH}). La formule générale s'écrit :
\[ \mathrm{H_2N{-}CHR{-}COOH} \quad \text{avec } \mathrm{R} = \text{chaîne latérale (atome H ou groupe alkyl, aryle, etc.)} \]
\item Pour savoir si le carbone $\alpha$ est \cle{asymétrique}, liste ses \textbf{quatre substituants} : \ce{-NH2}, \ce{-H}, \ce{-COOH} et \ce{-R}.
\item Conclus : le carbone $\alpha$ est asymétrique si et seulement si les quatre groupes sont \textbf{tous différents}, c'est-à-dire si $\mathrm{R} \neq \mathrm{H}$. Le seul acide $\alpha$-aminé protéinogène non chiral est la \cle{glycine} ($\mathrm{R} = \mathrm{H}$).
\item Marque le carbone asymétrique d'une étoile ($C^{*}$) dans la formule topologique.
\end{enumerate}
\end{methode}

\begin{exoresolu}[Acide aminé ou pas ? Carbone asymétrique ou pas ?]
On considère les trois composés suivants :
\begin{center}
A : \chemfig{H_2N-CH_2-CH_2-C(=[1]O)-[7]OH} \qquad
B : \chemfig{CH_3-CH(-[2]NH_2)-C(=[1]O)-[7]OH} \qquad
C : \chemfig{H_2N-CH_2-C(=[1]O)-[7]OH}
\end{center}
\begin{enumerate}
\item Lesquels sont des acides $\alpha$-aminés ? Justifier.
\item Parmi eux, lequel possède un carbone asymétrique ? Le repérer par une étoile.
\end{enumerate}
\tcblower
\textbf{1. Identification des acides $\alpha$-aminés.}

\textbf{Composé A} : le groupe \ce{-NH2} est porté par le carbone en position $\beta$ (deuxième carbone à partir du \ce{-COOH}), non par le carbone $\alpha$. Ce n'est \textbf{pas} un acide $\alpha$-aminé (c'est l'acide $\beta$-alanine, un acide $\beta$-aminé).

\textbf{Composé B} : le carbone adjacent au \ce{-COOH} porte à la fois \ce{-NH2} et \ce{-H} ; le groupe R est \ce{-CH3}. C'est bien un \textbf{acide $\alpha$-aminé} : l'\cle{alanine}.

\textbf{Composé C} : le carbone $\alpha$ porte \ce{-NH2}, \ce{-COOH} et \textbf{deux} atomes d'hydrogène ($\mathrm{R}=\mathrm{H}$). C'est bien un acide $\alpha$-aminé : la \cle{glycine}.

\textbf{2. Carbone asymétrique.}

Pour B (alanine), les quatre substituants du carbone $\alpha$ sont : \ce{-NH2}, \ce{-H}, \ce{-COOH}, \ce{-CH3} : tous différents $\Rightarrow$ carbone asymétrique :
\[ \chemfig{H_2N-\chemabove{C}{\scriptstyle *}(-[2]H)(-[6]CH_3)-C(=[1]O)-[7]OH} \]
Pour C (glycine), le carbone $\alpha$ porte deux H identiques : il n'est \textbf{pas} asymétrique.

\textbf{Conclusion} : B et C sont des acides $\alpha$-aminés ; seule l'alanine (B) possède un carbone asymétrique et est donc chirale.
\end{exoresolu}

\begin{savaistu}[Les acides aminés dans l'alimentation camerounaise]
Les protéines de nos aliments quotidiens (poisson \emph{kpém}, haricots \emph{haricots rouges}, viande, \emph{ndolé}, arachides) sont hydrolysées lors de la digestion en acides aminés libres. Certains sont \cle{essentiels} (non synthétisés par l'organisme : leucine, lysine, valine, etc.) et doivent être apportés par l'alimentation. La complémentarité céréales-légumineuses (riz + haricots, maïs + soja) assure un profil en acides aminés complet, base de la sécurité nutritionnelle.
\end{savaistu}

\coursec{Nomenclature et représentation de Fischer}

\courssub{Méthode 2 — Représentation de Fischer et configurations L et D}

\begin{methode}[Construire la représentation de Fischer et attribuer la configuration L ou D]
\begin{enumerate}
\item Place le \textbf{carbone asymétrique} ($C^{*}$) au centre d'une croix (liaisons verticales / horizontales).
\item Règle d'orientation (convention Fischer pour les acides aminés) : la \textbf{chaîne carbonée principale} est verticale, avec le groupe \ce{-COOH} \textbf{en haut} (carbone C1) et le groupe R \textbf{en bas}.
\item Place \ce{-NH2} et \ce{-H} sur l'horizontale (gauche/droite). Rappel : en projection de Fischer, les liaisons horizontales sortent du plan vers l'observateur, les liaisons verticales partent vers l'arrière.
\item Lecture de la configuration (avec \ce{-COOH} en haut) :
\begin{itemize}
\item \ce{-NH2} à \textbf{gauche} $\Rightarrow$ configuration \cle{L} (série L, la plus fréquente dans les protéines naturelles) ;
\item \ce{-NH2} à \textbf{droite} $\Rightarrow$ configuration \cle{D}.
\end{itemize}
\item \textbf{Interdiction} de faire pivoter la projection de \ang{90} dans le plan (cela inverserait la configuration) ; seule une rotation de \ang{180} la conserve.
\end{enumerate}
\end{methode}

\begin{exoresolu}[Configurations de l'alanine et de la valine]
\begin{enumerate}
\item Donner la représentation de Fischer de la L-alanine ($\mathrm{R} = \mathrm{-CH_3}$).
\item On donne ci-dessous la représentation de Fischer d'un acide $\alpha$-aminé. Identifier cet acide aminé et préciser sa configuration (L ou D).
\begin{center}
\chemfig{HOOC-[:-90]C(-[:0]NH_2)(-[:180]H)-[:-90]CH(-[:180]CH_3)(-[:0]CH_3)}
\end{center}
\end{enumerate}
\tcblower
\textbf{1. L-alanine.} Chaîne verticale : \ce{-COOH} en haut, \ce{-CH3} en bas ; configuration L $\Rightarrow$ \ce{-NH2} à gauche, \ce{-H} à droite :
\begin{center}
\chemfig{HOOC-[:-90]C(-[:180]NH_2)(-[:0]H)-[:-90]CH_3}
\end{center}
(Lecture : \ce{-NH2} à gauche $\Rightarrow$ L-alanine).

\textbf{2. Identification.} Le groupe R (en bas) est \ce{-CH(CH3)2} (groupe isopropyle) : il s'agit de la \textbf{valine}. Dans la représentation donnée, le \ce{-COOH} est en haut, R en bas, et le groupe \ce{-NH2} est à \textbf{droite} : c'est donc la \textbf{D-valine}.

\textbf{Conclusion} : la L-alanine a \ce{-NH2} à gauche ; la molécule proposée est la D-valine.
\end{exoresolu}

\begin{experience}[Mise en évidence des acides aminés : réaction à la ninhydrine]
\textbf{Objectif} : Détecter la présence de groupes \ce{-NH2} sur les acides $\alpha$-aminés.
\textbf{Matériel} : Tubes à essai, solution de ninhydrine (\SI{0,2}{\%} dans l'éthanol), bain-marie, glycine, alanine, eau distillée.
\textbf{Protocole} :
\begin{enumerate}
\item Déposer \SI{1}{mL} de solution aqueuse de glycine (\SI{0,1}{M}) dans un tube, \SI{1}{mL} d'alanine (\SI{0,1}{M}) dans un second, \SI{1}{mL} d'eau dans un troisième (témoin).
\item Ajouter \SI{1}{mL} de solution de ninhydrine dans chaque tube.
\item Chauffer au bain-marie bouillant pendant \SI{5}{min}.
\end{enumerate}
\textbf{Observations} : Les tubes contenant les acides aminés prennent une coloration \textbf{bleu-violet intense} (coloration de Ruhemann) ; le témoin reste incolore/jaune pâle.
\textbf{Interprétation} : La ninhydrine oxydativement déamine l'acide $\alpha$-aminé (libérant \ce{CO2}, \ce{NH3} et un aldéhyde) et forme un produit coloré. C'est un test de caractérisation et de dosage (colorimétrie) des acides aminés libres.
\end{experience}

\coursec{Propriétés acido-basiques : l'amphion}

\courssub{Méthode 3 — Écrire l'équilibre forme moléculaire / amphion}

\begin{methode}[Écrire l'équilibre entre la forme moléculaire et l'amphion]
\begin{enumerate}
\item Pars de la \textbf{forme moléculaire} (neutre) : $\mathrm{H_2N{-}CHR{-}COOH}$.
\item Identifie le \textbf{site donneur de proton} : le groupe carboxyle \ce{-COOH} (acide de Brønsted) ; et le \textbf{site accepteur} : le groupe amine \ce{-NH2} (base de Brønsted, grâce au doublet non liant de l'azote).
\item Le transfert de proton est \textbf{intramoléculaire} : le proton du \ce{-COOH} migre vers le \ce{-NH2}.
\item Écris l'équation-bilan avec une \textbf{double flèche} (équilibre) :
\[ \mathrm{H_2N{-}CHR{-}COOH} \ce{<=>} \mathrm{^+H_3N{-}CHR{-}COO^-} \]
\item Vérifie la structure de l'\cle{amphion} (ou \cle{zwitterion}) : charge $+$ portée par l'azote (\ce{-NH3+}), charge $-$ portée par l'oxygène du carboxylate (\ce{-COO-}) ; charge globale \textbf{nulle}.
\item Retiens : en solution aqueuse et à l'état solide, l'acide $\alpha$-aminé existe \textbf{quasi exclusivement} sous forme d'amphion (ce qui explique les températures de fusion élevées et la bonne solubilité dans l'eau).
\end{enumerate}
\end{methode}

\begin{exoresolu}[L'amphion de la glycine]
La glycine ($\mathrm{R} = \mathrm{H}$) est un solide blanc qui fond vers \SI{233}{\celsius} et qui est très soluble dans l'eau, contrairement à la plupart des acides carboxyliques de masse molaire voisine.
\begin{enumerate}
\item Écrire l'équation-bilan de l'équilibre entre la forme moléculaire et la forme ionique de la glycine.
\item Représenter l'amphion obtenu en précisant la localisation des charges.
\item Expliquer qualitativement la température de fusion élevée.
\end{enumerate}
\tcblower
\textbf{1. Équation-bilan.}
\[ \mathrm{H_2N{-}CH_2{-}COOH} \ce{<=>} \mathrm{^+H_3N{-}CH_2{-}COO^-} \]
Le proton du groupe carboxyle est transféré sur le doublet non liant de l'azote du groupe amine.

\textbf{2. Représentation de l'amphion.}
\begin{center}
\begin{popfigure}
\begin{tikzpicture}[scale=0.8, every node/.style={font=\small}]
\node (N) at (0,0) {$\ce{^{+}H3N}$};
\node (C) at (2,0) {$\ce{CH2}$};
\node (O1) at (4,0.5) {$\ce{O^{-}}$};
\node (O2) at (4,-0.5) {$\ce{O}$};
\draw[thick] (N) -- (C);
\draw[thick] (C) -- (3,0);
\draw[thick] (3,0) -- (O1);
\draw[thick] (3,0) -- (O2);
\draw[dashed, popblue] (3,0) -- (4,0) node[midway, above] {Résonance};
\end{tikzpicture}
\legende{Représentation de l'amphion de la glycine (ion zwitterion) avec délocalisation de la charge négative sur le carboxylate.}
\end{popfigure}
\end{center}
Charge $+$ sur l'azote (ion ammonium \ce{-NH3+}), charge $-$ délocalisée sur les deux oxygènes du groupe carboxylate \ce{-COO-}. Charge globale nulle : c'est un ion dipolaire.

\textbf{3. Interprétation.} À l'état solide, les amphions s'attirent fortement par interactions électrostatiques (liaisons ioniques entre têtes \ce{NH3+} et \ce{COO-} de molécules voisines), ce qui forme un réseau cohésif difficile à rompre : d'où une température de fusion élevée (\SI{233}{\celsius}) et une bonne solubilité dans l'eau, solvant polaire qui solvate les ions.

\textbf{Conclusion} : la glycine existe essentiellement sous la forme zwitterionique $\mathrm{^+H_3N{-}CH_2{-}COO^-}$, ce qui explique ses propriétés physiques de type « sel ».
\end{exoresolu}

\coursec{Comportement amphotère de l'amphion}

\courssub{Méthode 4 — Caractère amphotère de l'amphion}

\begin{methode}[Écrire les équations acide/base de l'amphion]
\begin{enumerate}
\item L'amphion est un \cle{ampholyte} : il peut \textbf{céder} un proton (comportement acide) ou en \textbf{capter} un (comportement basique).
\item \textbf{Comportement acide} (face à une base, ex. \ce{HO-}) : c'est le groupe \ce{-NH3+} qui cède son proton :
\[ \mathrm{^+H_3N{-}CHR{-}COO^-} + \ce{HO^-} \ce{->} \mathrm{H_2N{-}CHR{-}COO^-} + \ce{H2O} \]
L'espèce formée est un \textbf{anion} (charge $-1$), prédominant en milieu basique (pH > pI).
\item \textbf{Comportement basique} (face à un acide, ex. \ce{H3O+}) : c'est le groupe \ce{-COO-} qui capte le proton :
\[ \mathrm{^+H_3N{-}CHR{-}COO^-} + \ce{H3O^+} \ce{->} \mathrm{^+H_3N{-}CHR{-}COOH} + \ce{H2O} \]
L'espèce formée est un \textbf{cation} (charge $+1$), prédominant en milieu acide (pH < pI).
\item Vérifie toujours la \textbf{conservation des charges} et des atomes de part et d'autre de la flèche.
\end{enumerate}
\end{methode}

\begin{exoresolu}[L'amphion de l'alanine en milieu acide et en milieu basique]
On dissout de l'alanine dans l'eau ; elle s'y trouve sous forme d'amphion noté $\mathrm{^+H_3N{-}CH(CH_3){-}COO^-}$.
\begin{enumerate}
\item Écrire l'équation-bilan de la réaction de cet amphion avec les ions oxonium \ce{H3O+} d'une solution d'acide chlorhydrique. Nommer l'espèce formée.
\item Écrire l'équation-bilan de la réaction de cet amphion avec les ions hydroxyde \ce{HO^-} d'une solution d'hydroxyde de sodium. Nommer l'espèce formée.
\item Justifier le qualificatif d'\emph{ampholyte} donné à l'amphion.
\end{enumerate}
\tcblower
\textbf{1. En milieu acide.} L'amphion joue le rôle de \textbf{base} : le groupe carboxylate capte un proton :
\[ \mathrm{^+H_3N{-}CH(CH_3){-}COO^-} + \ce{H3O^+} \ce{->} \mathrm{^+H_3N{-}CH(CH_3){-}COOH} + \ce{H2O} \]
L'espèce formée est un \textbf{cation} (charge globale $+1$). Vérification des charges : à gauche $(+1-1)+(+1)=+1$ ; à droite $(+1)+0=+1$. Conservée.

\textbf{2. En milieu basique.} L'amphion joue le rôle d'\textbf{acide} : le groupe ammonium cède un proton :
\[ \mathrm{^+H_3N{-}CH(CH_3){-}COO^-} + \ce{HO^-} \ce{->} \mathrm{H_2N{-}CH(CH_3){-}COO^-} + \ce{H2O} \]
L'espèce formée est un \textbf{anion} (charge globale $-1$). Vérification : à gauche $0+(-1)=-1$ ; à droite $(-1)+0=-1$. Conservée.

\textbf{3. Justification.} L'amphion se comporte comme un acide vis-à-vis d'une base (\ce{HO-}) et comme une base vis-à-vis d'un acide (\ce{H3O+}) : c'est la définition d'un \textbf{ampholyte} (espèce amphotère).

\textbf{Conclusion} : selon le pH du milieu, l'alanine existe sous forme cationique (milieu acide), d'amphion (pH intermédiaire, voisin du point isoélectrique) ou anionique (milieu basique).
\end{exoresolu}

\begin{aretenir}[Formes prédominantes selon le pH]
Pour un acide $\alpha$-aminé simple (chaîne latérale non ionisable) :
\begin{itemize}
\item pH $<$ pI (milieu acide) : forme \textbf{cationique} $\mathrm{^+H_3N{-}CHR{-}COOH}$.
\item pH $=$ pI (point isoélectrique) : forme \textbf{amphion (zwitterion)} $\mathrm{^+H_3N{-}CHR{-}COO^-}$ (solubilité minimale, migration nulle en électrophorèse).
\item pH $>$ pI (milieu basique) : forme \textbf{anionique} $\mathrm{H_2N{-}CHR{-}COO^-}$.
\end{itemize}
Le pI (point isoélectrique) est le pH où la concentration moyenne des charges positives égale celle des charges négatives. Pour la glycine, pI $\approx$ 5,97 ; pour l'alanine, pI $\approx$ 6,00.
\end{aretenir}

\coursec{La liaison peptidique et les dipeptides}

\courssub{Méthode 5 — Écrire la formation d'un dipeptide et nommer les peptides}

\begin{methode}[Former et nommer un dipeptide]
\begin{enumerate}
\item La \cle{liaison peptidique} est une liaison \textbf{amide} \ce{-CO-NH-} qui se forme entre le groupe \ce{-COOH} du premier acide aminé et le groupe \ce{-NH2} du second, avec \textbf{élimination d'une molécule d'eau} (réaction de condensation).
\item Écris l'équation-bilan générale (souvent écrite avec une flèche simple $\ce{->}$ en synthèse préparative car on chasse l'eau ou on active le carboxyle, mais c'est un équilibre limité en milieu aqueux) :
\[ \mathrm{H_2N{-}CHR_1{-}COOH} + \mathrm{H_2N{-}CHR_2{-}COOH} \ce{->} \mathrm{H_2N{-}CHR_1{-}CO{-}NH{-}CHR_2{-}COOH} + \ce{H2O} \]
\item Repère dans le produit : l'extrémité \textbf{N-terminale} (groupe \ce{-NH2} libre, à gauche) et l'extrémité \textbf{C-terminale} (groupe \ce{-COOH} libre, à droite).
\item \textbf{Nomenclature} : on nomme le dipeptide en citant d'abord le résidu N-terminal (nom modifié en \emph{-yle} : alan$\rightarrow$alanyle, glycin$\rightarrow$glycyle, valin$\rightarrow$valyle) puis l'acide aminé C-terminal (nom inchangé). Exemple : \cle{glycylalanine} (Gly-Ala).
\item Attention à l'\textbf{ordre} : Ala-Gly et Gly-Ala sont deux dipeptides \textbf{différents} (isomères de constitution).
\end{enumerate}
\end{methode}

\begin{exoresolu}[Synthèse et dénombrement de dipeptides]
On fait réagir la glycine $\mathrm{H_2N{-}CH_2{-}COOH}$ (notée Gly) avec l'alanine $\mathrm{H_2N{-}CH(CH_3){-}COOH}$ (notée Ala) dans des conditions permettant la condensation.
\begin{enumerate}
\item Écrire l'équation-bilan de formation du dipeptide Gly-Ala et entourer la liaison peptidique.
\item Combien de dipeptides différents peut-on obtenir au total à partir de ce mélange de deux acides aminés ? Les nommer.
\item On réalise la condensation de \SI{7,5}{g} de glycine ($M = \SI{75}{g.mol^{-1}}$) avec un excès d'alanine, en supposant la réaction totale et sélective vers Gly-Ala ($M = \SI{146}{g.mol^{-1}}$). Calculer la masse de dipeptide obtenue.
\end{enumerate}
\tcblower
\textbf{1. Équation-bilan.}
\[ \mathrm{H_2N{-}CH_2{-}COOH} + \mathrm{H_2N{-}CH(CH_3){-}COOH} \ce{->} \mathrm{H_2N{-}CH_2{-}CO{-}NH{-}CH(CH_3){-}COOH} + \ce{H2O} \]
La liaison peptidique est le groupe \ce{-CO-NH-} central (surligné ci-dessus), formé entre le carboxyle de la glycine et l'amine de l'alanine, avec départ d'une molécule d'eau.

\textbf{2. Dénombrement.} Chaque extrémité peut être occupée par Gly ou Ala : $2 \times 2 = 4$ dipeptides :
\begin{center}
\begin{tabularx}{\linewidth}{|l|X|}
\hline
\rowcolor{popblueL} \textbf{Dipeptide} & \textbf{Nom systématique} \\
\hline
Gly-Gly & glycylglycine \\
\hline
Ala-Ala & alanylalanine \\
\hline
Gly-Ala & glycylalanine \\
\hline
Ala-Gly & alanylglycine \\
\hline
\end{tabularx}
\end{center}
Gly-Ala et Ala-Gly sont distincts car les extrémités N- et C-terminales sont inversées.

\textbf{3. Calcul de masse.} Quantité de glycine (réactif limitant) :
\[ n(\text{Gly}) = \frac{m}{M} = \frac{7,5}{75} = \SI{0,10}{mol} \]
Tableau d'avancement (réaction totale, sélective) :
\begin{center}
\begin{tabular}{|l|c|c|c|c|}
\hline
\rowcolor{popblueL} & Gly & Ala & Gly-Ala & \ce{H2O} \\
\hline
État initial (mol) & 0,10 & excès & 0 & 0 \\
\hline
État final (mol) & 0 & excès & 0,10 & 0,10 \\
\hline
\end{tabular}
\end{center}
D'après la stœchiométrie $1:1$ : $n(\text{Gly-Ala}) = \SI{0,10}{mol}$, d'où :
\[ m(\text{Gly-Ala}) = n \times M = 0,10 \times 146 = \SI{14,6}{g} \approx \SI{15}{g} \text{ (2 chiffres significatifs)} \]

\textbf{Conclusion} : on obtient \SI{14,6}{g} de glycylalanine.
\end{exoresolu}

\begin{popfigure}
\begin{tikzpicture}[scale=0.9, every node/.style={font=\small}]
% Peptide bond geometry
\node (N1) at (0,0) {$\ce{H2N}$};
\node (Ca1) at (1.2,0) {$\ce{CH2}$};
\node (C1) at (2.4,0) {$\ce{C}$};
\node (O1) at (2.4,0.7) {$\ce{O}$};
\node (N2) at (3.6,0) {$\ce{NH}$};
\node (Ca2) at (4.8,0) {$\ce{CH}$};
\node (C2) at (6.0,0) {$\ce{C}$};
\node (O2) at (6.0,0.7) {$\ce{O}$};
\node (OH) at (6.0,-0.7) {$\ce{OH}$};
\node (CH3) at (4.8,-1.0) {$\ce{CH3}$};

\draw[thick] (N1) -- (Ca1) -- (C1) -- (N2) -- (Ca2) -- (C2);
\draw[thick] (C1) -- (O1);
\draw[thick] (C2) -- (O2);
\draw[thick] (C2) -- (OH);
\draw[thick] (Ca2) -- (CH3);

% Highlight peptide bond
\draw[poporange, thick, decorate, decoration={brace, amplitude=5pt, mirror}] (2.4,-0.5) -- (3.6,-0.5) node[midway, below=6pt, poporange] {\textbf{Liaison peptidique}};
\draw[poporange, thick] (2.4,-0.3) -- (3.6,-0.3);
\end{tikzpicture}
\legende{Structure du dipeptide Gly-Ala (glycylalanine) avec mise en évidence de la liaison peptidique (amide) et des extrémités N-terminale (libre \ce{NH2}) et C-terminale (libre \ce{COOH}).}
\end{popfigure}

\coursec{Synthèse dirigée d'un dipeptide (complément Bac)}

\courssub{Méthode 6 — Stratégie de synthèse sélective (Protéger -- Activer -- Condenser -- Déprotéger)}

\begin{methode}[Stratégie de synthèse dirigée d'un dipeptide A-B]
\begin{enumerate}
\item \textbf{Problème} : la condensation directe d'un mélange de deux acides aminés A et B donne \textbf{quatre} dipeptides (A-A, B-B, A-B, B-A). Pour n'obtenir que le dipeptide souhaité A-B (A en N-terminal, B en C-terminal), il faut une \textbf{synthèse dirigée}.
\item \textbf{Étape 1 — Protéger} le groupe amine \ce{-NH2} de l'acide A (futur N-terminal) par un groupe protecteur (ex. \cle{Boc} = \textit{tert-butyloxycarbonyl}, \cle{Fmoc} = \textit{fluorenylméthyloxycarbonyl}, ou \cle{Cbz} = \textit{benzyloxycarbonyl}), \textbf{ET/OU} protéger le groupe \ce{-COOH} de l'acide B (futur C-terminal) par estérification (ex. ester méthylique ou \textit{tert}-butylique).
\item \textbf{Étape 2 — Activer} le groupe carboxyle de A (protégé sur l'amine) pour rendre la condensation possible et rapide (ex. chlorure d'acide, anhydride mixte, ou réactifs de couplage type DCC, EDC, HATU).
\item \textbf{Étape 3 — Condenser} A (protégé/activé) avec B (protégé sur l'acide) : une seule liaison peptidique peut se former, celle voulue (A-B).
\item \textbf{Étape 4 — Déprotéger} : régénérer les groupes \ce{-NH2} et \ce{-COOH} libres sans rompre la liaison peptidique (ex. TFA pour Boc, pipéridine pour Fmoc, hydrolyse basique douce pour l'ester).
\item Schéma réflexe à retenir : \emph{Protéger $\rightarrow$ Activer $\rightarrow$ Condenser $\rightarrow$ Déprotéger}.
\end{enumerate}
\end{methode}

\begin{exoresolu}[Synthèse dirigée de l'alanylglycine (Ala-Gly)]
Un laboratoire pharmaceutique de Douala souhaite préparer \textbf{uniquement} le dipeptide alanylglycine (Ala-Gly) à partir d'alanine et de glycine.
\begin{enumerate}
\item Expliquer pourquoi la simple condensation du mélange alanine + glycine ne convient pas.
\item Indiquer, pour obtenir sélectivement Ala-Gly, quel groupe fonctionnel de l'alanine doit être protégé et quel groupe de la glycine doit être bloqué. Justifier.
\item La liaison peptidique est une liaison amide. Citer une propriété de cette liaison qui impose des précautions lors de l'étape de déprotection.
\end{enumerate}
\tcblower
\textbf{1. Inconvénient de la condensation directe.} Chaque acide aminé possède à la fois un groupe donneur (\ce{-COOH}) et un groupe accepteur (\ce{-NH2}) : la condensation non dirigée conduit à un mélange de \textbf{quatre} dipeptides : Ala-Ala, Gly-Gly, Ala-Gly et Gly-Ala. Le rendement en Ala-Gly serait faible (\SI{25}{\%} statistique) et la séparation coûteuse.

\textbf{2. Stratégie de protection pour Ala-Gly.} Dans Ala-Gly, c'est le \ce{-COOH} de l'alanine qui réagit avec le \ce{-NH2} de la glycine. Il faut donc :
\begin{itemize}
\item \textbf{Protéger le groupe amine \ce{-NH2} de l'alanine} (ex. \ce{Boc-NH-Ala-OH}) pour l'empêcher de réagir comme nucléophile ;
\item \textbf{Protéger le groupe carboxyle \ce{-COOH} de la glycine} (ex. \ce{H-Gly-OMe} ou \ce{H-Gly-OtBu}) pour l'empêcher de réagir comme électrophile ;
\item Activer le \ce{-COOH} de l'alanine protégée (ex. en chlorure d'acide ou via DCC), condenser avec l'ester de glycine, puis déprotéger les deux extrémités (amine puis ester, ou simultanément).
\end{itemize}

\textbf{3. Précaution lors de la déprotection.} La liaison peptidique (amide) est \textbf{hydrolysable} en milieu acide fort concentré et chaud, ou en milieu basique fort et chaud. Lors de la déprotection, il faut choisir des conditions \textbf{orthogonales} (douces et sélectives) pour cliver les groupes protecteurs \textbf{sans hydrolyser la liaison peptidique} fraîchement formée (ex. TFA \SI{50}{\%} à froid pour Boc, hydrolyse basique \ce{LiOH} dilué à \SI{0}{\celsius} pour l'ester méthylique).

\textbf{Conclusion} : la séquence \emph{protéger -- activer -- condenser -- déprotéger} permet d'obtenir sélectivement Ala-Gly, à condition de préserver la liaison amide lors de la déprotection.
\end{exoresolu}

\begin{attention}[Les 5 erreurs qui coûtent des points au Bac]
\begin{enumerate}
\item \textbf{Confondre carbone $\alpha$ et carbone asymétrique.} Tout acide $\alpha$-aminé possède un carbone $\alpha$, mais la glycine ($\mathrm{R} = \mathrm{H}$) n'a \textbf{pas} de carbone asymétrique. Vérifie toujours que les quatre substituants sont différents avant de conclure à la chiralité.
\item \textbf{Inverser L et D en représentation de Fischer.} Le \ce{-COOH} doit être \textbf{en haut} ; alors seulement : \ce{-NH2} à gauche $=$ L, à droite $=$ D. Si tu places la chaîne autrement, la règle ne s'applique plus.
\item \textbf{Mal placer les charges de l'amphion.} La charge $+$ est sur l'azote (\ce{-NH3+}) et la charge $-$ sur l'oxygène (\ce{-COO-}) ; l'amphion est globalement \textbf{neutre}. Écrire $\ce{H2N-CHR-COO^-}$ comme amphion est faux.
\item \textbf{Oublier la molécule d'eau} dans l'équation de formation du dipeptide, ou écrire une flèche simple sans préciser qu'il s'agit d'une condensation (équilibre limité en milieu aqueux, déplacé par activation en synthèse).
\item \textbf{Croire que deux acides aminés ne donnent qu'un seul dipeptide.} Sans stratégie de protection, un mélange A + B donne \textbf{quatre} dipeptides ; et Ala-Gly $\neq$ Gly-Ala. Au Bac, toute question de « synthèse » exige la séquence protéger--activer--condenser--déprotéger.
\end{enumerate}
\end{attention}

\begin{exercice}[Entraînement Bac — Acides aminés et peptides]
\begin{enumerate}
\item On considère la sérine : $\mathrm{HO{-}CH_2{-}CH(NH_2){-}COOH}$.
\begin{enumerate}
\item Écrire sa formule topologique en faisant apparaître le carbone asymétrique. Donner sa configuration absolue (R ou S) en sachant que c'est un acide aminé L.
\item Écrire l'équation-bilan de sa réaction avec l'eau pour former l'amphion. Représenter l'amphion.
\end{enumerate}
\item On souhaite synthétiser le dipeptide \textbf{sérylglycine} (Ser-Gly).
\begin{enumerate}
\item Proposer une stratégie de synthèse dirigée (groupes à protéger, réactifs d'activation possibles).
\item Écrire l'équation-bilan de la condensation entre les dérivés protégés/activés.
\end{enumerate}
\item Un dosage par électrophorèse sur papier à pH = 6,0 sépare un mélange de glycine (pI=5,97), d'alanine (pI=6,00) et d'acide aspartique (pI=2,77). Prédire le sens de migration de chaque acide aminé (vers l'anode $+$, vers la cathode $-$, ou immobile) et justifier.
\end{enumerate}
\end{exercice}

\begin{corrige}[Exercice Entraînement Bac]
\textbf{1. Sérine.}
\begin{enumerate}
\item Carbone $\alpha$ asymétrique (4 substituants différents : \ce{H}, \ce{NH2}, \ce{COOH}, \ce{CH2OH}). Configuration L $\equiv$ (S) pour la sérine (règle de Cahn-Ingold-Prelog : priorités \ce{NH2} > \ce{COOH} > \ce{CH2OH} > \ce{H} ; \ce{NH2} à gauche en Fischer L $\Rightarrow$ sens trigonométrique $\Rightarrow$ S).
\item $\mathrm{HO{-}CH_2{-}CH(NH_2){-}COOH} \ce{<=>} \mathrm{HO{-}CH_2{-}CH(^{+}NH_3){-}COO^-}$.
\end{enumerate}
\textbf{2. Synthèse Ser-Gly.}
\begin{enumerate}
\item Protéger \ce{NH2} de la sérine (Boc-Ser-OH, protection de l'alcool \ce{OH} optionnelle mais recommandée ex. \ce{Bzl} ou \ce{tBu}). Protéger \ce{COOH} de la glycine (H-Gly-OMe). Activer \ce{COOH} de la sérine (DCC/HOBt ou EDC). Condenser. Déprotéger (TFA puis \ce{LiOH} ou \ce{NaOH} dilué).
\item $\ce{Boc-Ser(OBzl)-OH} \xrightarrow[\text{HOBt}]{\text{DCC}} \ce{Boc-Ser(OBzl)-Gly-OMe} \xrightarrow{\text{TFA}} \ce{H-Ser(OBzl)-Gly-OMe} \xrightarrow{\ce{LiOH}} \ce{H-Ser-Gly-OH}$.
\end{enumerate}
\textbf{3. Électrophorèse pH 6,0.}
\begin{itemize}
\item Glycine (pI 5,97) : pH > pI $\Rightarrow$ forme anionique prédominante $\Rightarrow$ migration vers l'\textbf{anode} ($+$).
\item Alanine (pI 6,00) : pH $\approx$ pI $\Rightarrow$ forme amphion prédominante (charge nulle) $\Rightarrow$ \textbf{immobile} (ou très faible migration).
\item Acide aspartique (pI 2,77) : pH > pI $\Rightarrow$ forme anionique (double charge $-$) $\Rightarrow$ migration vers l'\textbf{anode} ($+$), plus rapide que la glycine.
\end{itemize}
\end{corrige}

\newpage

\coursec{Exercices}

\courssub{Je vérifie mes connaissances}

\begin{exercice}[QCM : les acides $\alpha$-aminés]
Pour chaque affirmation, choisir la (ou les) bonne(s) réponse(s) en justifiant brièvement.
\begin{enumerate}
\item Un acide $\alpha$-aminé possède :
\begin{enumerate}
\item un groupe amine et un groupe carboxyle portés par deux carbones différents ;
\item un groupe amine et un groupe carboxyle portés par le même atome de carbone ;
\item uniquement un groupe carboxyle.
\end{enumerate}
\item La formule générale d'un acide $\alpha$-aminé s'écrit :
\begin{enumerate}
\item $R-CH(NH_2)-COOH$ ;
\item $R-CH(OH)-COOH$ ;
\item $R-CO-NH_2$.
\end{enumerate}
\item L'amphion (zwitterion) d'un acide $\alpha$-aminé :
\begin{enumerate}
\item porte une charge globale nulle ;
\item porte une charge positive et une charge négative ;
\item est toujours chargé positivement.
\end{enumerate}
\item La liaison peptidique est une liaison :
\begin{enumerate}
\item ester ; \item amide ; \item éther.
\end{enumerate}
\end{enumerate}
\end{exercice}

\begin{exercice}[Vrai ou faux ?]
Répondre par vrai ou faux en justifiant chaque réponse.
\begin{enumerate}
\item Tous les acides $\alpha$-aminés possèdent un carbone asymétrique.
\item La glycine \ce{H2N-CH2-COOH} est un acide $\alpha$-aminé chirale.
\item En solution aqueuse, un acide $\alpha$-aminé existe sous forme d'un équilibre entre sa forme moléculaire et sa forme ionique (amphion).
\item L'amphion peut réagir aussi bien avec un acide qu'avec une base : on dit qu'il a un comportement amphotère.
\item La réaction de formation d'un dipeptide à partir de deux acides $\alpha$-aminés s'accompagne de l'élimination d'une molécule d'eau.
\end{enumerate}
\end{exercice}

\begin{exercice}[Texte à trous]
Compléter le texte suivant avec les mots ou expressions : \emph{amphion}, \emph{carbone asymétrique}, \emph{amide}, \emph{zwitterion}, \emph{condensation}, \emph{peptidique}, \emph{amphotère}, \emph{$\alpha$}.

Un acide $\alpha$-aminé est un composé organique qui possède un groupe carboxyle \ce{-COOH} et un groupe amine \ce{-NH2} fixés sur le carbone en position \ldots\ldots\ldots\ldots du groupe carboxyle. Lorsque ce carbone porte quatre groupes différents, c'est un \ldots\ldots\ldots\ldots et la molécule est chirale. En solution aqueuse, l'acide $\alpha$-aminé existe sous forme d'un ion dipolaire appelé \ldots\ldots\ldots\ldots ou \ldots\ldots\ldots\ldots, de formule $R-CH(NH_3^+)-COO^-$. Cet ion peut céder ou capter un proton : il présente un comportement \ldots\ldots\ldots\ldots. La réaction de \ldots\ldots\ldots\ldots entre deux acides $\alpha$-aminés, avec élimination d'une molécule d'eau, conduit à un dipeptide dans lequel les deux résidus sont unis par une liaison \ldots\ldots\ldots\ldots, qui est une liaison \ldots\ldots\ldots\ldots.
\end{exercice}

\begin{exercice}[Définitions]
Définir brièvement et précisément les termes suivants, en illustrant chacun par un exemple :
\begin{enumerate}
\item acide $\alpha$-aminé ;
\item amphion (zwitterion) ;
\item liaison peptidique ;
\item dipeptide ; extrémités N-terminale et C-terminale d'un peptide.
\end{enumerate}
\end{exercice}

\courssub{Je m'entraîne}

\begin{exercice}[L'alanine : nomenclature et stéréochimie]
On donne la formule semi-développée : \ce{CH3-CH(NH2)-COOH}.
\begin{enumerate}
\item Nommer ce composé en nomenclature systématique et donner son nom usuel ainsi que son abréviation.
\item Repérer, s'il existe, le carbone asymétrique de la molécule (le marquer d'un astérisque) et justifier.
\item Donner les représentations de Fischer des configurations L et D de cet acide $\alpha$-aminé.
\end{enumerate}
\end{exercice}

\begin{exercice}[Amphion de la valine]
La valine a pour formule semi-développée \ce{(CH3)2CH-CH(NH2)-COOH}.
\begin{enumerate}
\item Écrire l'équation-bilan mettant en évidence l'équilibre entre la forme moléculaire et la forme ionique (amphion) de la valine en solution aqueuse.
\item Représenter la formule semi-développée de l'amphion de la valine en précisant les charges portées.
\item L'amphion de la valine possède-t-il un carbone asymétrique ? Justifier.
\end{enumerate}
\end{exercice}

\begin{exercice}[Comportement amphotère de l'alanine]
L'amphion de l'alanine, de formule \ce{CH3-CH(NH3+)-COO-}, est mis en présence :
\begin{itemize}
\item d'une solution d'acide chlorhydrique \ce{(H3O+ + Cl-)} ;
\item puis, dans une seconde expérience, d'une solution d'hydroxyde de sodium \ce{(Na+ + HO-)}.
\end{itemize}
\begin{enumerate}
\item Écrire l'équation-bilan de la réaction de l'amphion avec les ions oxonium \ce{H3O+}. Quel caractère (acide ou basique) de l'amphion est mis en évidence ? Identifier le couple acide/base correspondant.
\item Écrire l'équation-bilan de la réaction de l'amphion avec les ions hydroxyde \ce{HO-}. Quel caractère de l'amphion est mis en évidence ? Identifier le couple acide/base correspondant.
\item Conclure sur le comportement amphotère de l'alanine.
\end{enumerate}
\end{exercice}

\begin{exercice}[Masse molaire d'un dipeptide]
On considère le dipeptide Ala--Ala obtenu par condensation de deux molécules d'alanine \ce{CH3-CH(NH2)-COOH}.

\emph{Données :} $M(\text{Ala}) = \SI{89}{g.mol^{-1}}$ ; $M(\ce{H2O}) = \SI{18}{g.mol^{-1}}$ ; $M(\ce{C}) = 12$ ; $M(\ce{H}) = 1$ ; $M(\ce{O}) = 16$ ; $M(\ce{N}) = \SI{14}{g.mol^{-1}}$.
\begin{enumerate}
\item Écrire l'équation-bilan de la formation du dipeptide Ala--Ala et donner sa formule brute.
\item Calculer la masse molaire du dipeptide Ala--Ala à partir de sa formule brute.
\item Vérifier la conservation de la masse : montrer que $M(\text{Ala--Ala}) = 2\,M(\text{Ala}) - M(\ce{H2O})$.
\end{enumerate}
\end{exercice}

\courssub{Je me prépare au Bac}

\begin{exercice}[Type Bac : synthèse du dipeptide Gly--Ala]
Une unité de production de compléments alimentaires à Bafoussam souhaite synthétiser le dipeptide Gly--Ala à partir de la glycine \ce{H2N-CH2-COOH} et de l'alanine \ce{CH3-CH(NH2)-COOH}.

\emph{Données :} $M(\text{Gly}) = \SI{75}{g.mol^{-1}}$ ; $M(\text{Ala}) = \SI{89}{g.mol^{-1}}$ ; $M(\ce{H2O}) = \SI{18}{g.mol^{-1}}$.
\begin{enumerate}
\item Donner la formule générale d'un acide $\alpha$-aminé. La glycine possède-t-elle un carbone asymétrique ? L'alanine en possède-t-elle un ? Justifier.
\item Représenter l'amphion de la glycine et écrire l'équation-bilan de l'équilibre entre la forme moléculaire et la forme ionique de la glycine en solution aqueuse.
\item Écrire l'équation-bilan de la réaction de formation du dipeptide Gly--Ala. Entourer la liaison peptidique dans la formule obtenue et donner le nom de cette fonction chimique.
\item Préciser, sur la formule du dipeptide, l'extrémité N-terminale et l'extrémité C-terminale.
\item La synthèse est réalisée \emph{sans protection} des fonctions amine et carboxyle : citer, avec leurs abréviations, les trois autres dipeptides susceptibles de se former dans le mélange réactionnel.
\item Calculer la masse molaire du dipeptide Gly--Ala. Quelle masse de dipeptide peut-on espérer obtenir, au maximum, à partir de \SI{15,0}{g} de glycine et d'un excès d'alanine, si le rendement de la synthèse est de $60\,\%$ ?
\end{enumerate}
\end{exercice}

\begin{exercice}[Type Bac : hydrolyse d'un tripeptide]
Au cours du contrôle qualité d'un hydrolysat de protéines destiné à l'alimentation infantile, un laboratoire de Douala isole un tripeptide $T$. L'hydrolyse totale de $T$ fournit trois acides $\alpha$-aminés : la glycine (Gly), l'alanine (Ala) et la valine (Val), chacun en proportion d'une mole par mole de tripeptide.

\emph{Données :} $M(\text{Gly}) = \SI{75}{g.mol^{-1}}$ ; $M(\text{Ala}) = \SI{89}{g.mol^{-1}}$ ; $M(\text{Val}) = \SI{117}{g.mol^{-1}}$ ; $M(\ce{H2O}) = \SI{18}{g.mol^{-1}}$.
\begin{enumerate}
\item Rappeler la définition d'un acide $\alpha$-aminé et donner la formule semi-développée de la valine, sachant que sa chaîne latérale $R$ est le groupe \ce{-CH(CH3)2}.
\item Parmi les trois acides $\alpha$-aminés obtenus, lesquels possèdent un carbone asymétrique ? Les identifier en marquant le carbone asymétrique d'un astérisque sur les formules.
\item Proposer deux enchaînements possibles du tripeptide $T$ et les nommer à l'aide des abréviations (par exemple Gly--Ala--Val).
\item Écrire l'équation-bilan de l'hydrolyse totale du tripeptide Gly--Ala--Val. Combien de molécules d'eau sont consommées par molécule de tripeptide hydrolysée ?
\item Calculer la masse molaire du tripeptide Gly--Ala--Val. Vérifier le résultat à l'aide de la relation $M(T) = M(\text{Gly}) + M(\text{Ala}) + M(\text{Val}) - 2\,M(\ce{H2O})$.
\item On hydrolyse totalement \SI{4,83}{g} de ce tripeptide. Calculer la quantité de matière de tripeptide hydrolysé, puis la masse de valine obtenue, l'hydrolyse étant supposée totale.
\end{enumerate}
\end{exercice}

\begin{exercice}[Type Bac : dosage et propriétés acido-basiques de l'alanine]
Un technicien d'une savonnerie--pharmacie artisanale de Yaoundé prépare une solution aqueuse d'alanine, acide $\alpha$-aminé utilisé comme complément. Il dissout une masse $m = \SI{4,45}{g}$ d'alanine \ce{CH3-CH(NH2)-COOH} dans de l'eau distillée pour obtenir $V = \SI{500}{mL}$ de solution $S$.

\emph{Données :} $M(\text{Ala}) = \SI{89}{g.mol^{-1}}$. Couples acide/base de l'alanine : $\ce{CH3-CH(NH3+)-COOH}/\ce{CH3-CH(NH3+)-COO-}$ : $pK_{a1} = 2,3$ ; $\ce{CH3-CH(NH3+)-COO-}/\ce{CH3-CH(NH2)-COO-}$ : $pK_{a2} = 9,9$.
\begin{enumerate}
\item Calculer la concentration molaire $C$ de la solution $S$.
\item Écrire l'équation-bilan de l'équilibre entre la forme moléculaire et la forme ionique (amphion) de l'alanine en solution aqueuse.
\item Tracer l'axe des $pK_a$ de l'alanine et placer les domaines de prédominance des trois formes : cation, amphion et anion. Quelle est la forme prédominante à $pH = 6$ ?
\item On ajoute à un échantillon de la solution $S$ quelques gouttes d'acide chlorhydrique concentré : le $pH$ devient égal à $1,0$. Quelle est la forme prédominante de l'alanine ? Écrire l'équation-bilan de la réaction de l'amphion avec les ions \ce{H3O+} et préciser le caractère mis en évidence.
\item À un second échantillon, on ajoute une solution d'hydroxyde de sodium jusqu'à $pH = 12$. Quelle est la forme prédominante ? Écrire l'équation-bilan de la réaction de l'amphion avec les ions \ce{HO-} et préciser le caractère mis en évidence.
\item Conclure sur le comportement amphotère de l'alanine et expliquer pourquoi on dit que l'amphion est un ampholyte.
\end{enumerate}
\end{exercice}

\newpage

\coursec{Corrigés des exercices}

\begin{corrige}[Exercice 1 — QCM : les acides $\alpha$-aminés]
\begin{enumerate}
\item \textbf{Réponse b.} Par définition, un acide $\alpha$-aminé possède le groupe amine \ce{-NH2} et le groupe carboxyle \ce{-COOH} portés par le \cle{même atome de carbone}, le carbone $\alpha$ (carbone voisin du carboxyle). La réponse a décrirait un acide $\beta$- ou $\gamma$-aminé ; la réponse c un simple acide carboxylique.
\item \textbf{Réponse a.} La formule générale est $R-CH(NH_2)-COOH$, où $R$ est un atome d'hydrogène (glycine) ou un groupe alkyle. La formule b est celle d'un acide $\alpha$-hydroxylé (comme l'acide lactique) et c celle d'un amide.
\item \textbf{Réponses a et b.} L'amphion $R-CH(NH_3^+)-COO^-$ porte \textbf{à la fois} une charge positive (sur l'azote) et une charge négative (sur l'oxygène du carboxylate) ; sa \textbf{charge globale est nulle}. C'est un ion dipolaire, d'où le nom de \emph{zwitterion} (« ion hybride » en allemand). Il n'est jamais chargé positivement dans son ensemble : c est faux.
\item \textbf{Réponse b.} La liaison peptidique est la liaison \ce{-CO-NH-}, c'est-à-dire une liaison \cle{amide}, formée entre le groupe carboxyle d'un acide $\alpha$-aminé et le groupe amine d'un autre. Une liaison ester serait \ce{-CO-O-} et une liaison éther \ce{-C-O-C-}.
\end{enumerate}
\end{corrige}

\begin{corrige}[Exercice 2 — Vrai ou faux ?]
\begin{enumerate}
\item \textbf{Faux.} Tous les acides $\alpha$-aminés possèdent un carbone asymétrique \emph{sauf la glycine} \ce{H2N-CH2-COOH}, dont le carbone $\alpha$ porte deux atomes d'hydrogène identiques. La glycine n'est donc pas chirale.
\item \textbf{Faux.} La glycine est bien un acide $\alpha$-aminé (les groupes \ce{NH2} et \ce{COOH} sont portés par le même carbone), mais ce carbone $\alpha$ porte deux hydrogènes : il n'est pas asymétrique, donc la molécule est \textbf{achirale}.
\item \textbf{Vrai.} En solution aqueuse s'établit l'équilibre de transfert interne de proton :
\[ \ce{R-CH(NH2)-COOH <=> R-CH(NH3+)-COO-} \]
La forme ionique (amphion) est très largement prédominante.
\item \textbf{Vrai.} L'amphion peut capter un proton (caractère basique, via \ce{COO-}) ou en céder un (caractère acide, via \ce{NH3+}) : c'est un \cle{ampholyte}, au comportement \cle{amphotère}.
\item \textbf{Vrai.} La formation d'un dipeptide est une \cle{condensation} : le groupe \ce{COOH} de l'un et le groupe \ce{NH2} de l'autre s'unissent avec élimination d'une molécule d'eau :
\[ \ce{H2N-CHR-COOH + H2N-CHR^1-COOH -> H2N-CHR-CO-NH-CHR^1-COOH + H2O} \]
\end{enumerate}
\end{corrige}

\begin{corrige}[Exercice 3 — Texte à trous]
Un acide $\alpha$-aminé est un composé organique qui possède un groupe carboxyle \ce{-COOH} et un groupe amine \ce{-NH2} fixés sur le carbone en position \textbf{$\alpha$} du groupe carboxyle. Lorsque ce carbone porte quatre groupes différents, c'est un \textbf{carbone asymétrique} et la molécule est chirale. En solution aqueuse, l'acide $\alpha$-aminé existe sous forme d'un ion dipolaire appelé \textbf{amphion} ou \textbf{zwitterion}, de formule $R-CH(NH_3^+)-COO^-$. Cet ion peut céder ou capter un proton : il présente un comportement \textbf{amphotère}. La réaction de \textbf{condensation} entre deux acides $\alpha$-aminés, avec élimination d'une molécule d'eau, conduit à un dipeptide dans lequel les deux résidus sont unis par une liaison \textbf{peptidique}, qui est une liaison \textbf{amide}.
\end{corrige}

\begin{corrige}[Exercice 4 — Définitions]
\begin{enumerate}
\item \textbf{Acide $\alpha$-aminé} : composé organique possédant un groupe carboxyle \ce{-COOH} et un groupe amine \ce{-NH2} portés par le même atome de carbone, appelé carbone $\alpha$. Formule générale : $R-CH(NH_2)-COOH$. \emph{Exemple :} l'alanine \ce{CH3-CH(NH2)-COOH}.
\item \textbf{Amphion (zwitterion)} : ion dipolaire de formule $R-CH(NH_3^+)-COO^-$, portant une charge positive et une charge négative mais de charge globale nulle ; il résulte du transfert interne du proton du groupe carboxyle vers le groupe amine. \emph{Exemple :} \ce{CH3-CH(NH3+)-COO-} pour l'alanine.
\item \textbf{Liaison peptidique} : liaison covalente \ce{-CO-NH-} (fonction amide) reliant le carbone du carboxyle d'un résidu d'acide $\alpha$-aminé à l'azote du groupe amine du résidu suivant. \emph{Exemple :} dans Ala--Gly, la liaison \ce{-CO-NH-} entre le \ce{C} de l'alanine et le \ce{N} de la glycine.
\item \textbf{Dipeptide} : molécule formée par condensation de deux acides $\alpha$-aminés, unis par une liaison peptidique, avec élimination d'eau. \emph{Exemple :} Ala--Ala. Un peptide possède toujours :
\begin{itemize}
\item une \textbf{extrémité N-terminale} : le résidu dont le groupe amine \ce{-NH2} est resté libre (écrit à gauche par convention) ;
\item une \textbf{extrémité C-terminale} : le résidu dont le groupe carboxyle \ce{-COOH} est resté libre (écrit à droite).
\end{itemize}
\end{enumerate}
\end{corrige}

\begin{corrige}[Exercice 5 — L'alanine : nomenclature et stéréochimie]
\begin{enumerate}
\item La chaîne principale comporte trois carbones avec un groupe carboxyle : acide propanoïque ; le groupe amine est en position 2. Nom systématique : \textbf{acide 2-aminopropanoïque}. Nom usuel : \textbf{alanine} ; abréviation : \textbf{Ala}.
\item Le carbone 2 (carbone $\alpha$) porte quatre groupes différents : \ce{-H}, \ce{-NH2}, \ce{-CH3} et \ce{-COOH}. C'est un \cle{carbone asymétrique} : \ce{CH3-C*H(NH2)-COOH}. La molécule est donc chirale.
\item Dans la représentation de Fischer d'un acide $\alpha$-aminé, le groupe \ce{-COOH} est placé en haut, la chaîne $R$ en bas ; la configuration est \textbf{L} si \ce{-NH2} est à gauche, \textbf{D} s'il est à droite :
\begin{center}
\chemfig{C(-[2]COOH)(-[6]CH_3)(-[4]NH_2)(-[0]H)}
\hspace{2cm}
\chemfig{C(-[2]COOH)(-[6]CH_3)(-[0]NH_2)(-[4]H)}
\end{center}
\begin{center}
\begin{tabular}{cc}
\textbf{L-alanine} (\ce{NH2} à gauche) & \textbf{D-alanine} (\ce{NH2} à droite)
\end{tabular}
\end{center}
\emph{(En projection de Fischer, les liaisons horizontales sortent du plan vers le lecteur, les liaisons verticales s'éloignent.)}
\end{enumerate}
\end{corrige}

\begin{corrige}[Exercice 6 — Amphion de la valine]
\begin{enumerate}
\item En solution aqueuse, le proton du groupe carboxyle est transféré sur le groupe amine ; il s'établit l'équilibre :
\[ \ce{(CH3)2CH-CH(NH2)-COOH <=> (CH3)2CH-CH(NH3+)-COO-} \]
\item Formule semi-développée de l'amphion :
\[ \underbrace{\ce{(CH3)2CH}}_{R}-\underset{+}{\overset{\displaystyle\ce{NH3}}{\underset{|}{\ce{CH}}}}-\ce{COO^-} \qquad\text{soit}\qquad \chemfig{(CH_3)_2CH-CH(-[2]NH_3^{+})-COO^{-}} \]
L'azote porte la charge $+$ (quatre liaisons), l'un des oxygènes du carboxylate porte la charge $-$ ; la charge globale est nulle.
\item \textbf{Oui.} Le carbone $\alpha$ de l'amphion porte quatre groupes différents : \ce{-H}, \ce{-NH3+}, \ce{-CH(CH3)2} et \ce{-COO-}. Il reste asymétrique : la formation de l'amphion ne détruit pas la chiralité de la molécule.
\end{enumerate}
\end{corrige}

\begin{corrige}[Exercice 7 — Comportement amphotère de l'alanine]
\begin{enumerate}
\item Avec les ions oxonium, le groupe carboxylate \ce{-COO-} capte un proton :
\[ \ce{CH3-CH(NH3+)-COO- + H3O+ -> CH3-CH(NH3+)-COOH + H2O} \]
L'amphion joue le rôle de \textbf{base} : c'est son \cle{caractère basique} qui est mis en évidence. Couple acide/base : $\ce{CH3-CH(NH3+)-COOH}/\ce{CH3-CH(NH3+)-COO-}$ (l'amphion est la base du couple).
\item Avec les ions hydroxyde, le groupe ammonium \ce{-NH3+} cède un proton :
\[ \ce{CH3-CH(NH3+)-COO- + HO- -> CH3-CH(NH2)-COO- + H2O} \]
L'amphion joue le rôle d'\textbf{acide} : c'est son \cle{caractère acide} qui est mis en évidence. Couple acide/base : $\ce{CH3-CH(NH3+)-COO-}/\ce{CH3-CH(NH2)-COO-}$ (l'amphion est l'acide du couple).
\item \textbf{Conclusion :} l'amphion de l'alanine réagit aussi bien avec un acide fort qu'avec une base forte : il possède un comportement \cle{amphotère} (c'est un ampholyte). Cette double réactivité explique le pouvoir tampon des solutions d'acides aminés.
\end{enumerate}
\end{corrige}

\begin{corrige}[Exercice 8 — Masse molaire d'un dipeptide]
\begin{enumerate}
\item Équation-bilan de la condensation :
\[ \ce{CH3-CH(NH2)-COOH + CH3-CH(NH2)-COOH -> CH3-CH(NH2)-CO-NH-CH(CH3)-COOH + H2O} \]
Formule brute de l'alanine : \ce{C3H7NO2}. Le dipeptide a pour formule brute :
\[ \ce{2 C3H7NO2 - H2O = C6H12N2O3} \]
\item Masse molaire du dipeptide \ce{C6H12N2O3} :
\[ M = 6\times 12 + 12\times 1 + 2\times 14 + 3\times 16 = 72 + 12 + 28 + 48 = \SI{160}{g.mol^{-1}} \]
\item Vérification de la conservation de la masse :
\[ 2\,M(\text{Ala}) - M(\ce{H2O}) = 2\times 89 - 18 = 178 - 18 = \SI{160}{g.mol^{-1}} \]
On retrouve bien $M(\text{Ala--Ala}) = \SI{160}{g.mol^{-1}}$ : la masse se conserve, la condensation éliminant une molécule d'eau par liaison peptidique formée.
\end{enumerate}
\end{corrige}

\begin{corrige}[Exercice 9 — Type Bac : synthèse du dipeptide Gly--Ala]
\textbf{Barème indicatif (sur 10) :} 1) 2 pts ; 2) 2 pts ; 3) 2 pts ; 4) 1 pt ; 5) 1,5 pt ; 6) 1,5 pt.

\begin{enumerate}
\item Formule générale d'un acide $\alpha$-aminé : $R-CH(NH_2)-COOH$.
\begin{itemize}
\item \textbf{Glycine} \ce{H2N-CH2-COOH} : le carbone $\alpha$ porte \emph{deux} atomes d'hydrogène identiques ; il n'est \textbf{pas} asymétrique.
\item \textbf{Alanine} \ce{CH3-CH(NH2)-COOH} : le carbone $\alpha$ porte \ce{-H}, \ce{-NH2}, \ce{-CH3}, \ce{-COOH}, quatre groupes différents : c'est un \textbf{carbone asymétrique}.
\end{itemize}
\item Amphion de la glycine : \ce{H3N+-CH2-COO-}. Équilibre en solution aqueuse :
\[ \ce{H2N-CH2-COOH <=> H3N+-CH2-COO-} \]
\item Équation-bilan de la formation de Gly--Ala :
\[ \ce{H2N-CH2-COOH + CH3-CH(NH2)-COOH -> H2N-CH2-\boxed{CO-NH}-CH(CH3)-COOH + H2O} \]
La liaison encadrée \ce{-CO-NH-} est la \textbf{liaison peptidique} ; la fonction chimique correspondante est la fonction \textbf{amide}.
\item Sur la formule \ce{H2N-CH2-CO-NH-CH(CH3)-COOH} :
\begin{itemize}
\item \textbf{extrémité N-terminale} : le résidu glycine, à gauche, dont le groupe \ce{-NH2} est libre ;
\item \textbf{extrémité C-terminale} : le résidu alanine, à droite, dont le groupe \ce{-COOH} est libre.
\end{itemize}
\item Sans protection des fonctions, chaque groupe amine peut réagir avec chaque groupe carboxyle : quatre dipeptides sont possibles. Outre \textbf{Gly--Ala}, on peut obtenir : \textbf{Ala--Gly}, \textbf{Gly--Gly} et \textbf{Ala--Ala}.
\item Formule brute de Gly--Ala : \ce{C2H5NO2 + C3H7NO2 - H2O = C5H10N2O3}.
\[ M(\text{Gly--Ala}) = 5\times 12 + 10\times 1 + 2\times 14 + 3\times 16 = 60 + 10 + 28 + 48 = \SI{146}{g.mol^{-1}} \]
Vérification : $75 + 89 - 18 = \SI{146}{g.mol^{-1}}$. \checkmark

Quantité de glycine (réactif limitant, alanine en excès) :
\[ n(\text{Gly}) = \frac{m}{M} = \frac{15,0}{75} = \SI{0,200}{mol} \]
La stœchiométrie étant 1:1, on peut former au maximum $n_{\max} = \SI{0,200}{mol}$ de dipeptide. Avec un rendement $\eta = 60\,\%$ :
\[ n_{\text{réel}} = \eta \times n_{\max} = 0,60 \times 0,200 = \SI{0,120}{mol} \]
\[ m(\text{Gly--Ala}) = n_{\text{réel}} \times M = 0,120 \times 146 = \SI{17,5}{g} \]
\end{enumerate}

\textbf{Points de vigilance :} ne pas oublier la molécule d'eau dans l'équation-bilan ; dans le calcul de masse molaire, soustraire $M(\ce{H2O})$ une seule fois ; bien distinguer Gly--Ala et Ala--Gly, qui sont deux dipeptides différents (l'ordre des résidus compte).
\end{corrige}

\begin{corrige}[Exercice 10 — Type Bac : hydrolyse d'un tripeptide]
\textbf{Barème indicatif (sur 10) :} 1) 1,5 pt ; 2) 1,5 pt ; 3) 1 pt ; 4) 2 pts ; 5) 1,5 pt ; 6) 2,5 pts.

\begin{enumerate}
\item Un \textbf{acide $\alpha$-aminé} est un composé organique possédant un groupe amine \ce{-NH2} et un groupe carboxyle \ce{-COOH} portés par le même atome de carbone (le carbone $\alpha$). Avec $R = \ce{-CH(CH3)2}$, la valine s'écrit :
\[ \ce{(CH3)2CH-CH(NH2)-COOH} \]
\item \textbf{Glycine} \ce{H2N-CH2-COOH} : pas de carbone asymétrique (deux H sur le carbone $\alpha$). \textbf{Alanine} et \textbf{valine} en possèdent un :
\[ \ce{CH3-C*H(NH2)-COOH} \qquad \ce{(CH3)2CH-C*H(NH2)-COOH} \]
Dans chaque cas, le carbone $\alpha$ étoilé porte quatre groupes différents.
\item Deux enchaînements possibles (parmi six) : \textbf{Gly--Ala--Val} et \textbf{Val--Gly--Ala} (on acceptera aussi Ala--Gly--Val, Gly--Val--Ala, Val--Ala--Gly, Ala--Val--Gly).
\item L'hydrolyse totale est la réaction inverse de la condensation ; elle consomme une molécule d'eau par liaison peptidique rompue, soit \textbf{deux molécules d'eau} pour un tripeptide :
\[ \ce{Gly-Ala-Val + 2 H2O -> Gly + Ala + Val} \]
soit, en formules :
\[ \ce{H2N-CH2-CO-NH-CH(CH3)-CO-NH-CH(CH(CH3)2)-COOH + 2 H2O -> H2N-CH2-COOH + CH3-CH(NH2)-COOH + (CH3)2CH-CH(NH2)-COOH} \]
\item Formule brute du tripeptide : \ce{C2H5NO2 + C3H7NO2 + C5H11NO2 - 2 H2O = C10H19N3O4}.
\[ M(T) = 10\times 12 + 19\times 1 + 3\times 14 + 4\times 16 = 120 + 19 + 42 + 64 = \SI{245}{g.mol^{-1}} \]
Vérification :
\[ M(T) = M(\text{Gly}) + M(\text{Ala}) + M(\text{Val}) - 2\,M(\ce{H2O}) = 75 + 89 + 117 - 36 = \SI{245}{g.mol^{-1}} \quad\checkmark \]
\item Quantité de tripeptide hydrolysé :
\[ n(T) = \frac{m}{M} = \frac{4,83}{245} = \SI{1,97e-2}{mol} \approx \SI{19,7}{mmol} \]
L'hydrolyse étant totale, la stœchiométrie donne $n(\text{Val}) = n(T) = \SI{1,97e-2}{mol}$. D'où :
\[ m(\text{Val}) = n(\text{Val}) \times M(\text{Val}) = 1,97\times 10^{-2} \times 117 \approx \SI{2,31}{g} \]
\end{enumerate}

\textbf{Points de vigilance :} compter \emph{deux} molécules d'eau pour un tripeptide (trois résidus = deux liaisons peptidiques) ; ne pas confondre hydrolyse (consomme l'eau) et condensation (libère l'eau) ; vérifier systématiquement la masse molaire par les deux méthodes.
\end{corrige}

\begin{corrige}[Exercice 11 — Type Bac : dosage et propriétés acido-basiques de l'alanine]
\textbf{Barème indicatif (sur 10) :} 1) 1,5 pt ; 2) 1,5 pt ; 3) 2 pts ; 4) 2 pts ; 5) 2 pts ; 6) 1 pt.

\begin{enumerate}
\item Quantité d'alanine dissoute :
\[ n = \frac{m}{M} = \frac{4,45}{89} = \SI{5,00e-2}{mol} \]
Concentration molaire de la solution $S$ ($V = \SI{500}{mL} = \SI{0,500}{L}$) :
\[ C = \frac{n}{V} = \frac{5,00\times 10^{-2}}{0,500} = \SI{0,100}{mol.L^{-1}} \]
\item Équilibre entre la forme moléculaire et l'amphion :
\[ \ce{CH3-CH(NH2)-COOH <=> CH3-CH(NH3+)-COO-} \]
\item Axe des $pK_a$ et domaines de prédominance :
\begin{popfigure}
\begin{tikzpicture}[scale=1.0]
\draw[-{Stealth}, thick] (0,0) -- (14.5,0) node[below left=2pt and -30pt] {$pH$};
\foreach \x/\lab in {2.3/{$pK_{a1}=2,3$}, 9.9/{$pK_{a2}=9,9$}}{
  \draw[thick] (\x,-0.12) -- (\x,0.12);
  \node[below] at (\x,-0.15) {\lab};
}
\node[above, popblue, font=\small\bfseries] at (1.15,0.35) {cation};
\node[above, popgreen, font=\small\bfseries] at (6.1,0.35) {amphion};
\node[above, poporange, font=\small\bfseries] at (12.2,0.35) {anion};
\node[below=14pt, popblue, font=\scriptsize, align=center] at (1.15,0) {\ce{CH3-CH(NH3+)-COOH}};
\node[below=14pt, popgreen, font=\scriptsize, align=center] at (6.1,0) {\ce{CH3-CH(NH3+)-COO-}};
\node[below=14pt, poporange, font=\scriptsize, align=center] at (12.2,0) {\ce{CH3-CH(NH2)-COO-}};
\draw[popink, thick, dashed] (6,-0.6) -- (6,0.9);
\node[popink, font=\scriptsize, anchor=west] at (6.05,0.75) {$pH=6$};
\end{tikzpicture}
\legende{Diagramme de prédominance des trois formes de l'alanine.}
\end{popfigure}

À $pH = 6$, on a $pK_{a1} < pH < pK_{a2}$ : la forme prédominante est l'\textbf{amphion} \ce{CH3-CH(NH3+)-COO-}.
\item À $pH = 1,0 < pK_{a1}$, la forme prédominante est le \textbf{cation} \ce{CH3-CH(NH3+)-COOH}. Réaction de l'amphion avec les ions oxonium :
\[ \ce{CH3-CH(NH3+)-COO- + H3O+ -> CH3-CH(NH3+)-COOH + H2O} \]
L'amphion capte un proton : c'est son \textbf{caractère basique} qui est mis en évidence (couple $\ce{CH3-CH(NH3+)-COOH}/\ce{CH3-CH(NH3+)-COO-}$).
\item À $pH = 12 > pK_{a2}$, la forme prédominante est l'\textbf{anion} \ce{CH3-CH(NH2)-COO-}. Réaction de l'amphion avec les ions hydroxyde :
\[ \ce{CH3-CH(NH3+)-COO- + HO- -> CH3-CH(NH2)-COO- + H2O} \]
L'amphion cède un proton : c'est son \textbf{caractère acide} qui est mis en évidence (couple $\ce{CH3-CH(NH3+)-COO-}/\ce{CH3-CH(NH2)-COO-}$).
\item \textbf{Conclusion :} l'amphion de l'alanine réagit avec les acides (comme une base) et avec les bases (comme un acide) : il présente un comportement \cle{amphotère}. On dit qu'il est un \cle{ampholyte} car il appartient, selon le couple considéré, soit au rôle de base (couple cation/amphion), soit au rôle d'acide (couple amphion/anion) : il est à la fois l'acide d'un couple et la base d'un autre.
\end{enumerate}

\textbf{Points de vigilance :} bien convertir le volume en litres ($\SI{500}{mL} = \SI{0,500}{L}$) ; sur l'axe des $pK_a$, l'espèce la plus protonée (cation) est à gauche, la plus déprotonée (anion) à droite ; ne pas inverser les caractères acide et basique : capter un proton = comportement basique.
\end{corrige}

\newpage

\coursec{Fiche bilan}

\coursec{Les acides $\alpha$-aminés et la synthèse des peptides}

\courssub{Structure et définition des acides $\alpha$-aminés}

\begin{definition}[Acide $\alpha$-aminé]
Un \cle{acide $\alpha$-aminé} est une molécule organique portant deux groupes fonctionnels sur le \emph{même} atome de carbone, appelé \cle{carbone $\alpha$} :
\begin{itemize}
  \item un groupe \cle{amine} \ce{-NH2} (fonction basique) ;
  \item un groupe \cle{carboxyle} \ce{-COOH} (fonction acide).
\end{itemize}
Sa formule générale s'écrit : $\ce{H2N-CHR-COOH}$ où \cle{R} représente la chaîne latérale (différente pour chaque acide aminé).
\end{definition}

\begin{popfigure}
\centering
\begin{tikzpicture}[
  atom/.style={circle, draw=popdark, fill=popblue!20, thick, minimum size=1.8em, inner sep=1pt, font=\small},
  bond/.style={thick, draw=popdark},
  labelstyle/.style={font=\small\bfseries, text=popdark}
]
  % Carbone alpha central
  \node[atom] (Calpha) at (0,0) {C$_\alpha$};
  % Groupes
  \node[labelstyle] (NH2) at (-1.8, 1.2) {\ce{NH2}};
  \node[labelstyle] (H) at (1.8, 1.2) {\ce{H}};
  \node[labelstyle] (COOH) at (-1.8, -1.2) {\ce{COOH}};
  \node[labelstyle] (R) at (1.8, -1.2) {R};
  % Liaisons
  \draw[bond] (Calpha) -- (-1.2, 0.8) node[midway, above left, font=\tiny] {};
  \draw[bond] (Calpha) -- (1.2, 0.8);
  \draw[bond] (Calpha) -- (-1.2, -0.8);
  \draw[bond] (Calpha) -- (1.2, -0.8);
  % Étoiles pour chiralité
  \node[font=\Large, text=poporange] at (0, -2.5) {$\ast$};
  \node[font=\small, align=center, text=poporange] at (0, -3.2) {Carbone asymétrique\\ (sauf si R = H)};
\end{tikzpicture}
\legende{Structure tétraédrique d'un acide $\alpha$-aminé. Le carbone $\alpha$ porte 4 groupes différents : \ce{NH2}, \ce{COOH}, \ce{H}, R. C'est un carbone chiral (asymétrique) pour tous les acides $\alpha$-aminés \textbf{sauf la glycine} (R = H).}
\end{popfigure}

\begin{propriete}[Les 20 acides $\alpha$-aminés protéinogènes (nomenclature)]
On retient surtout les plus courants. La nomenclature systématique (IUPAC) les nomme comme des \emph{acides 2-aminocarboxyliques}.
\begin{center}
\begin{tabularx}{\linewidth}{|l|l|l|X|l|}
\hline
\rowcolor{popblueL} \textbf{Trivial} & \textbf{Abrév. 3 lettres} & \textbf{1 lettre} & \textbf{Chaîne R} & \textbf{Nom systématique (IUPAC)} \\
\hline
Glycine & Gly & G & \ce{-H} & Acide 2-aminoéthanoïque \\
\hline
Alanine & Ala & A & \ce{-CH3} & Acide 2-aminopropanoïque \\
\hline
Valine & Val & V & \ce{-CH(CH3)2} & Acide 2-amino-3-méthylbutanoïque \\
\hline
Leucine & Leu & L & \ce{-CH2-CH(CH3)2} & Acide 2-amino-4-méthylpentanoïque \\
\hline
Sérine & Ser & S & \ce{-CH2OH} & Acide 2-amino-3-hydroxypropanoïque \\
\hline
Cystéine & Cys & C & \ce{-CH2SH} & Acide 2-amino-3-sulfanylpropanoïque \\
\hline
Acide aspartique & Asp & D & \ce{-CH2COOH} & Acide 2-aminobutanedioïque \\
\hline
Acide glutamique & Glu & E & \ce{-CH2CH2COOH} & Acide 2-aminopentanedioïque \\
\hline
Lysine & Lys & K & \ce{-CH2CH2CH2CH2NH2} & Acide 2,6-diaminohexanoïque \\
\hline
Phénylalanine & Phe & F & \ce{-CH2-Ph} & Acide 2-amino-3-phénylpropanoïque \\
\hline
\end{tabularx}
\end{center}
\emph{Remarque :} Seule la \textbf{glycine} (R = H) n'est pas chirale. Les 19 autres possèdent un carbone $\alpha$ asymétrique et existent sous deux énantiomères (séries L et D). Seuls les \textbf{acides $\alpha$-aminés de série L} entrent dans la composition des protéines naturelles.
\end{propriete}

\courssub{Nomenclature et représentation de Fischer}

\begin{methode}[Tracer la projection de Fischer d'un acide $\alpha$-aminé]
\begin{enumerate}
  \item Placer la chaîne carbonée principale \textbf{verticalement} avec le groupe le plus oxydé (\ce{COOH}) \textbf{en haut}.
  \item Le carbone $\alpha$ est au centre de la croix. Les liaisons \textbf{horizontales} sortent du plan (vers l'observateur), les liaisons \textbf{verticales} rentrent dans le plan.
  \item Placer le groupe \ce{R} \textbf{en bas} (prolongement de la chaîne).
  \item Placer \ce{H} et \ce{NH2} sur les liaisons horizontales :
  \begin{itemize}
    \item \ce{NH2} à \textbf{GAUCHE} $\rightarrow$ Configuration \cle{L} (naturelle).
    \item \ce{NH2} à \textbf{DROITE} $\rightarrow$ Configuration \cle{D}.
  \end{itemize}
\end{enumerate}
\end{methode}

\begin{popfigure}
\centering
\begin{tikzpicture}[
  fischer/.style={thick, draw=popdark},
  label/.style={font=\small\bfseries, text=popdark},
  chiral/.style={circle, draw=poporange, thick, minimum size=2.5em, inner sep=1pt}
]
  % L-Alanine
  \begin{scope}[xshift=-4cm]
    \node[chiral] (Cl) at (0,0) {};
    \draw[fischer] (0,1.5) -- (0,0.7) node[midway, right, label] {\ce{COOH}};
    \draw[fischer] (0,-0.7) -- (0,-1.5) node[midway, right, label] {\ce{CH3}};
    \draw[fischer] (-0.7,0) -- (-1.5,0) node[midway, above, label] {\ce{NH2}};
    \draw[fischer] (0.7,0) -- (1.5,0) node[midway, above, label] {\ce{H}};
    \node[below=0.5cm, font=\bfseries, text=popblue] at (0,-1.5) {L-Alanine};
    \node[above=0.5cm, font=\small, text=popmuted] at (0,1.5) {Série L};
  \end{scope}
  % D-Alanine
  \begin{scope}[xshift=4cm]
    \node[chiral] (Cd) at (0,0) {};
    \draw[fischer] (0,1.5) -- (0,0.7) node[midway, right, label] {\ce{COOH}};
    \draw[fischer] (0,-0.7) -- (0,-1.5) node[midway, right, label] {\ce{CH3}};
    \draw[fischer] (-0.7,0) -- (-1.5,0) node[midway, above, label] {\ce{H}};
    \draw[fischer] (0.7,0) -- (1.5,0) node[midway, above, label] {\ce{NH2}};
    \node[below=0.5cm, font=\bfseries, text=poporange] at (0,-1.5) {D-Alanine};
    \node[above=0.5cm, font=\small, text=popmuted] at (0,1.5) {Série D};
  \end{scope}
\end{tikzpicture}
\legende{Projections de Fischer de l'alanine. \textbf{Règle mnémotechnique} : « \textbf{L} = \ce{NH2} à Gauche (Left) ». Le groupe \ce{COOH} est toujours en haut, le groupe R (\ce{CH3} ici) en bas.}
\end{popfigure}

\courssub{Propriétés acido-basiques : l'amphion}

\begin{definition}[Amphion (ou zwitterion)]
En milieu aqueux, l'acide $\alpha$-aminé existe majoritairement sous forme d'un \cle{ion dipolaire} appelé \cle{amphion} (ou zwitterion), résultant d'un transfert de proton \textbf{intramoléculaire} du groupe carboxyle vers le groupe amine :
\[ \ce{H2N-CHR-COOH <=> ^{+}H3N-CHR-COO^{-}} \]
\begin{itemize}
  \item Le groupe \ce{-COOH} (acide) a cédé \ce{H+} $\rightarrow$ \ce{-COO^{-}}.
  \item Le groupe \ce{-NH2} (base) a capté \ce{H+} $\rightarrow$ \ce{-NH3^{+}}.
\end{itemize}
La molécule est globalement neutre mais porte des charges localisées.
\end{definition}

\begin{propriete}[Équilibre forme moléculaire / forme ionique]
L'équilibre entre la forme « neutre » (minoritaire) et l'amphion (majoritaire en solution) s'écrit :
\[ \ce{H2N-CHR-COOH <=> ^{+}H3N-CHR-COO^{-}} \]
Cet équilibre est fortement déplacé vers la droite (forme ionique) dans l'eau. La constante d'équilibre $K = \frac{[\text{amphion}]}{[\text{forme moléculaire}]}$ est très grande ($\sim 10^{10}$ pour la glycine).
\end{propriete}

\courssub{Comportement amphotère de l'amphion}

\begin{propriete}[Caractère amphotère de l'amphion]
L'amphion \ce{^{+}H3N-CHR-COO^{-}} peut réagir \textbf{soit comme un acide} (donneur de \ce{H+}), \textbf{soit comme une base} (accepteur de \ce{H+}).
\end{propriete}

\begin{exemplebox}[Équations-bilans du caractère amphotère (à connaître par cœur)]
Soit l'amphion de l'alanine : \ce{^{+}H3N-CH(CH3)-COO^{-}}.
\begin{enumerate}
  \item \textbf{Caractère acide} (réaction avec l'eau, don de \ce{H+} par \ce{-NH3^{+}}) :
  \[ \ce{^{+}H3N-CH(CH3)-COO^{-} + H2O <=> H2N-CH(CH3)-COO^{-} + H3O^{+}} \]
  \emph{Produit : forme anionique (milieu basique).}
  \item \textbf{Caractère basique} (réaction avec l'ion oxonium, capture de \ce{H+} par \ce{-COO^{-}}) :
  \[ \ce{^{+}H3N-CH(CH3)-COO^{-} + H3O^{+} <=> ^{+}H3N-CH(CH3)-COOH + H2O} \]
  \emph{Produit : forme cationique (milieu acide).}
  \item \textbf{Alternative pour le caractère basique} (réaction avec l'ion hydroxyde) :
  \[ \ce{^{+}H3N-CH(CH3)-COO^{-} + HO^{-} -> H2N-CH(CH3)-COO^{-} + H2O} \]
  \emph{Attention : ici l'amphion agit comme acide vis-à-vis de \ce{HO-}. La définition « base = capteur de \ce{H+} » privilégie l'équation (2) avec \ce{H3O+}.}
\end{enumerate}
\end{exemplebox}

\begin{attention}[Pièges fréquents sur les équations amphotères]
\begin{itemize}
  \item \textbf{Charges :} Vérifier systématiquement la conservation des charges (somme à gauche = somme à droite).
  \item \textbf{Atomes :} Compter les H, C, N, O. L'eau \ce{H2O} et \ce{H3O+}/\ce{HO-} sont des réactifs/produits, pas des spectateurs.
  \item \textbf{Flèches :} Utiliser \ce{<=>} pour les équilibres (réactions acide/base faibles), \ce{->} si réaction totale (ex: avec \ce{HO-} en excès).
  \item \textbf{Formules :} Ne pas confondre \ce{H2N-CHR-COO^{-}} (anion, forme basique) et \ce{^{+}H3N-CHR-COOH} (cation, forme acide).
\end{itemize}
\end{attention}

\begin{popfigure}
\centering
\begin{tikzpicture}[
  scale=0.9,
  axis/.style={->, thick, draw=popdark},
  tick/.style={thick, draw=popdark},
  zone/.style={fill=#1, opacity=0.3, draw=#1!80!black},
  pka/.style={font=\small\bfseries, text=popdark},
  form/.style={font=\small, align=center, text=popdark}
]
  % Axes
  \draw[axis] (0,0) -- (12,0) node[right] {pH};
  \draw[axis] (0,0) -- (0,3.5) node[above] {Forme prédominante};
  % Graduations pKa
  \draw[tick] (2.3,0) -- (2.3,-0.15) node[below] {p$K_{\text{a1}}$ $\approx$ 2,3};
  \draw[tick] (9.7,0) -- (9.7,-0.15) node[below] {p$K_{\text{a2}}$ $\approx$ 9,7};
  % Zones
  \fill[zone=popblue] (0,0.2) rectangle (2.3,1.5);
  \fill[zone=poppurple] (2.3,0.2) rectangle (9.7,1.5);
  \fill[zone=poporange] (9.7,0.2) rectangle (12,1.5);
  % Légendes formes
  \node[form, text=popblue] at (1.15, 0.8) {Cation\\ \ce{^{+}H3N-CHR-COOH}\\(Acide)};
  \node[form, text=poppurple] at (6, 0.8) {Amphion (Zwitterion)\\ \ce{^{+}H3N-CHR-COO^{-}}\\(Neutre global)};
  \node[form, text=poporange] at (10.85, 0.8) {Anion\\ \ce{H2N-CHR-COO^{-}}\\(Base)};
  % Flèches équilibres
  \draw[<->, thick, popblue] (1.15, 1.7) -- node[above, font=\tiny] {$\ce{+ H2O <=> H3O+}$} (3.5, 1.7);
  \draw[<->, thick, poporange] (6, 2.2) -- node[above, font=\tiny] {$\ce{+ H2O <=> H3O+}$} (10.85, 2.2);
  % pH physiologique
  \draw[dashed, thick, popgreen] (7.4,0) -- (7.4,3);
  \node[above, font=\small\bfseries, text=popgreen] at (7.4, 3) {pH sanguin $\approx$ 7,4};
\end{tikzpicture}
\legende{Diagramme de prédominance des formes d'un acide $\alpha$-aminé (ex: alanine) en fonction du pH. p$K_{\text{a1}}$ (groupe \ce{COOH}) $\approx$ 2,3 ; p$K_{\text{a2}}$ (groupe \ce{NH3+}) $\approx$ 9,7. Au pH physiologique (7,4), l'amphion est la forme quasi exclusive.}
\end{popfigure}

\courssub{La liaison peptidique et les dipeptides}

\begin{definition}[Liaison peptidique]
La \cle{liaison peptidique} est une liaison amide $\mathrm{-CO-NH-}$ formée par condensation entre le groupe carboxyle $\mathrm{-COOH}$ d'un acide aminé et le groupe amine $\mathrm{-NH_2}$ d'un autre acide aminé, avec élimination d'une molécule d'eau.
\[ \mathrm{H_2N{-}CHR{-}COOH} + \mathrm{H_2N{-}CHR'{-}COOH} \longrightarrow \mathrm{H_2N{-}CHR{-}CO{-}NH{-}CHR'{-}COOH} + \ce{H2O} \]
\end{definition}

\begin{propriete}[Structure d'un dipeptide]
Un dipeptide possède deux extrémités distinctes :
\begin{itemize}
  \item \cle{Extrémité N-terminale} (libre) : groupe amine \ce{H2N-} (porte le résidu N-terminal).
  \item \cle{Extrémité C-terminale} (libre) : groupe carboxyle \ce{-COOH} (porte le résidu C-terminal).
\end{itemize}
La liaison peptidique \ce{-CO-NH-} est \textbf{plane} et \textbf{rigide} (caractère partiel de double liaison par mésomérie), ce qui structure la chaîne polypeptidique.
\end{propriete}

\begin{methode}[Nommer un dipeptide (nomenclature)]
\begin{enumerate}
  \item Identifier le résidu \textbf{N-terminal} (à gauche, amine libre).
  \item Modifier son nom en remplaçant la terminaison \emph{-ine} (ou \emph{-ique}) par \emph{-yle} (ex: glycine $\rightarrow$ \textbf{glycyl}, alanine $\rightarrow$ \textbf{alanyl}, valine $\rightarrow$ \textbf{valyl}).
  \item Écrire le nom du résidu \textbf{C-terminal} (à droite, carboxyle libre) \textbf{sans modification}.
  \item Relier les deux noms par un trait d'union.
\end{enumerate}
\emph{Exemple :} \ce{Gly-Ala} = \textbf{glycyl-alanine}. \ce{Ala-Gly} = \textbf{alanyl-glycine}. Ce sont deux dipeptides \textbf{différents} (isomères de constitution).
\end{methode}

\begin{exoresolu}[Formation du dipeptide Glycyl-Alanine]
\textbf{Énoncé :} Écrire l'équation-bilan de la formation du dipeptide glycyl-alanine à partir de la glycine et de l'alanine. Entourer la liaison peptidique formée.
\tcblower
\textbf{Solution :}
\begin{itemize}
  \item Réactifs : Glycine \ce{H2N-CH2-COOH} + Alanine \ce{H2N-CH(CH3)-COOH}.
  \item Produits : Glycyl-Alanine \ce{H2N-CH2-CO-NH-CH(CH3)-COOH} + \ce{H2O}.
\end{itemize}
\[ \ce{H2N-CH2-COOH + H2N-CH(CH3)-COOH -> H2N-CH2-CO-NH-CH(CH3)-COOH + H2O} \]
\chemfig{H_2N-CH_2-C(=[1]O)-NH-CH(-[6]CH_3)-C(=[1]O)OH} \quad \text{La liaison peptidique \ce{-CO-NH-} est encadrée.}
\end{exoresolu}

\begin{attention}[Ordre des acides aminés]
Deux acides aminés différents (ex: Gly et Ala) donnent \textbf{deux dipeptides distincts} :
\begin{itemize}
  \item \ce{Gly-Ala} (Glycyl-Alanine) : Gly en N-terminal, Ala en C-terminal.
  \item \ce{Ala-Gly} (Alanyl-Glycine) : Ala en N-terminal, Gly en C-terminal.
\end{itemize}
Ils ont la même formule brute mais des propriétés physico-chimiques et biologiques différentes. Ne jamais inverser l'ordre dans le nom !
\end{attention}

\courssub{Synthèse dirigée d'un dipeptide (complément Bac)}

\begin{savaistu}[Pourquoi une synthèse « dirigée » ?]
Si l'on mélange simplement deux acides aminés différents (A et B) et qu'on chauffe pour éliminer l'eau, on obtient un \textbf{mélange de 4 dipeptides} : A-A, B-B, A-B, B-A. Pour obtenir \textbf{un seul dipeptide pur} (ex: A-B), il faut \textbf{bloquer sélectivement} les fonctions qui ne doivent pas réagir. C'est le principe de la \cle{synthèse dirigée} (ou synthèse en phase liquide), base de la synthèse peptidique en phase solide (Merrifield, Prix Nobel 1984).
\end{savaistu}

\begin{methode}[Étapes de la synthèse dirigée d'un dipeptide A-B]
\begin{enumerate}
  \item \textbf{Protection de l'amine de A (composant N-terminal)} \\
  On transforme \ce{NH2} en un groupe protégé stable (ex: \cle{Boc} = \ce{tert-butyloxycarbonyl} ou \cle{Z} = \ce{benzyloxycarbonyl}).
  \[ \ce{H2N-CHR_A-COOH ->[Boc_2O, base] Boc-NH-CHR_A-COOH} \]
  \item \textbf{Protection du carboxyle de B (composant C-terminal)} \\
  On estérifie le \ce{COOH} (ex: ester méthylique ou benzyle).
  \[ \ce{H2N-CHR_B-COOH ->[MeOH, H+] H2N-CHR_B-COOMe} \]
  \item \textbf{Activation du carboxyle de A} \\
  On rend le \ce{COOH} de A plus réactif (ex: chlorure d'acide via \ce{SOCl2}, ou anhydride mixte, ou activation par \cle{DCC} = dicyclohexylcarbodiimide).
  \[ \ce{Boc-NH-CHR_A-COOH ->[DCC] Boc-NH-CHR_A-CO-O-C(=O)-N(Cy)2} \quad \text{(ester actif)} \]
  \item \textbf{Condensation (Couplage)} \\
  On met en présence l'amine libre de B (déprotégée si besoin, ou on utilise B avec \ce{NH2} libre dès le départ) et l'acide activé de A. La liaison peptidique se forme.
  \[ \ce{Boc-NH-CHR_A-CO-} \text{(activé)} \ce{+ H2N-CHR_B-COOMe -> Boc-NH-CHR_A-CO-NH-CHR_B-COOMe} \]
  \item \textbf{Déprotection finale} \\
  On retire les groupes protecteurs (Boc par acide fort \ce{TFA} ou \ce{HCl} ; Ester par hydrolyse basique \ce{NaOH} ou \ce{H2/Pd} pour Z/OBn) pour obtenir le dipeptide libre \ce{H2N-CHR_A-CO-NH-CHR_B-COOH}.
\end{enumerate}
\end{methode}

\begin{exemplebox}[Schéma récapitulatif : Synthèse de l'Alanyl-Glycine (Ala-Gly)]
\begin{center}
\begin{tikzpicture}[node distance=1.2cm and 0.5cm, every node/.style={font=\small, align=center}]
  \node (startA) {Alanine (A)\\\ce{H2N-CH(CH3)-COOH}};
  \node (protA) [right=of startA] {Protection \ce{NH2} (Boc)\\\ce{Boc-NH-CH(CH3)-COOH}};
  \node (actA) [right=of protA] {Activation \ce{COOH} (DCC)\\\ce{Boc-NH-CH(CH3)-CO-} (actif)};
  \node (startB) [below=of startA] {Glycine (B)\\\ce{H2N-CH2-COOH}};
  \node (protB) [below=of protA] {Protection \ce{COOH} (MeOH/H+)\\\ce{H2N-CH2-COOMe}};
  \node (couple) [right=of actA] {Couplage\\\ce{Boc-NH-CH(CH3)-CO-NH-CH2-COOMe}};
  \node (final) [right=of couple] {Déprotection (TFA + NaOH)\\\ce{H2N-CH(CH3)-CO-NH-CH2-COOH}\\(Alanyl-Glycine)};
  \draw[->, thick, popblue] (startA) -- (protA);
  \draw[->, thick, popblue] (protA) -- (actA);
  \draw[->, thick, popblue] (startB) -- (protB);
  \draw[->, thick, popblue] (actA) -- node[above] {+ \ce{H2N-CH2-COOMe}} (couple);
  \draw[->, thick, popblue] (protB) -- (couple);
  \draw[->, thick, popgreen] (couple) -- (final);
\end{tikzpicture}
\end{center}
\end{exemplebox}

\begin{aretenir}[Checklist Bac - Acides $\alpha$-aminés et Peptides]
\begin{itemize}
  \item[$\square$] Écrire la formule générale $\ce{H2N-CHR-COOH}$ et identifier le carbone $\alpha$.
  \item[$\square$] Repérer la chiralité : carbone asymétrique si R $\neq$ H (glycine achirale). Marquer l'étoile $\ast$.
  \item[$\square$] Citer 3 à 5 acides aminés courants avec leur chaîne R (Gly, Ala, Val, Ser, Lys...).
  \item[$\square$] Dessiner la projection de Fischer : \ce{COOH} en haut, R en bas. \ce{NH2} à gauche = \textbf{L} (naturel), à droite = \textbf{D}.
  \item[$\square$] Écrire l'équilibre forme moléculaire $\rightleftharpoons$ amphion : $\ce{H2N-CHR-COOH <=> ^{+}H3N-CHR-COO^{-}}$.
  \item[$\square$] Écrire les 2 équations du caractère amphotère de l'amphion :
  \begin{itemize}
    \item Acide : $\ce{^{+}H3N-CHR-COO^{-} + H2O <=> H2N-CHR-COO^{-} + H3O^{+}}$
    \item Base : $\ce{^{+}H3N-CHR-COO^{-} + H3O^{+} <=> ^{+}H3N-CHR-COOH + H2O}$
  \end{itemize}
  \item[$\square$] Reconnaître la liaison peptidique \ce{-CO-NH-} (amide) dans une chaîne.
  \item[$\square$] Écrire l'équation-bilan de formation d'un dipeptide (+ \ce{H2O}).
  \item[$\square$] Nommer un dipeptide : résidu N-terminal en \emph{-yle}, résidu C-terminal inchangé (ex: \textbf{glycyl-alanine}).
  \item[$\square$] Savoir que 2 AA différents $\rightarrow$ 2 dipeptides possibles (ordre N $\rightarrow$ C).
  \item[$\square$] Connaître les 4 étapes de la synthèse dirigée : Protection (N de A, C de B), Activation (C de A), Couplage, Déprotection.
  \item[$\square$] Vérifier la conservation des atomes et des charges dans toutes les équations.
\end{itemize}
\end{aretenir}

\findecours

\end{document}
